\documentclass[10pt,a4paper,landscape]{article}
\usepackage[landscape,margin=1.2cm]{geometry}
\usepackage[ngerman]{babel}
\usepackage{fontspec}
\usepackage{longtable}
\usepackage{array}
\usepackage[table]{xcolor}
\usepackage{ragged2e}

\definecolor{diffred}{RGB}{180, 0, 0}
\definecolor{diffgreen}{RGB}{0, 130, 0}
\setlength{\LTpre}{0pt}
\setlength{\LTpost}{0pt}
\setlength{\tabcolsep}{3pt}
\renewcommand{\arraystretch}{1.1}

\newcolumntype{L}[1]{>{\RaggedRight\arraybackslash}p{#1}}

\begin{document}
\scriptsize

\begin{center}
\begin{tabular}{|p{2.8cm}|p{6.9cm}|p{13.0cm}|}
\hline
\textbf{Markierungsfarbe} & \textbf{Änderungen RefE 2024} & \textbf{Änderungen RefE 2026} \\
\hline
\textcolor{diffgreen}{Grün} & - & Wurde im Vergleich zum RefE 2024 (oder geltendem Recht, falls RefE 2024 nicht vorhanden / unverändert) hinzugefügt \\
\hline
\textcolor{diffred}{Rot} & Wurde im Vergleich zum RefE 2026 gelöscht & - \\
\hline
\end{tabular}
\end{center}
\vspace{0.3cm}

\begin{center}
{\Large\bfseries Synopse IKJHG - Vergleich der Referentenentwürfe 2024 und 2026}
\end{center}
\vspace{0.5cm}

\begin{longtable}{|L{8.2cm}|L{8.2cm}|L{8.2cm}|}
\hline
\multicolumn{1}{|c|}{\cellcolor{gray!15}\textbf{Geltendes Recht (kombiniert)}} & \multicolumn{2}{c|}{\cellcolor{gray!15}\textbf{Änderungen durch den Referentenentwurf}} \\
\hline
\textbf{Synopsis 2024/2026} & \textbf{Synopsis 2024} & \textbf{Synopsis 2026} \\
\hline
\endhead

- Aus Synopsis 2024 - \newline \textbf{( - SGB VIII)  \newline vom: 11.09.2012 - zuletzt geändert durch  \newline Art. 5 v. 8.5.2024 I Nr. 152 } \newline - Aus Synopsis 2026 - \newline \textbf{( - SGB VIII)  \newline vom: 11.09.2012 - zuletzt geändert durch  \newline Art. 2 v. 3.4.2025 I Nr. 107 } & \textbf{( - SGB VIII)  \newline vom: 11.09.2012 - zuletzt geändert durch  \newline Art. }\textbf{\textcolor{diffred}{5}}\textbf{ v. }\textbf{\textcolor{diffred}{8}}\textbf{.}\textbf{\textcolor{diffred}{5}}\textbf{.202}\textbf{\textcolor{diffred}{4}}\textbf{ I Nr. 1}\textbf{\textcolor{diffred}{52}} & \textbf{( - SGB VIII)  \newline vom: 11.09.2012 - zuletzt geändert durch  \newline Art. }\textbf{\textcolor{diffgreen}{2}}\textbf{ v. }\textbf{\textcolor{diffgreen}{3}}\textbf{.}\textbf{\textcolor{diffgreen}{4}}\textbf{.202}\textbf{\textcolor{diffgreen}{5}}\textbf{ I Nr. 1}\textbf{\textcolor{diffgreen}{07}} \\
\hline
\rowcolor{gray!25}
\textbf{\S{} 1} &  & \textbf{\S{} 1} \\
\hline
Recht auf Erziehung, Elternverantwortung,  \newline Jugendhilfe &  & \textbf{Recht auf }\textbf{\textcolor{diffgreen}{Förderung der Entwicklung,  \newline auf Erziehung und auf Teilhabe}} \\
\hline
(1) Jeder junge Mensch hat ein Recht auf  \newline Förderung seiner Entwicklung und auf  \newline Erziehung zu einer selbstbestimmten,  \newline eigenverantwortlichen und  \newline gemeinschaftsfähigen Persönlichkeit &  & (1) Jeder junge Mensch hat ein Recht auf  \newline Förderung seiner Entwicklung und auf  \newline Erziehung zu einer selbstbestimmten,  \newline eigenverantwortlichen und  \newline gemeinschaftsfähigen Persönlichkeit \textbf{\textcolor{diffgreen}{und  \newline auf Förderung seiner vollen, wirksamen  \newline und gleichberechtigten Teilhabe am  \newline Leben in der Gesellschaft}} \\
\hline
(2) Pflege und Erziehung der Kinder sind  \newline das natürliche Recht der Eltern und die  \newline zuvörderst ihnen obliegende Pflicht. Über  \newline ihre Betätigung wacht die staatliche  \newline Gemeinschaft. &  & (2) unverändert \\
\hline
(3) Jugendhilfe soll zur Verwirklichung des  \newline Rechts nach Absatz 1 insbesondere &  & (3) unverändert \\
\hline
1. junge Menschen in ihrer individuellen  \newline und sozialen Entwicklung fördern und dazu  \newline beitragen, Benachteiligungen zu vermeiden  \newline oder abzubauen, &  & 1. unverändert \\
\hline
2. jungen Menschen ermöglichen oder  \newline erleichtern, entsprechend ihrem Alter und  ihrer individuellen Fähigkeiten in allen sie  \newline betreffenden Lebensbereichen  \newline selbstbestimmt zu interagieren und damit  \newline gleichberechtigt am Leben in der  \newline Gesellschaft teilhaben zu können, &  & 2. unverändert \\
\hline
3. Eltern und andere Erziehungsberechtigte  \newline bei der Erziehung beraten und  \newline unterstützen, &  & 3. unverändert \\
\hline
4. Kinder und Jugendliche vor Gefahren für  \newline ihr Wohl schützen &  & 4. unverändert \\
\hline
5. dazu beitragen, positive  \newline Lebensbedingungen für junge Menschen  \newline und ihre Familien sowie eine kinder- und  \newline familienfreundliche Umwelt zu erhalten  \newline oder zu schaffen &  & 5. unverändert \\
\hline
... & ... & ... \\
\hline
1.  \newline Angebote der Jugendarbeit, der  \newline Jugendsozialarbeit, der Schulsozialarbeit  \newline und des erzieherischen Kinder- und  \newline Jugendschutzes (\S{}\S{} 11 bis 14), & 1.  \newline Angebote der Jugendarbeit, der  \newline Jugendsozialarbeit, der Schulsozialarbeit  \newline und des erzieherischen Kinder- und  \newline Jugendschutzes (\S{}\S{} 11 bis 14), & \textbf{1. }\textbf{\textcolor{diffgreen}{Angebote der Beratung für Kinder  \newline und Jugendliche (\S{} 8 Absatz 3) und der  \newline Beratung, Vermittlung und Klärung in  \newline Konflikten durch Ombudsstellen (\S{} 9a)  \newline sowie die Unterstützung und Begleitung  \newline durch einen Verfahrenslotsen (\S{} 10b), }} \textbf{\textcolor{diffgreen}{2}}\textcolor{diffgreen}{.}  \newline Angebote der Jugendarbeit, der  \newline Jugendsozialarbeit, der Schulsozialarbeit  \newline und des erzieherischen Kinder- und  \newline Jugendschutzes (\S{}\S{} 11 bis 14), \\
\hline
2.  \newline Angebote zur Förderung der Erziehung in  \newline der Familie (\S{}\S{} 16 bis 21), & \textcolor{diffred}{2}.  \newline Angebote zur Förderung der Erziehung in  \newline der Familie (\S{}\S{} 16 bis 21), & \textbf{\textcolor{diffgreen}{3}}.  \newline Angebote zur Förderung der Erziehung in  \newline der Familie (\S{}\S{} 16 bis 21), \\
\hline
3.  \newline Angebote zur Förderung von Kindern in  \newline Tageseinrichtungen und in  \newline Kindertagespflege (\S{}\S{} 22 bis 25), & \textcolor{diffred}{3}.  \newline Angebote zur Förderung von Kindern in  \newline Tageseinrichtungen und in  \newline Kindertagespflege (\S{}\S{} 22 bis 25), & \textbf{\textcolor{diffgreen}{4}}.  \newline Angebote zur Förderung von Kindern in  \newline Tageseinrichtungen und in  \newline Kindertagespflege (\S{}\S{} 22 bis 25), \\
\hline
4.  \newline Hilfe zur Erziehung und ergänzende  \newline Leistungen (\S{}\S{} 27 bis 35, 36, 37, 39, 40), & \textcolor{diffred}{4.  \newline Hilfe} zur Erziehung und \textcolor{diffred}{ergänzende  \newline Leistungen} \textbf{(\S{}\S{} 27 bis 40),} & \textbf{\textcolor{diffgreen}{5.  \newline Leistungen zur Entwicklung,}}\textbf{ zur  \newline Erziehung und }\textbf{\textcolor{diffgreen}{zur Teilhabe}} \textbf{(\S{}\S{} 27 bis  \newline 40),} \\
\hline
... & ... & ... \\
\hline
 & \textbf{b) Leistungen der Eingliederungshilfe  \newline für Kinder und Jugendliche mit  \newline Behinderungen und ergänzende  \newline Leistungen (\S{}\S{} 27 Absatz 3 }\textbf{\textcolor{diffred}{bis 3b,}}\textbf{ 35a  \newline bis 40), } & \textbf{b) Leistungen der Eingliederungshilfe  \newline für Kinder und Jugendliche mit  \newline Behinderungen und ergänzende  \newline Leistungen (\S{}\S{} 27 Absatz 3 }\textbf{\textcolor{diffgreen}{und 4,}}\textbf{ 35a  \newline bis 40), } \\
\hline
5.  \newline Hilfe für seelisch behinderte Kinder und  \newline Jugendliche und ergänzende Leistungen  \newline (\S{}\S{} 35a bis 37, 39, 40), & 5. \textcolor{diffred}{entfällt} & 5.  \newline \textcolor{diffgreen}{Hilfe für seelisch behinderte Kinder und  \newline Jugendliche und ergänzende Leistungen  \newline (\S{}\S{} 35a bis 37, 39, 40),} \\
\hline
6.  \newline Hilfe für junge Volljährige und  \newline Nachbetreuung (den \S{}\S{} 41 und 41a). & \textbf{\textcolor{diffred}{5}}\textbf{.}  \newline Hilfe für junge Volljährige und  \newline Nachbetreuung (den \S{}\S{} 41 und 41a). & \textbf{\textcolor{diffgreen}{6}}.  \newline Hilfe für junge Volljährige und  \newline Nachbetreuung (den \S{}\S{} 41 und 41a). \\
\hline
... & ... & ... \\
\hline
\rowcolor{gray!25}
\textbf{\S{} 3} &  & \textbf{\S{} 3} \\
\hline
Freie und öffentliche Jugendhilfe &  & Freie und öffentliche Jugendhilfe \\
\hline
(1) Die Jugendhilfe ist gekennzeichnet  \newline durch die Vielfalt von Trägern  \newline unterschiedlicher Wertorientierungen und  \newline die Vielfalt von Inhalten, Methoden und  \newline Arbeitsformen. &  & (1) unverändert \\
\hline
(2) Leistungen der Jugendhilfe werden von  \newline Trägern der freien Jugendhilfe und von  \newline Trägern der öffentlichen Jugendhilfe  \newline erbracht. Leistungsverpflichtungen, die  \newline durch dieses Buch begründet werden,  \newline richten sich an die Träger der öffentlichen  \newline Jugendhilfe. &  & (2) unverändert \\
\hline
(3) Andere Aufgaben der Jugendhilfe  \newline werden von Trägern der öffentlichen  \newline Jugendhilfe wahrgenommen. Soweit dies  \newline ausdrücklich bestimmt ist, können Träger  \newline der freien Jugendhilfe diese Aufgaben  \newline wahrnehmen oder mit ihrer Ausführung  \newline betraut werden. &  & (3) unverändert \\
\hline
 &  & \textbf{\textcolor{diffgreen}{(4) Die Träger der öffentlichen  \newline Jugendhilfe sind für Leistungen nach \S{}  \newline 2 Absatz 2 Nummer 4 Buchstabe b  \newline Rehabilitationsträger im Sinne von \S{} 6  \newline Absatz 1 Nummer 6 des Neunten  \newline Buches. \S{} 7 des Neunten Buches ist zu  \newline beachten.}} \\
\hline
... & ... & ... \\
\hline
 & \textbf{(3) }\textbf{\textcolor{diffred}{Bei der Entscheidung nach Absatz 2  \newline ist zunächst die Zumutbarkeit einer  \newline Abweichung}}\textbf{ von den Wünschen des  \newline Leistungsberechtigten }\textbf{\textcolor{diffred}{zu prüfen. Dabei  \newline sin}}\textbf{ die persönlichen, familiären und  \newline örtlichen Umstände einschließlich der  \newline gewünschten Wohnform angemessen  \newline zu berücksichtigen. Bei Unzumutbarkeit  \newline einer abweichenden  \newline Leistungsgestaltung ist ein  \newline Kostenvergleich nicht vorzunehmen.  \newline Für Leistungsberechtigte nach \S{} 27  \newline Absatz 3 gilt im Übrigen \S{} 104 Absatz }\textbf{\textcolor{diffred}{3  \newline Satz 3 und}}\textbf{ 4 des Neunten Buches  \newline entsprechend. } & \textbf{(3) }\textbf{\textcolor{diffgreen}{Eine}}\textbf{ von den Wünschen des  \newline Leistungsberechtigten }\textbf{\textcolor{diffgreen}{abweichende  \newline Leistung kommt nicht in Betracht, wenn  \newline diese Abweichung für den  \newline Leistungsberechtigten unzumutbar ist.  \newline Bei der Prüfung der Zumutbarkeit sind}}\textbf{  \newline die persönlichen, familiären und  \newline örtlichen Umstände einschließlich der  \newline gewünschten Wohnform angemessen  \newline zu berücksichtigen. Bei Unzumutbarkeit  \newline einer abweichenden  \newline Leistungsgestaltung ist ein  \newline Kostenvergleich nicht vorzunehmen.  \newline Für Leistungsberechtigte nach \S{} 27  \newline Absatz 3 gilt im Übrigen \S{} 104 Absatz 4  \newline des Neunten Buches entsprechend. } \\
\hline
\rowcolor{gray!25}
\textbf{\S{} 8a} &  & \textbf{\S{} 8a} \\
\hline
Schutzauftrag bei Kindeswohlgefährdung &  & Schutzauftrag bei Kindeswohlgefährdung \\
\hline
(1) Werden dem Jugendamt gewichtige  \newline Anhaltspunkte für die Gefährdung des  \newline Wohls eines Kindes oder Jugendlichen  \newline bekannt, so hat es das Gefährdungsrisiko  \newline im Zusammenwirken mehrerer Fachkräfte  \newline einzuschätzen. Soweit der wirksame  \newline Schutz dieses Kindes oder dieses  \newline Jugendlichen nicht in Frage gestellt wird,  \newline hat das Jugendamt die  \newline Erziehungsberechtigten sowie das Kind  \newline oder den Jugendlichen in die  \newline Gefährdungseinschätzung einzubeziehen  \newline und, sofern dies nach fachlicher  \newline Einschätzung erforderlich ist,   \newline 1. sich dabei einen unmittelbaren Eindruck  \newline von dem Kind und von seiner persönlichen  \newline Umgebung zu verschaffen sowie   \newline 2. Personen, die gemäß \S{} 4 Absatz 3 des  \newline Gesetzes zur Kooperation und Information  \newline im Kinderschutz dem Jugendamt Daten  \newline übermittelt haben, in geeigneter Weise an  der Gefährdungseinschätzung zu  \newline beteiligen. Hält das Jugendamt zur  \newline Abwendung der Gefährdung die  \newline Gewährung von Hilfen für geeignet und  \newline notwendig, so hat es diese den  \newline Erziehungsberechtigten anzubieten. &  & (1) Werden dem Jugendamt gewichtige  \newline Anhaltspunkte für die Gefährdung des  \newline Wohls eines Kindes oder Jugendlichen  \newline bekannt, so hat es das Gefährdungsrisiko  \newline im Zusammenwirken mehrerer Fachkräfte  \newline \textbf{\textcolor{diffgreen}{und, soweit dies notwendig ist, unter  \newline Beteiligung betroffener Einrichtungen,  \newline Dienste oder anderer Stellen}}  \newline einzuschätzen. Soweit der wirksame  \newline Schutz dieses Kindes oder dieses  \newline Jugendlichen nicht in Frage gestellt wird,  \newline hat das Jugendamt die  \newline Erziehungsberechtigten sowie das Kind  \newline oder den Jugendlichen in die  \newline Gefährdungseinschätzung einzubeziehen  \newline und, sofern dies nach fachlicher  \newline Einschätzung erforderlich ist,   \newline 1. sich dabei einen unmittelbaren Eindruck  \newline von dem Kind und von seiner persönlichen  \newline Umgebung zu verschaffen sowie   \newline 2. Personen, die gemäß \S{} 4 Absatz 3 des  \newline Gesetzes zur Kooperation und Information  \newline im Kinderschutz dem Jugendamt Daten  übermittelt haben, in geeigneter Weise an  \newline der Gefährdungseinschätzung zu  \newline beteiligen. Hält das Jugendamt zur  \newline Abwendung der Gefährdung die  \newline Gewährung von Hilfen für geeignet und  \newline notwendig, so hat es diese den  \newline Erziehungsberechtigten anzubieten. \\
\hline
Absätze 2 bis 6 &  & unverändert \\
\hline
... & ... & ... \\
\hline
(5) Soweit Leistungen zum Lebensunterhalt  \newline nach \S{} 39 erbracht werden, gehen sie den  \newline Leistungen zum Lebensunterhalt nach \S{} 93  \newline des Vierzehnten Buches vor. & \textbf{(5) Die Leistungen nach diesem Buch  \newline gehen Leistungen nach dem Zwölften  \newline Buch vor. Abweichend von Satz 1 gehen  \newline Leistungen nach }\textbf{\textcolor{diffred}{\S{} 27a Absatz 1 in  \newline Verbindung mit}}\textbf{ \S{} 34 Absatz 6 des  \newline Zwölften Buches den Leistungen nach  \newline diesem Buch vor.}\textbf{\textcolor{diffred}{1}} & \textbf{(5) Die Leistungen nach diesem Buch  \newline gehen Leistungen nach dem Zwölften  \newline Buch vor. Abweichend von Satz 1 gehen  \newline Leistungen }\textbf{\textcolor{diffgreen}{für den Lebensunterhalt   \newline 1.}}\textbf{ nach }\textbf{\textcolor{diffgreen}{dem Dritten Kapitel des Zwölften  \newline Buches zur Deckung der Bedarfe nach  \newline dem Ersten, Zweiten, Vierten und  \newline Fünften Abschnitt, von Leistungen zur  \newline Deckung von Bedarfen des Dritten  \newline Abschnittes nur diejenigen nach}}\textbf{ \S{} 34  \newline Absatz 6 }\textbf{\textcolor{diffgreen}{des Zwölften Buches  und  \newline 2. nach dem Vierten Kapitel}}\textbf{ des  \newline Zwölften Buches   \newline den Leistungen nach diesem Buch vor. } \\
\hline
 & \textbf{\textcolor{diffred}{(6)2 Soweit Leistungen zum  \newline Lebensunterhalt nach \S{} 39 erbracht  \newline werden, gehen sie den Leistungen zum  \newline Lebensunterhalt nach \S{} 93 des  \newline Vierzehnten Buches vor. }} &  \\
\hline
\rowcolor{gray!25}
\textbf{\S{} 10a} &  & \textbf{\S{} 10a} \\
\hline
Beratung &  & Beratung \\
\hline
Absatz 1 &  & unverändert \\
\hline
(2) Die Beratung umfasst insbesondere  \newline 1.   die Familiensituation oder die  \newline persönliche Situation des jungen  \newline Menschen, Bedarfe, vorhandene  \newline Ressourcen sowie mögliche Hilfen,   \newline 2. die Leistungen der Kinder- und  \newline Jugendhilfe einschließlich des Zugangs  \newline zum Leistungssystem,   \newline 3. die Leistungen anderer Leistungsträger,   \newline 4. mögliche Auswirkungen und Folgen  \newline einer Hilfe,   \newline 5. die Verwaltungsabläufe,   \newline 6. Hinweise auf Leistungsanbieter und  \newline andere Hilfemöglichkeiten im Sozialraum  und auf Möglichkeiten zur  \newline Leistungserbringung,   \newline 7. Hinweise auf andere Beratungsangebote  \newline im Sozialraum.   \newline Soweit erforderlich, gehört zur Beratung  \newline auch Hilfe bei der Antragstellung, bei der  \newline Klärung weiterer zuständiger  \newline Leistungsträger, bei der Inanspruchnahme  \newline von Leistungen sowie bei der Erfüllung von  \newline Mitwirkungspflichten. &  & (2) Die Beratung umfasst insbesondere  \newline 1.   die Familiensituation oder die  \newline persönliche Situation des jungen  \newline Menschen, Bedarfe, vorhandene  \newline Ressourcen sowie mögliche Hilfen,   \newline 2. die Leistungen der Kinder- und  \newline Jugendhilfe einschließlich des Zugangs  \newline zum Leistungssystem,   \newline 3. die Leistungen anderer Leistungsträger,   \newline 4. mögliche Auswirkungen und Folgen  \newline einer Hilfe,   \newline 5. die Verwaltungsabläufe,   \newline 6. Hinweise auf Leistungsanbieter und  \newline andere Hilfemöglichkeiten im Sozialraum  und auf Möglichkeiten zur  \newline Leistungserbringung,   \newline 7. Hinweise auf andere Beratungsangebote  \newline im \textcolor{diffgreen}{Sozialraum,} \newline \textbf{\textcolor{diffgreen}{8. eine gebotene Budgetberatung.}} Soweit erforderlich, gehört zur Beratung  \newline auch Hilfe bei der Antragstellung, bei der  \newline Klärung weiterer zuständiger  \newline Leistungsträger, bei der Inanspruchnahme  \newline von Leistungen sowie bei der Erfüllung von  \newline Mitwirkungspflichten. \\
\hline
Absatz 3 &  & unverändert \\
\hline
... & ... & ... \\
\hline
(1) Junge Menschen, die Leistungen der  \newline Eingliederungshilfe wegen einer  \newline Behinderung oder wegen einer drohenden  \newline Behinderung geltend machen oder bei  \newline denen solche Leistungsansprüche in  \newline Betracht kommen, sowie ihre Mütter, Väter,  \newline Personensorge- und  \newline Erziehungsberechtigten haben bei der  \newline Antragstellung, Verfolgung und  \newline Wahrnehmung dieser Leistungen Anspruch  \newline auf Unterstützung und Begleitung durch  \newline einen Verfahrenslotsen. Der  \newline Verfahrenslotse soll die  \newline Leistungsberechtigten bei der  \newline Verwirklichung von Ansprüchen auf  \newline Leistungen der Eingliederungshilfe  \newline unabhängig unterstützen sowie auf die  \newline Inanspruchnahme von Rechten hinwirken.  \newline Diese Leistung wird durch den örtlichen  \newline Träger der öffentlichen Jugendhilfe  \newline erbracht. & (1) Junge Menschen, die Leistungen \textbf{zur  \newline Teilhabe }wegen einer Behinderung oder  \newline wegen einer drohenden Behinderung  \newline geltend machen oder bei denen solche  \newline Leistungsansprüche in Betracht kommen,  \newline sowie ihre Mütter, Väter, Personensorge-  \newline und Erziehungsberechtigten haben bei der  \newline Antragstellung, Verfolgung und  \newline Wahrnehmung dieser Leistungen Anspruch  \newline auf Unterstützung und Begleitung durch  \newline einen Verfahrenslotsen. Der  \newline Verfahrenslotse soll die  \newline Leistungsberechtigten bei der  \newline Verwirklichung von Ansprüchen auf  \newline Leistungen \textbf{zur Teilhabe} \textcolor{diffred}{unabhängig}  \newline unterstützen sowie auf die  \newline Inanspruchnahme von Rechten hinwirken.  \newline \textcolor{diffred}{Diese} Leistung wird durch den örtlichen  \newline Träger der öffentlichen Jugendhilfe  \newline erbracht. & (1) Junge Menschen, die Leistungen \textbf{zur  \newline Teilhabe }wegen einer Behinderung oder  \newline wegen einer drohenden Behinderung  \newline geltend machen oder bei denen solche  \newline Leistungsansprüche in Betracht kommen,  \newline sowie ihre Mütter, Väter, Personensorge-  \newline und Erziehungsberechtigten haben bei der  \newline Antragstellung, Verfolgung und  \newline Wahrnehmung dieser Leistungen Anspruch  \newline auf Unterstützung und Begleitung durch  \newline einen Verfahrenslotsen. Der  \newline Verfahrenslotse soll die  \newline Leistungsberechtigten bei der  \newline Verwirklichung von Ansprüchen auf  \newline Leistungen \textbf{zur Teilhabe} unterstützen  \newline sowie auf die Inanspruchnahme von  \newline Rechten hinwirken. \textbf{\textcolor{diffgreen}{Der Verfahrenslotse  \newline soll auf Wunsch der in Satz 1 genannten  \newline Person auch zu Ansprüchen im Rahmen  \newline der Pflegeversicherung und deren  \newline Inanspruchnahme beraten und die  \newline Anspruchsberechtigten unterstützen.  \newline Die}}\textbf{ Leistung }\textbf{\textcolor{diffgreen}{nach Satz 1 und 2}}\textbf{ wird  \newline durch den örtlichen Träger der  \newline öffentlichen Jugendhilfe }\textbf{\textcolor{diffgreen}{funktionell,  \newline organisatorisch und personell getrennt  \newline von seinen übrigen Aufgaben}}\textbf{ erbracht.} \\
\hline
... & ... & ... \\
\hline
Abschnitt 4:   \newline Hilfe zur Erziehung, Eingliederungshilfe für  \newline seelisch behinderte Kinder und  \newline Jugendliche, Hilfe für junge Volljährige & \textcolor{diffred}{Abschnitt 4:} \newline \textbf{\textcolor{diffred}{Leistungen zur Entwicklung, zur  \newline Erziehung und zur Teilhabe, Hilfen  \newline für junge Volljährige}} &  \\
\hline
Unterabschnitt 1:   \newline Hilfe zur Erziehung & \textcolor{diffred}{Unterabschnitt 1:}\textbf{\textcolor{diffred}{   \newline Leistungen zur Entwicklung, zur  \newline Erziehung und zur Teilhabe  }} &  \\
\hline
\rowcolor{gray!25}
\textbf{\S{} 13} &  & \textbf{\S{} 13} \\
\hline
Jugendsozialarbeit &  & Jugendsozialarbeit \\
\hline
Absätze 1 und 2 &  & unverändert \\
\hline
(3) Jungen Menschen kann während der  \newline Teilnahme an schulischen oder beruflichen  \newline Bildungsmaßnahmen oder bei der  \newline beruflichen Eingliederung Unterkunft in  \newline sozialpädagogisch begleiteten  \newline Wohnformen angeboten werden. In diesen  \newline Fällen sollen auch der notwendige  \newline Unterhalt des jungen Menschen  \newline sichergestellt und Krankenhilfe nach  \newline Maßgabe des \S{} 40 geleistet werden &  & (3) Jungen Menschen kann \textbf{\textcolor{diffgreen}{in Ergänzung  \newline oder unabhängig von Hilfen oder  \newline Maßnahmen nach Absatz 1 oder 2}} während der Teilnahme an schulischen  \newline oder beruflichen Bildungsmaßnahmen oder  \newline bei der beruflichen Eingliederung  \newline Unterkunft in sozialpädagogisch  \newline begleiteten Wohnformen angeboten  \newline werden. In diesen Fällen sollen auch der  \newline notwendige Unterhalt des jungen  \newline Menschen sichergestellt und Krankenhilfe  \newline nach Maßgabe des \S{} 40 geleistet werden \\
\hline
Absatz 4 &  & unverändert \\
\hline
\rowcolor{gray!25}
\textbf{\S{} 18} &  & \textbf{\S{} 18} \\
\hline
Beratung und Unterstützung bei der  \newline Ausübung der Personensorge und des  \newline Umgangsrechts &  & Beratung und Unterstützung bei der  \newline Ausübung der Personensorge und des  \newline Umgangsrechts \\
\hline
Absätze 1 bis 2 &  & unverändert \\
\hline
(3) Kinder und Jugendliche haben  \newline Anspruch auf Beratung und Unterstützung  \newline bei der Ausübung des Umgangsrechts  \newline nach \S{} 1684 Absatz 1 des Bürgerlichen  \newline Gesetzbuchs. Sie sollen darin unterstützt  \newline werden, dass die Personen, die nach  \newline Maßgabe der \S{}\S{} 1684, 1685 und 1686a  \newline des Bürgerlichen Gesetzbuchs zum  \newline Umgang mit ihnen berechtigt sind, von  \newline diesem Recht zu ihrem Wohl Gebrauch  \newline machen. Eltern, andere  \newline Umgangsberechtigte sowie Personen, in  \newline deren Obhut sich das Kind befindet, haben  \newline Anspruch auf Beratung und Unterstützung  \newline bei der Ausübung des Umgangsrechts. Bei  der Befugnis, Auskunft über die  \newline persönlichen Verhältnisse des Kindes zu  \newline verlangen, bei der Herstellung von  \newline Umgangskontakten und bei der Ausführung  \newline gerichtlicher oder vereinbarter  \newline Umgangsregelungen soll vermittelt und in  \newline geeigneten Fällen Hilfestellung geleistet  \newline werden. &  & (3) Kinder und Jugendliche haben  \newline Anspruch auf Beratung und Unterstützung  \newline bei der Ausübung des Umgangsrechts  \newline nach \S{} 1684 Absatz 1 des Bürgerlichen  \newline Gesetzbuchs. Sie sollen darin unterstützt  \newline werden, dass die Personen, die nach  \newline Maßgabe der \S{}\S{} 1684, 1685 und 1686a  \newline des Bürgerlichen Gesetzbuchs zum  \newline Umgang mit ihnen berechtigt sind, von  \newline diesem Recht zu ihrem Wohl Gebrauch  \newline machen. Eltern, andere  \newline Umgangsberechtigte sowie Personen, in  \newline deren Obhut sich das Kind befindet, haben  \newline Anspruch auf Beratung und Unterstützung  \newline bei der Ausübung des Umgangsrechts. Bei  der Befugnis, Auskunft über die  \newline persönlichen Verhältnisse zu verlangen,  \newline bei der Herstellung von Umgangskontakten  \newline und bei der Ausführung gerichtlicher oder  \newline vereinbarter Umgangsregelungen soll  \newline vermittelt und in geeigneten Fällen  \newline Hilfestellung geleistet \textcolor{diffgreen}{werden}\textbf{\textcolor{diffgreen}{; der  \newline begleitete Umgang soll an einem für die  \newline sichere Ausübung des Umgangsrechts  \newline geeigneten Ort stattfinden.}} \\
\hline
Absatz 4 &  & unverändert \\
\hline
Abschnitt 4:   \newline Hilfe zur Erziehung, Eingliederungshilfe für  \newline seelisch behinderte Kinder und  \newline Jugendliche, Hilfe für junge Volljährige &  & Abschnitt 4:   \newline \textbf{\textcolor{diffgreen}{Le}}\textbf{i}\textbf{\textcolor{diffgreen}{stung}}\textbf{e}\textbf{\textcolor{diffgreen}{n}}\textbf{ zur }\textbf{\textcolor{diffgreen}{Entwicklung, zur  \newline Erziehung}}\textbf{ und }\textbf{\textcolor{diffgreen}{zur Teilhabe, Hilfen}}\textbf{  \newline für junge Volljährige} \\
\hline
Unterabschnitt 1:   \newline Hilfe zur Erziehung &  & Unterabschnitt 1: \textbf{\textcolor{diffgreen}{Leistungen zur Entwicklung,}}\textbf{ zur  \newline Erziehung }\textbf{\textcolor{diffgreen}{und zur Teilhabe}} \\
\hline
... & ... & ... \\
\hline
(1) Ein Personensorgeberechtigter hat bei  \newline der Erziehung eines Kindes oder eines  \newline Jugendlichen Anspruch auf Hilfe (Hilfe zur  \newline Erziehung), wenn eine dem Wohl des  \newline Kindes oder des Jugendlichen  \newline entsprechende Erziehung nicht  \newline gewährleistet ist und die Hilfe für seine  \newline Entwicklung geeignet und notwendig ist. & \textbf{(1) Zur Verwirklichung des Rechts eines  \newline jeden jungen Menschen auf Förderung  \newline seiner }\textbf{\textcolor{diffred}{Entwicklung und}}\textbf{ auf Erziehung  \newline zu einer selbstbestimmten,  \newline eigenverantwortlichen und  \newline gemeinschaftsfähigen Persönlichkeit  \newline durch Leistungen zur Entwicklung, zur  \newline Erziehung und zur Teilhabe haben  \newline Kinder, Jugendliche oder  \newline Personensorgeberechtigte einen  \newline Anspruch auf Hilfe zur Erziehung }\textbf{\textcolor{diffred}{un}}\textbf{d  \newline auf Leistungen der Eingliederungshilfe  \newline nach Maßgabe der Absätze 2 und 3. } & \textbf{(1) Zur Verwirklichung des Rechts eines  \newline jeden jungen Menschen auf Förderung  \newline seiner }\textbf{\textcolor{diffgreen}{Entwicklung,}}\textbf{ auf Erziehung zu  \newline einer selbstbestimmten,  \newline eigenverantwortlichen und  \newline gemeinschaftsfähigen Persönlichkeit} \newline \textbf{\textcolor{diffgreen}{und auf Förderung seiner vollen  \newline wirksamen und gleichberechtigten  \newline Teilhabe am Leben in der Gesellschaft}}\textbf{  \newline durch Leistungen zur Entwicklung, zur  \newline Erziehung und zur Teilhabe haben  \newline Kinder, Jugendliche oder  \newline Personensorgeberechtigte einen  \newline Anspruch auf Hilfe zur Erziehung }\textbf{\textcolor{diffgreen}{o}}\textbf{d}\textbf{\textcolor{diffgreen}{er}}\textbf{  \newline auf Leistungen der Eingliederungshilfe  \newline nach Maßgabe der Absätze 2 und 3. } \\
\hline
(2) Hilfe zur Erziehung wird insbesondere  \newline nach Maßgabe der \S{}\S{} 28 bis 35 gewährt.  \newline Art und Umfang der Hilfe richten sich nach  \newline dem erzieherischen Bedarf im Einzelfall;  \newline dabei soll das engere soziale Umfeld des  \newline Kindes oder des Jugendlichen einbezogen  \newline werden. Unterschiedliche Hilfearten  \newline können miteinander kombiniert werden,  \newline sofern dies dem erzieherischen Bedarf des  \newline Kindes oder Jugendlichen im Einzelfall  \newline entspricht. & \textbf{(2) Personensorgeberechtigte haben  \newline einen Anspruch auf Hilfe zur Erziehung,  \newline wenn und solange eine dem Kindeswohl  \newline entsprechende Erziehung nicht  \newline gewährleistet ist und die Hilfe für die  \newline Entwicklung de}\textbf{\textcolor{diffred}{r}}\textbf{ jungen Menschen  \newline geeignet und notwendig ist. }\textbf{\textcolor{diffred}{Liegen die  \newline Voraussetzungen nach Satz 1 vor,  \newline haben auch Jugendliche einen  \newline Anspruch auf Hilfe zur Erziehung, die  \newline außerhalb des Elternhauses erbracht  \newline wird.}} & \textbf{(2) Personensorgeberechtigte haben  \newline einen Anspruch auf Hilfe zur Erziehung,  \newline wenn und solange eine dem Kindeswohl  \newline entsprechende Erziehung nicht } \textbf{gewährleistet ist und die Hilfe für die  \newline Entwicklung de}\textbf{\textcolor{diffgreen}{s}}\textbf{ jungen Menschen  \newline geeignet und notwendig ist.} \\
\hline
... & ... & ... \\
\hline
(3) Hilfe zur Erziehung umfasst  \newline insbesondere die Gewährung  \newline pädagogischer und damit verbundener  \newline therapeutischer Leistungen. Bei Bedarf soll  \newline sie Ausbildungs- und  \newline Beschäftigungsmaßnahmen im Sinne des  \newline \S{} 13 Absatz 2 einschließen und kann mit  \newline anderen Leistungen nach diesem Buch  \newline kombiniert werden. Die in der Schule oder  \newline Hochschule wegen des erzieherischen  \newline Bedarfs erforderliche Anleitung und  \newline Begleitung können als Gruppenangebote  \newline an Kinder oder Jugendliche gemeinsam  \newline erbracht werden, soweit dies dem Bedarf  \newline des Kindes oder Jugendlichen im Einzelfall  \newline entspricht. & \textbf{(3) Kinder oder Jugendliche mit  \newline Behinderungen im Sinne von \S{} 7 Absatz} \newline \textbf{\textcolor{diffred}{2, die in der gleichberechtigten Teilhabe  \newline an der Gesellschaft eingeschränkt sind  \newline oder von einer solchen Behinderung  \newline bedroht sind,}}\textbf{ haben einen Anspruch auf  \newline Leistungen der Eingliederungshilfe,  \newline wenn und solange diese Leistungen  \newline nach der Besonderheit des Einzelfalles  \newline geeignet und notwendig sind, }\textbf{\textcolor{diffred}{den  \newline jungen Menschen eine individuelle  \newline Lebensführung}}\textbf{ zu }\textbf{\textcolor{diffred}{ermöglichen, die der  \newline Würde des Menschen entspricht, ihre  \newline volle, wirksame und gleichberechtigte  \newline Teilhabe am Leben in der Gesellschaft  \newline zu fördern und sie zu befähigen, ihre  \newline Lebensplanung und -führung möglichst  \newline selbstbestimmt und eigenverantwortlich  \newline wahrnehmen zu können.}} & \textbf{(3) Kinder oder Jugendliche mit  \newline Behinderungen }\textbf{\textcolor{diffgreen}{oder von Behinderung  \newline bedrohte Kinder oder Jugendliche}}\textbf{ im  \newline Sinne von \S{} 7 Absatz }\textbf{\textcolor{diffgreen}{2}}\textbf{ haben einen  \newline Anspruch auf Leistungen der  \newline Eingliederungshilfe, wenn und solange  \newline diese Leistungen nach der Besonderheit  \newline des Einzelfalles geeignet und notwendig  \newline sind, }\textbf{\textcolor{diffgreen}{die Aufgaben der  \newline Eingliederungshilfe nach \S{} 90 des  \newline Neunten Buches}}\textbf{ zu }\textbf{\textcolor{diffgreen}{erfüllen.}} \\
\hline
 & \textbf{\textcolor{diffred}{(3a) Maßgeblich für die Eignung und  \newline Notwendigkeit der Leistungen der  \newline Eingliederungshilfe sind insbesondere  \newline die Wechselwirkungen der geistigen,  \newline seelischen, körperlichen oder  \newline Sinnesbeeinträchtigungen mit  \newline einstellungs- und umweltbedingten  \newline Barrieren im Einzelfall und deren  \newline konkrete Auswirkungen auf die Teilhabe  \newline der jungen Menschen an der  \newline Gesellschaft. }} &  \\
\hline
 & \textbf{\textcolor{diffred}{(3b) Von einer Behinderung bedroht  \newline sind Kinder oder Jugendliche, bei denen  \newline der Eintritt einer Behinderung nach  \newline fachlicher Erkenntnis mit hoher  \newline Wahrscheinlichkeit zu erwarten ist. }} &  \\
\hline
(4) Wird ein Kind oder eine Jugendliche  \newline während ihres Aufenthalts in einer  \newline Einrichtung oder einer Pflegefamilie selbst  \newline Mutter eines Kindes, so umfasst die Hilfe  \newline zur Erziehung auch die Unterstützung bei  \newline der Pflege und Erziehung dieses Kindes. & \textbf{(4) }\textbf{\textcolor{diffred}{Die Bundesregierung kann durch  \newline Rechtsverordnung}}\textbf{ mit }\textbf{\textcolor{diffred}{Zustimmung des  \newline Bundesrates Bestimmungen über}}\textbf{ die} \newline \textbf{\textcolor{diffred}{Konkr}}\textbf{e}\textbf{\textcolor{diffred}{t}}\textbf{i}\textbf{\textcolor{diffred}{si}}\textbf{e}\textbf{\textcolor{diffred}{rung}}\textbf{ der} \newline \textbf{\textcolor{diffred}{Leistungsberechtigung nach dieser  \newline Vorschrift erlassen.}} & \textbf{(4) }\textbf{\textcolor{diffgreen}{Maßgeblich für die Eignung und  \newline Notwendigkeit der Leistungen der  \newline Eingliederungshilfe sind die  \newline Wechselwirkungen der geistigen,  \newline seelischen, körperlichen oder  \newline Sinnesbeeinträchtigungen}}\textbf{ mit} \newline \textbf{\textcolor{diffgreen}{einstellungs- und umweltbedingten  \newline Barrieren im Einzelfall und deren  \newline konkrete Auswirkungen auf}}\textbf{ die }\textbf{\textcolor{diffgreen}{T}}\textbf{ei}\textbf{\textcolor{diffgreen}{lhab}}\textbf{e  \newline der }\textbf{\textcolor{diffgreen}{jungen Menschen an der  \newline Gesellschaft. Andere Leistungen der  \newline Eingliederungshilfe können gewährt  \newline werden.}} \\
\hline
 & \textbf{(5) }\textbf{\textcolor{diffred}{Besteht gleichzeitig ein Anspruch auf  \newline Hilfe zur Erziehung}}\textbf{ nach Absatz }\textbf{\textcolor{diffred}{2}}\textbf{ und} \newline \textbf{\textcolor{diffred}{ein Anspruch auf Leistungen der  \newline Eingliederungshilfe nach}}\textbf{ Absatz }\textbf{\textcolor{diffred}{3, so  \newline sollen Einrichtungen, Dienste und  \newline Personen die Hilfe und Leistungen  \newline erbringen, die geeignet sind, sowohl  \newline den erzieherischen Bedarf zu decken als  \newline auch die Aufgaben der  \newline Eingliederungshilfe zu erfüllen.}} & \textbf{(5) }\textbf{\textcolor{diffgreen}{Geeignete Leistungen können  \newline gewährt werden, wenn die  \newline Voraussetzungen der Notwendigkeit der  \newline Leistungen}}\textbf{ nach Absatz }\textbf{\textcolor{diffgreen}{3}}\textbf{ und Absatz }\textbf{\textcolor{diffgreen}{4  \newline nicht vorliegen.}} \\
\hline
 &  & \textbf{\textcolor{diffgreen}{(6) Besteht gleichzeitig ein Anspruch auf  \newline Hilfe zur Erziehung nach Absatz 2 und  \newline ein Anspruch auf Leistungen der  \newline Eingliederungshilfe nach Absatz 3, so  \newline sollen Einrichtungen, Dienste und  \newline Personen die Hilfe und Leistungen  \newline erbringen, die geeignet sind, sowohl  \newline den erzieherischen Bedarf zu decken als  \newline auch die Aufgaben der  \newline Eingliederungshilfe zu erfüllen.}} \\
\hline
... & ... & ... \\
\hline
 & \textbf{(2) Ist eine Erziehung des Kindes oder  \newline Jugendlichen außerhalb des  \newline Elternhauses erforderlich, so entfällt der  \newline Anspruch auf Hilfe zur Erziehung nicht  \newline dadurch, dass eine andere  \newline unterhaltspflichtige Person bereit ist,  \newline diese Aufgabe zu }\textbf{\textcolor{diffred}{übernehmen; die}}\textbf{  \newline Gewährung von Hilfe zur Erziehung  \newline setzt in diesem Fall voraus, dass diese  \newline Person bereit und geeignet ist, den  \newline Hilfebedarf in Zusammenarbeit mit dem  \newline Träger der öffentlichen Jugendhilfe  \newline nach Maßgabe der \S{}\S{} 36, 36a und 3}\textbf{\textcolor{diffred}{8}}\textbf{ zu  \newline decken. } & \textbf{(2)} \textbf{Ist eine Erziehung des Kindes oder  \newline Jugendlichen außerhalb des  \newline Elternhauses erforderlich, so entfällt der  \newline Anspruch auf Hilfe zur Erziehung nicht  \newline dadurch, dass eine andere  \newline unterhaltspflichtige Person bereit ist,  \newline diese Aufgabe zu }\textbf{\textcolor{diffgreen}{übernehmen. Die}}\textbf{  \newline Gewährung von Hilfe zur Erziehung  \newline setzt in diesem Fall voraus, dass diese  \newline Person bereit und geeignet ist, den  \newline Hilfebedarf in Zusammenarbeit mit dem  \newline Träger der öffentlichen Jugendhilfe  \newline nach Maßgabe der \S{}\S{} 36, 36a und 3}\textbf{\textcolor{diffgreen}{9}}\textbf{ zu  \newline decken.  } \\
\hline
 & \textbf{(3) Hilfe zur Erziehung umfasst  \newline insbesondere die Gewährung  \newline pädagogischer und damit verbundener  \newline therapeutischer Leistungen. Bei Bedarf  \newline soll sie Ausbildungs- und  \newline Beschäftigungsmaßnahmen im Sinne  \newline des \S{} 13 Absatz 2 einschließen und  \newline kann mit anderen Leistungen nach  \newline diesem Buch kombiniert werden. }\textbf{\textcolor{diffred}{Die in  \newline der Schule oder Hochschule wegen des  \newline erzieherischen Bedarfs erforderliche  \newline Anleitung und Begleitung können als  \newline Gruppenangebote an Kinder oder  \newline Jugendliche gemeinsam erbracht  \newline werden, soweit dies dem Bedarf des  \newline Kindes oder Jugendlichen im Einzelfall  \newline entspricht, für diesen nach \S{} 5 Absatz 3  \newline zumutbar ist und mit  \newline Leistungserbringern entsprechende  \newline Vereinbarungen bestehen. Die  \newline Leistungen nach Satz 1 sind auf  \newline Wunsch des Leistungsberechtigten  \newline gemeinsam zu erbringen.}} & \textbf{(3) Hilfe zur Erziehung umfasst  \newline insbesondere die Gewährung  \newline pädagogischer und damit verbundener  \newline therapeutischer Leistungen. Bei Bedarf  \newline soll sie Ausbildungs- und  \newline Beschäftigungsmaßnahmen im Sinne  \newline des \S{} 13 Absatz 2 einschließen und  \newline kann mit anderen Leistungen nach  \newline diesem Buch kombiniert werden.  } \\
\hline
 & \textbf{(4) }\textbf{\textcolor{diffred}{Wird ein Kind}}\textbf{ oder }\textbf{\textcolor{diffred}{eine Jugendliche  \newline während ihres Aufenthalts in einer  \newline Einrichtung}}\textbf{ oder }\textbf{\textcolor{diffred}{einer Pflegefamilie  \newline selbst Mutter eines Kindes, so umfasst  \newline die Hilfe zur Erziehung auch die  \newline Unterstützung bei der Pflege}}\textbf{ und} \newline \textbf{\textcolor{diffred}{Erziehung dieses Kindes.}} & \textbf{(4) }\textbf{\textcolor{diffgreen}{Sofern infrastrukturelle Angebote}}\textbf{  \newline oder }\textbf{\textcolor{diffgreen}{Regelangebote insbesondere nach  \newline \S{}\S{} 16 bis 18, \S{}\S{} 22 bis 25}}\textbf{ oder }\textbf{\textcolor{diffgreen}{\S{} 13 im  \newline Hinblick auf den Bedarf des Kindes oder  \newline des Jugendlichen im Einzelfall  \newline geeigneter oder gleichermaßen geeignet  \newline sind, werden diese vorrangig gewährt.  \newline Bei Jugendlichen}}\textbf{ und }\textbf{\textcolor{diffgreen}{jungen  \newline Volljährigen werden Hilfen oder  \newline Maßnahmen nach \S{} 13 vorrangig  \newline gewährt, wenn sie gleichermaßen  \newline geeignet sind.}} \\
\hline
 &  & \textbf{\textcolor{diffgreen}{(5)}} \textbf{\textcolor{diffgreen}{Die in einer Tageseinrichtung für  \newline Kinder nach \S{} 22 Absatz 1 Satz 1 oder in  \newline der Schule oder Hochschule wegen des  \newline erzieherischen Bedarfs erforderliche  \newline Anleitung und Begleitung werden als  \newline infrastrukturelle Angebote nach \S{} 80a  \newline gewährt; die Vorschriften zur Hilfe- und  \newline Leistungsplanung (\S{}\S{} 36 bis 38d) finden  \newline keine Anwendung. Nur wenn dem  \newline erzieherischen Bedarf im Einzelfall  \newline ausschließlich durch eine an dem  \newline jeweiligen Kind oder Jugendlichen  \newline erbrachte Anleitung und Begleitung  in  \newline einer Tageseinrichtung für Kinder nach  \newline \S{} 22 Absatz 1 Satz 1 oder in der Schule  \newline oder Hochschule entsprochen werden  \newline kann, besteht nach \S{} 27 Absatz 2 ein  \newline Anspruch auf diese Einzelhilfe.}} \\
\hline
 &  & \textbf{\textcolor{diffgreen}{(6)}} \textbf{\textcolor{diffgreen}{Wird ein Kind oder eine Jugendliche  \newline während des Aufenthalts in einer  \newline Einrichtung oder einer Pflegefamilie  \newline selbst Mutter eines Kindes, so umfasst  \newline die Hilfe zur Erziehung auch die  \newline Unterstützung bei der Pflege und  \newline Erziehung dieses Kindes.}} \\
\hline
... & ... & ... \\
\hline
(1) Kinder oder Jugendliche haben  \newline Anspruch auf Eingliederungshilfe, wenn & \textbf{(1) Besteht ein Anspruch auf Leistungen  \newline der Eingliederungshilfe nach \S{} 27  \newline Absatz 3, werden diese insbesondere  \newline nach Maßgabe der Absätze 2 }\textbf{\textcolor{diffred}{und 3}}\textbf{  \newline sowie der \S{}\S{} 35b bis 35i und }\textbf{\textcolor{diffred}{den  \newline Kapiteln}}\textbf{ 9 bis 13 des Teils 1 des  \newline Neunten Buches gewährt; die Kapitel 3  \newline bis 6 des Teils 2 des Neunten Buches  \newline gelten im Übrigen entsprechend. Art  \newline und Umfang der Leistungen richten sich  \newline nach dem Ergebnis der Prüfung gemäß  \newline \S{} 27 Absatz }\textbf{\textcolor{diffred}{3a}}\textbf{ und bestimmen sich  \newline nach der Besonderheit des Einzelfalles,  \newline insbesondere nach dem individuellen  \newline Bedarf, den persönlichen Verhältnissen,  \newline dem engeren sozialen Umfeld, dem  \newline Sozialraum und den eigenen Kräften  \newline und Mitteln; dabei ist auch die  \newline Wohnform zu würdigen.  \newline Unterschiedliche Leistungsarten  \newline können miteinander kombiniert werden,  \newline sofern dies dem Bedarf des Kindes oder  \newline Jugendlichen im Einzelfall entspricht. } & \textbf{(1) Besteht ein Anspruch auf Leistungen  \newline der Eingliederungshilfe nach \S{} 27  \newline Absatz 3, werden diese insbesondere  \newline nach Maßgabe der Absätze 2 }\textbf{\textcolor{diffgreen}{bis 7}}\textbf{  \newline sowie der \S{}\S{} 35b bis 35i und }\textbf{\textcolor{diffgreen}{der Kapitel}}\textbf{  \newline 9 bis 13 des Teils 1 des Neunten Buches  \newline gewährt; }\textbf{\textcolor{diffgreen}{\S{} 107 sowie}}\textbf{ die Kapitel 3 bis 6 } \textbf{des Teils 2 des Neunten Buches gelten  \newline im Übrigen entsprechend. Art und  \newline Umfang der Leistungen richten sich  \newline nach dem Ergebnis der Prüfung gemäß  \newline \S{} 27 Absatz }\textbf{\textcolor{diffgreen}{4}}\textbf{ und bestimmen sich nach  \newline der Besonderheit des Einzelfalles,  \newline insbesondere nach dem individuellen  \newline Bedarf, den persönlichen Verhältnissen,  \newline dem engeren sozialen Umfeld, dem  \newline Sozialraum und den eigenen Kräften  \newline und Mitteln; dabei ist auch die  \newline Wohnform zu würdigen.  \newline Unterschiedliche Leistungsarten }\textbf{\textcolor{diffgreen}{der  \newline Eingliederungshilfe}}\textbf{ können miteinander  \newline kombiniert werden, sofern dies dem  \newline Bedarf des Kindes oder Jugendlichen  \newline im Einzelfall entspricht. } \\
\hline
... & ... & ... \\
\hline
 & \textbf{\textcolor{diffred}{5.   \newline in Einrichtungen über Tag und Nacht  \newline sowie sonstigen Wohnformen erbracht.  \newline Dabei sollen Einrichtungen, Dienste  \newline oder Personen die Leistungen erbringen  \newline oder in Anspruch genommen werden, in  \newline denen Kinder oder Jugendliche mit  \newline Behinderungen und Kinder oder  \newline Jugendliche ohne Behinderungen  \newline gemeinsam Leistungen erhalten  \newline können, wenn die Aufgaben der  \newline Eingliederungshilfe erfüllt werden  \newline können; die besonderen Bedürfnisse  \newline von Kindern oder Jugendlichen mit  \newline Behinderungen und von Kindern oder  \newline Jugendlichen, die von einer  \newline Behinderung bedroht sind, sind zu  \newline berücksichtigen. }} &  \\
\hline
... & ... & ... \\
\hline
 & \textbf{\textcolor{diffred}{4}}\textbf{. durch geeignete Pflegepersonen und } & \textbf{\textcolor{diffgreen}{3}}\textbf{. durch geeignete Pflegepersonen und } \\
\hline
 & \textbf{\textcolor{diffred}{5}}\textbf{. in Einrichtungen über Tag und Nacht  \newline sowie sonstigen Wohnformen erbracht.  \newline Dabei sollen Einrichtungen, Dienste  \newline oder Personen die Leistungen erbringen  \newline oder in Anspruch genommen werden, in  \newline denen Kinder oder Jugendliche mit  \newline Behinderungen und Kinder oder  \newline Jugendliche ohne Behinderungen  \newline gemeinsam Leistungen erhalten  \newline können, wenn die Aufgaben der  \newline Eingliederungshilfe erfüllt werden  \newline können; die besonderen Bedürfnisse  \newline von Kindern oder Jugendlichen mit  \newline Behinderungen und von Kindern oder  \newline Jugendlichen, die von einer  \newline Behinderung bedroht sind, sind zu  \newline berücksichtigen. } & \textbf{\textcolor{diffgreen}{4}}\textbf{. in Einrichtungen über Tag und Nacht  \newline sowie sonstigen Wohnformen erbracht.   \newline Dabei sollen Einrichtungen, Dienste  \newline oder Personen die Leistungen erbringen  \newline oder in Anspruch genommen werden, in  \newline denen Kinder oder Jugendliche mit  \newline Behinderungen und Kinder oder  \newline Jugendliche ohne Behinderungen  \newline gemeinsam Leistungen erhalten  \newline können, wenn die Aufgaben der  \newline Eingliederungshilfe erfüllt werden  \newline können; die besonderen Bedürfnisse  \newline von Kindern oder Jugendlichen mit  \newline Behinderungen und von Kindern oder  \newline Jugendlichen, die von einer  \newline Behinderung bedroht sind, sind zu  \newline berücksichtigen. } \\
\hline
... & ... & ... \\
\hline
 & \textbf{(6) Die Leistungen der  \newline Eingliederungshilfe werden auf Antrag  \newline auch }\textbf{\textcolor{diffred}{als Teil eines Persönlichen  \newline Budgets}}\textbf{ ausgeführt. }\textbf{\textcolor{diffred}{Die Betroffenen}}\textbf{  \newline sind entsprechend zu beraten. Die  \newline Vorschrift zum Persönlichen Budget  \newline nach \S{} 29 des Neunten Buches ist  \newline insoweit anzuwenden. } & \textbf{(6) Die Leistungen der  \newline Eingliederungshilfe werden auf Antrag  \newline auch }\textbf{\textcolor{diffgreen}{durch ein Persönliches Budget}}\textbf{  \newline ausgeführt. }\textbf{\textcolor{diffgreen}{Der Leistungsberechtigte  \newline und der Personensorgeberechtigte}}\textbf{ sind  \newline entsprechend zu beraten. Die Vorschrift  \newline zum Persönlichen Budget nach \S{} 29 des  \newline Neunten Buches ist insoweit  \newline anzuwenden.  } \\
\hline
 &  & \textbf{\textcolor{diffgreen}{(7) \S{} 103 des Neunten Buches gilt  \newline entsprechend.}} \\
\hline
... & ... & ... \\
\hline
 & \textbf{(3) Leistungsberechtigte haben  \newline entsprechend den Bestimmungen der  \newline gesetzlichen Krankenversicherung die  \newline freie Wahl unter den Ärztinnen oder  \newline Ärzten und }\textbf{\textcolor{diffred}{Zahnärztinnen oder}}\textbf{  \newline Zahnärzten sowie unter den  \newline Krankenhäusern und Vorsorge- und  \newline Rehabilitationseinrichtungen. } & \textbf{(3) Leistungsberechtigte haben  \newline entsprechend den Bestimmungen der  \newline gesetzlichen Krankenversicherung die  \newline freie Wahl unter den Ärztinnen oder  \newline Ärzten und Zahnärzten sowie unter den  \newline Krankenhäusern und Vorsorge- und  \newline Rehabilitationseinrichtungen. } \\
\hline
 & \textbf{(4) Bei der Erbringung von Leistungen  \newline zur medizinischen Rehabilitation sind  \newline die Regelungen, die für die gesetzlichen  \newline Krankenkassen nach dem Vierten  \newline Kapitel des Fünften Buches gelten, mit  \newline Ausnahme des Dritten Titels des  \newline Zweiten Abschnitts anzuwenden.} \newline \textbf{\textcolor{diffred}{Ärztinnen oder}}\textbf{ Ärzte,} \newline \textbf{\textcolor{diffred}{Psychotherapeutinnen oder}}\textbf{  \newline Psychotherapeuten im Sinne des \S{} 28  \newline Absatz 3 Satz 1 des Fünften Buches und} \newline \textbf{\textcolor{diffred}{Zahnärztinnen oder}}\textbf{ Zahnärzte haben für  \newline ihre Leistungen Anspruch auf die  \newline Vergütung, welche die  \newline Ortskrankenkasse, in deren Bereich }\textbf{\textcolor{diffred}{die  \newline Ärztin oder}}\textbf{ der Arzt, }\textbf{\textcolor{diffred}{die  \newline Psychotherapeutin oder}}\textbf{ der  \newline Psychotherapeut }\textbf{\textcolor{diffred}{oder die Zahnärztin}}\textbf{  \newline oder der Zahnarzt niedergelassen ist,  \newline für ihre Mitglieder zahlt. } & \textbf{(4) Bei der Erbringung von Leistungen  \newline zur medizinischen Rehabilitation sind  \newline die Regelungen, die für die gesetzlichen  \newline Krankenkassen nach dem Vierten  \newline Kapitel des Fünften Buches gelten, mit  \newline Ausnahme des Dritten Titels des  \newline Zweiten Abschnitts anzuwenden. Ärzte, } \textbf{Psychotherapeuten im Sinne des \S{} 28  \newline Absatz 3 Satz 1 des Fünften Buches und  \newline Zahnärzte haben für ihre Leistungen  \newline Anspruch auf die Vergütung, welche die  \newline Ortskrankenkasse, in deren Bereich der  \newline Arzt, der Psychotherapeut oder der  \newline Zahnarzt niedergelassen ist, für ihre  \newline Mitglieder zahlt. } \\
\hline
... & ... & ... \\
\hline
 & \textbf{(1) Früherkennung und Frühförderung  \newline für Kinder mit Behinderungen und von  \newline Behinderung}\textbf{\textcolor{diffred}{en}}\textbf{ bedrohte Kinder sollen  \newline auf der Grundlage eines ganzheitlichen  \newline und interdisziplinären Konzepts unter  \newline Berücksichtigung und Einbeziehung des  \newline engeren sozialen Umfelds des Kindes  \newline eine drohende oder bereits eingetretene  \newline Behinderung zum frühestmöglichen  \newline Zeitpunkt erkennen }\textbf{\textcolor{diffred}{o}}\textbf{d}\textbf{\textcolor{diffred}{er}}\textbf{ die  \newline Behinderung durch gezielte Förder- und  \newline Behandlungsmaßnahmen ausgleichen  \newline oder mildern. Die Leistungen der  \newline Früherkennung und Frühförderung  \newline bestimmen sich nach \S{}\S{} 42 Absatz 2  \newline Nummer 2, 46 des Neunten }\textbf{\textcolor{diffred}{Buches.}} & \textbf{(1) Früherkennung und Frühförderung  \newline für Kinder mit Behinderungen und von  \newline Behinderung bedrohte Kinder sollen auf  \newline der Grundlage eines ganzheitlichen und  \newline interdisziplinären Konzepts unter  \newline Berücksichtigung und Einbeziehung des  \newline engeren sozialen Umfelds des Kindes  \newline eine drohende oder bereits eingetretene  \newline Behinderung zum frühestmöglichen  \newline Zeitpunkt erkennen }\textbf{\textcolor{diffgreen}{un}}\textbf{d die  \newline Behinderung durch gezielte Förder- und  \newline Behandlungsmaßnahmen ausgleichen  \newline oder mildern. Die Leistungen der  \newline Früherkennung und Frühförderung  \newline bestimmen sich nach \S{}\S{} 42 Absatz 2  \newline Nummer 2, 46 }\textbf{\textcolor{diffgreen}{und 79}}\textbf{ des Neunten} \newline \textbf{\textcolor{diffgreen}{Buches; \S{} 27 Absatz 3 findet insoweit  \newline keine Anwendung.}} \\
\hline
 & \textbf{(2) Die Vorschriften der }\textbf{\textcolor{diffred}{Verordnung zur  \newline Früherkennung und Frühförderung  \newline behinderter und von Behinderung  \newline bedrohter Kinder}}\textbf{ finden Anwendung. }\textbf{\textcolor{diffred}{In  \newline Verbindung mit schulvorbereitenden  \newline Maßnahmen der Schulträger werden die  \newline Leistungen ebenfalls als  \newline Komplexleistung erbracht.}} & \textbf{(2) Die Vorschriften der} \newline \textbf{\textcolor{diffgreen}{Frühförderungsverordnung}}\textbf{ finden  \newline Anwendung.  } \\
\hline
 & \textbf{(3) Die Vorschriften zur Hilfe- und  \newline Leistungsplanung (\S{}\S{} 36 bis 38d) finden  \newline bei Leistungen }\textbf{\textcolor{diffred}{de}}\textbf{r Früherkennung und  \newline Frühförderung keine Anwendung. } & \textbf{(3) Die Vorschriften zur Hilfe- und  \newline Leistungsplanung (\S{}\S{} 36 bis 38d) finden  \newline bei Leistungen }\textbf{\textcolor{diffgreen}{zu}}\textbf{r Früherkennung und  \newline Frühförderung keine Anwendung. }\textbf{\textcolor{diffgreen}{An }} \textbf{\textcolor{diffgreen}{die Stelle der Hilfe- und  \newline Leistungsplanung und des Hilfe- und  \newline Leistungsplans nach den \S{}\S{} 36 bis 38d  \newline tritt der Förder- und Behandlungsplan  \newline nach \S{} 7 der  \newline Frühförderungsverordnung.}} \\
\hline
... & ... & ... \\
\hline
 & \textbf{2. Hilfen zur schulischen oder  \newline hochschulischen Ausbildung oder  \newline Weiterbildung für einen Beruf. Die Hilfen  \newline nach Satz 1 Nummer 1 schließen  \newline Leistungen zur Unterstützung  \newline schulischer Ganztagsangebote in der  \newline offenen Form ein, die im Einklang mit  \newline dem Bildungs- und Erziehungsauftrag  \newline der Schule stehen und unter deren  \newline Aufsicht und Verantwortung ausgeführt  \newline werden, an den stundenplanmäßigen  \newline Unterricht anknüpfen und in der Regel  \newline in den Räumlichkeiten der Schule oder  \newline in deren Umfeld durchgeführt werden.  \newline Hilfen nach Satz 1 Nummer 1 umfassen  \newline auch heilpädagogische und sonstige  \newline Maßnahmen, wenn die Maßnahmen  \newline erforderlich und geeignet sind, der  \newline leistungsberechtigten Person den  \newline Schulbesuch zu ermöglichen oder zu  \newline erleichtern. Hilfen zu einer schulischen  \newline oder hochschulischen Ausbildung nach  \newline Satz 1 Nummer 2 können erneut  \newline erbracht werden, wenn dies aus  \newline behinderungsbedingten Gründen  \newline erforderlich ist. Hilfen nach Satz 1  \newline umfassen auch Gegenstände und  \newline Hilfsmittel, die wegen der  \newline gesundheitlichen Beeinträchtigung zur  \newline Teilhabe an Bildung erforderlich sind.  \newline Voraussetzung für eine  \newline Hilfsmittelversorgung ist, dass die  \newline leistungsberechtigte Person das  \newline Hilfsmittel bedienen kann. Die  \newline Versorgung mit Hilfsmitteln schließt  \newline eine notwendige Unterweisung im  \newline Gebrauch und eine notwendige  \newline Instandhaltung oder Änderung ein. Die  \newline Ersatzbeschaffung des Hilfsmittels  \newline erfolgt, wenn sie infolge der  \newline körperlichen Entwicklung }\textbf{\textcolor{diffred}{der  \newline leistungsberechtigten Person}}\textbf{ notwendig  \newline ist oder wenn das Hilfsmittel aus  \newline anderen Gründen ungeeignet oder  \newline unbrauchbar geworden ist. } & \textbf{2. Hilfen zur schulischen oder  \newline hochschulischen Ausbildung oder  \newline Weiterbildung für einen Beruf.   \newline Die Hilfen nach Satz 1 Nummer 1  \newline schließen Leistungen zur Unterstützung  \newline schulischer Ganztagsangebote in der  \newline offenen Form ein, die im Einklang mit  \newline dem Bildungs- und Erziehungsauftrag  \newline der Schule stehen und unter deren  \newline Aufsicht und Verantwortung ausgeführt  \newline werden, an den stundenplanmäßigen  \newline Unterricht anknüpfen und in der Regel  \newline in den Räumlichkeiten der Schule oder  \newline in deren Umfeld durchgeführt werden.  \newline Hilfen nach Satz 1 Nummer 1 umfassen  \newline auch heilpädagogische und sonstige  \newline Maßnahmen, wenn die Maßnahmen  \newline erforderlich und geeignet sind, der  \newline leistungsberechtigten Person den  \newline Schulbesuch zu ermöglichen oder zu  \newline erleichtern. Hilfen zu einer schulischen  \newline oder hochschulischen Ausbildung nach  \newline Satz 1 Nummer 2 können erneut  \newline erbracht werden, wenn dies aus  \newline behinderungsbedingten Gründen  \newline erforderlich ist. Hilfen nach Satz 1  \newline umfassen auch Gegenstände und } \textbf{Hilfsmittel, die wegen der  \newline gesundheitlichen Beeinträchtigung zur  \newline Teilhabe an Bildung erforderlich sind.  \newline Voraussetzung für eine  \newline Hilfsmittelversorgung ist, dass die  \newline leistungsberechtigte Person das  \newline Hilfsmittel bedienen kann. Die  \newline Versorgung mit Hilfsmitteln schließt  \newline eine notwendige Unterweisung im  \newline Gebrauch und eine notwendige  \newline Instandhaltung oder Änderung ein. Die  \newline Ersatzbeschaffung des Hilfsmittels  \newline erfolgt, wenn sie infolge der  \newline körperlichen Entwicklung }\textbf{\textcolor{diffgreen}{des  \newline Leistungsberechtigten}}\textbf{ notwendig ist  \newline oder wenn das Hilfsmittel aus anderen  \newline Gründen ungeeignet oder unbrauchbar  \newline geworden ist. } \\
\hline
 & \textbf{(2) Hilfen nach Absatz 1 Satz 1 Nummer  \newline 2 werden erbracht für eine schulische  \newline oder hochschulische berufliche  \newline Weiterbildung, die in einem zeitlichen  \newline Zusammenhang an eine duale,  \newline schulische oder hochschulische  \newline Berufsausbildung anschließt, in  \newline dieselbe fachliche Richtung weiterführt  \newline und es dem Leistungsberechtigten  \newline ermöglicht, das von ihm angestrebte  \newline Berufsziel zu erreichen. Hilfen für ein  \newline Masterstudium werden abweichend von  \newline Satz 1 Nummer 2 auch erbracht, wenn  \newline das Masterstudium auf ein zuvor  \newline abgeschlossenes Bachelorstudium  \newline aufbaut und dieses interdisziplinär  \newline ergänzt, ohne in dieselbe Fachrichtung  \newline weiterzuführen. Aus  \newline behinderungsbedingten oder aus  \newline anderen, nicht von der  \newline leistungsberechtigten Person  \newline beeinflussbaren gewichtigen Gründen  \newline kann von Satz 1 Nummer 1 abgewichen  \newline werden. } & \textbf{(2) Hilfen nach Absatz 1 Satz 1 Nummer  \newline 2 werden erbracht für eine schulische  \newline oder hochschulische berufliche  \newline Weiterbildung, die} \newline \textbf{\textcolor{diffgreen}{1.}}\textbf{ in einem zeitlichen Zusammenhang  \newline an eine duale, schulische oder  \newline hochschulische Berufsausbildung  \newline anschließt,} \newline \textbf{\textcolor{diffgreen}{2.}}\textbf{ in dieselbe fachliche Richtung  \newline weiterführt und} \newline \textbf{\textcolor{diffgreen}{3.}}\textbf{ es dem Leistungsberechtigten  \newline ermöglicht, das von ihm angestrebte  \newline Berufsziel zu erreichen.   \newline Hilfen für ein Masterstudium werden  \newline abweichend von Satz 1 Nummer 2 auch  \newline erbracht, wenn das Masterstudium auf  \newline ein zuvor abgeschlossenes  \newline Bachelorstudium aufbaut und dieses  \newline interdisziplinär ergänzt, ohne in  \newline dieselbe Fachrichtung weiterzuführen.  \newline Aus behinderungsbedingten oder aus  \newline anderen, nicht von der  \newline leistungsberechtigten Person  \newline beeinflussbaren gewichtigen Gründen  \newline kann von Satz 1 Nummer 1 abgewichen  \newline werden. } \\
\hline
... & ... & ... \\
\hline
 & \textbf{(4) Die in der Schule oder Hochschule  \newline wegen der Behinderung erforderliche  \newline Anleitung und Begleitung }\textbf{\textcolor{diffred}{könn}}\textbf{en als} \newline \textbf{\textcolor{diffred}{Gruppenangebote}}\textbf{ an }\textbf{\textcolor{diffred}{Kinder und  \newline Jugendliche gemeinsam erbracht  \newline werden, soweit dies}}\textbf{ dem }\textbf{\textcolor{diffred}{Bedarf des  \newline Kindes}}\textbf{ oder Jugendlichen }\textbf{\textcolor{diffred}{im Einzelfall  \newline entspricht, für diesen}}\textbf{ nach \S{} }\textbf{\textcolor{diffred}{5}}\textbf{ Absatz 3} \newline \textbf{\textcolor{diffred}{zumutbar ist und mit  \newline Leistungserbringern entsprechende  \newline Vereinbarungen bestehen. Die  \newline Leistungen nach Satz 1 sind}}\textbf{ auf} \newline \textbf{\textcolor{diffred}{Wunsch der Leistungsberechtigten  \newline gemeinsam zu erbringen.}} & \textbf{(4) Die in der Schule oder Hochschule  \newline wegen der Behinderung erforderliche  \newline Anleitung und Begleitung }\textbf{\textcolor{diffgreen}{werd}}\textbf{en als} \newline \textbf{\textcolor{diffgreen}{infrastrukturelle Angebote nach \S{} 80a  \newline gewährt; die Vorschriften zur Hilfe- und  \newline Leistungsplanung (\S{}\S{} 36 bis 38d) finden  \newline keine Anwendung. Nur wenn dem  \newline Ergebnis der Prüfung gemäß \S{} 27  \newline Absatz 4 und der Besonderheit des  \newline Einzelfalles ausschließlich durch eine}}\textbf{  \newline an dem }\textbf{\textcolor{diffgreen}{jeweiligen Kind}}\textbf{ oder  \newline Jugendlichen }\textbf{\textcolor{diffgreen}{erbrachte Anleitung und  \newline Begleitung in der Schule oder  \newline Hochschule entsprochen werden kann,  \newline besteht}}\textbf{ nach \S{} }\textbf{\textcolor{diffgreen}{27}}\textbf{ Absatz 3 }\textbf{\textcolor{diffgreen}{ein  \newline Anspruch}}\textbf{ auf }\textbf{\textcolor{diffgreen}{diese Einzelhilfe.}} \\
\hline
... & ... & ... \\
\hline
 & \textbf{2. Leistungen bei anderen  \newline Leistungsanbietern nach den \S{}\S{} 6}\textbf{\textcolor{diffred}{0}}\textbf{ und  \newline 62 des Neunten Buches, } & \textbf{2. Leistungen bei anderen  \newline Leistungsanbietern nach den \S{}\S{} 6}\textbf{\textcolor{diffgreen}{1}}\textbf{ und  \newline 62 des Neunten Buches, } \\
\hline
... & ... & ... \\
\hline
 & \textbf{(4) }\textbf{\textcolor{diffred}{Zur Ermöglichung}}\textbf{ der} \newline \textbf{\textcolor{diffred}{gemeinschaftlichen Mittagsverpflegung}}\textbf{  \newline in }\textbf{\textcolor{diffred}{der Verantwortung}}\textbf{ einer }\textbf{\textcolor{diffred}{W}}\textbf{er}\textbf{\textcolor{diffred}{ks}}\textbf{t}\textbf{\textcolor{diffred}{att}}\textbf{ für} \newline \textbf{\textcolor{diffred}{behinderte Menschen, einem anderen  \newline Leistungsanbieter oder dem  \newline Leistungserbringer vergleichbarer  \newline anderer tagesstrukturierender  \newline Maßnahmen}}\textbf{ werden }\textbf{\textcolor{diffred}{die erforderliche  \newline sächliche Ausstattung, die personelle  \newline Ausstattung und die erforderlichen  \newline betriebsnotwendigen Anlagen des  \newline Leistungserbringers übernommen.}} & \textbf{(4) }\textbf{\textcolor{diffgreen}{Die in einer Tageseinrichtung für  \newline Kinder nach \S{} 22 Absatz 1 Satz 1 wegen}}\textbf{  \newline der }\textbf{\textcolor{diffgreen}{Behinderung erforderliche Anleitung  \newline und Begleitung werden als  \newline infrastrukturelle Angebote nach \S{} 80a }} \textbf{\textcolor{diffgreen}{gewährt; die Vorschriften zur Hilfe- und  \newline Leistungsplanung (\S{}\S{} 36 bis 38d) finden  \newline keine Anwendung. Nur wenn dem  \newline Ergebnis der Prüfung gemäß \S{} 27  \newline Absatz 4 und der Besonderheit des  \newline Einzelfalles ausschließlich durch eine  \newline an dem jeweiligen Kind erbrachte  \newline Anleitung und Begleitung}}\textbf{  in einer} \newline \textbf{\textcolor{diffgreen}{Tag}}\textbf{e}\textbf{\textcolor{diffgreen}{sein}}\textbf{r}\textbf{\textcolor{diffgreen}{ich}}\textbf{t}\textbf{\textcolor{diffgreen}{ung}}\textbf{ für }\textbf{\textcolor{diffgreen}{Kinder nach \S{} 22  \newline Absatz 1 Satz 1 entsprochen}}\textbf{ werden} \newline \textbf{\textcolor{diffgreen}{kann, besteht nach \S{} 27 Absatz 2 ein  \newline Anspruch auf diese Einzelhilfe.}} \\
\hline
 & \textbf{(5) }\textbf{\textcolor{diffred}{In besonderen Wohnformen}}\textbf{ des }\textbf{\textcolor{diffred}{\S{}  \newline 42a Absatz 2 Satz 1 Nummer 2 und Satz  \newline 3 des Zwölften Buches werden  \newline Aufwendungen für Wohnraum oberhalb  \newline der Angemessenheitsgrenze nach \S{} 42a  \newline Absatz 6 des Zwölften Buches  \newline übernommen, sofern dies wegen der  \newline besonderen Bedürfnisse des Kindes  \newline oder Jugendlichen mit Behinderungen  \newline erforderlich ist. \S{}\S{} 78a bis 78g sind  \newline anzuwenden.}} & \textbf{(5) }\textbf{\textcolor{diffgreen}{Zur Ermöglichung der  \newline gemeinschaftlichen Mittagsverpflegung  \newline in der Verantwortung einer Werkstatt für  \newline behinderte Menschen, einem anderen  \newline Leistungsanbieter oder dem  \newline Leistungserbringer vergleichbarer  \newline anderer tagesstrukturierender  \newline Maßnahmen werden die erforderliche  \newline sächliche Ausstattung, die personelle  \newline Ausstattung und die erforderlichen  \newline betriebsnotwendigen Anlagen}}\textbf{ des} \newline \textbf{\textcolor{diffgreen}{Leistungserbringers übernommen.}} \\
\hline
... & ... & ... \\
\hline
 & \textbf{1.   \newline die Leistungsberechtigten zusätzlich zu  \newline den }\textbf{\textcolor{diffred}{in\S{}}}\textbf{ 83 Absatz 2 des Neunten  \newline Buches genannten Voraussetzungen  \newline zur Teilhabe am Leben in der  \newline Gemeinschaft ständig auf die Nutzung  \newline eines Kraftfahrzeugs angewiesen sind  \newline und } & \textbf{1.   \newline die Leistungsberechtigten zusätzlich zu  \newline den }\textbf{\textcolor{diffgreen}{in \S{}}}\textbf{ 83 Absatz 2 des Neunten  \newline Buches genannten Voraussetzungen  \newline zur Teilhabe am Leben in der  \newline Gemeinschaft ständig auf die Nutzung  \newline eines Kraftfahrzeugs angewiesen sind  \newline und } \\
\hline
... & ... & ... \\
\hline
 & \textbf{6.   \newline zur Erreichbarkeit einer  \newline Ansprechperson unabhängig von einer  \newline konkreten Inanspruchnahme (\S{} 35f  \newline Absatz 2 Nummer 2 in Verbindung mit \S{}  \newline 78 Absatz 6 des Neunten Buches)  \newline können an mehrere  \newline Leistungsberechtigte gemeinsam  \newline erbracht werden, soweit dies nach \S{} 5  \newline Absatz 3 für die Leistungsberechtigten  \newline zumutbar ist und mit  \newline Leistungserbringern entsprechende  \newline Vereinbarungen bestehen. Maßgeblich  \newline sind die Ermittlungen und  \newline Feststellungen nach \S{}\S{} 36 bis }\textbf{\textcolor{diffred}{38d}} & \textbf{6.   \newline zur Erreichbarkeit einer  \newline Ansprechperson unabhängig von einer  \newline konkreten Inanspruchnahme (\S{} 35f  \newline Absatz 2 Nummer 2 in Verbindung mit \S{}  \newline 78 Absatz 6 des Neunten Buches)   \newline können an mehrere  \newline Leistungsberechtigte gemeinsam  \newline erbracht werden, soweit dies nach \S{} 5  \newline Absatz 3 für die Leistungsberechtigten  \newline zumutbar ist und mit  \newline Leistungserbringern entsprechende  \newline Vereinbarungen bestehen. Maßgeblich  \newline sind die Ermittlungen und  \newline Feststellungen nach \S{}\S{} 36 bis }\textbf{\textcolor{diffgreen}{38d. \S{} 35f  \newline Absatz 4 bleibt unberührt.}} \\
\hline
 & \textbf{(3) }\textbf{\textcolor{diffred}{Die}}\textbf{ Leistungen nach Absatz 2 sind} \newline \textbf{\textcolor{diffred}{auf Wunsch der Leistungsberechtigten}}\textbf{  \newline gemeinsam zu erbringen, soweit die  \newline Teilhabeziele erreicht werden können. } & \textbf{(3) }\textbf{\textcolor{diffgreen}{Auf Wunsch der  \newline Leistungsberechtigten sind die}}\textbf{  \newline Leistungen nach Absatz 2 sind }\textbf{\textcolor{diffgreen}{an  \newline mehrere Leistungsberechtigte}}\textbf{  \newline gemeinsam zu erbringen, soweit die  \newline Teilhabeziele erreicht werden können.  } \\
\hline
\textbf{Unterabschnitt 3 } & \textbf{Unterabschnitt 4   \newline Gemeinsame Vorschriften für  \newline Leistungen zur Entwicklung, zur  \newline Erziehung und zur Teilhabe } & \textbf{Unterabschnitt 4   \newline Gemeinsame Vorschriften für  \newline Leistungen zur Entwicklung, zur  \newline Erziehung und zur Teilhabe } \\
\hline
... & ... & ... \\
\hline
\textbf{Mitwirkung, Hilfeplan} & \textbf{Grundsätze der Hilfe- und  \newline Leistungsplanung } & \textbf{Grundsätze der Hilfe- und  \newline Leistungsplanung } \\
\hline
... & ... & ... \\
\hline
 & \textbf{3.   \newline Feststellungen über den individuellen  \newline Bedarf des Kindes oder Jugendlichen  \newline unter Einbeziehung seines engeren  \newline sozialen Umfelds}\textbf{\textcolor{diffred}{;}} & \textbf{3.   \newline Feststellungen über den individuellen  \newline Bedarf des Kindes oder Jugendlichen  \newline unter Einbeziehung seines engeren  \newline sozialen Umfelds}\textbf{\textcolor{diffgreen}{,}} \\
\hline
 & \textbf{4.} \newline \textbf{\textcolor{diffred}{Auswahl}}\textbf{ der }\textbf{\textcolor{diffred}{zu gewährenden}}\textbf{ Art der  \newline Hilfe oder Leistung }\textbf{\textcolor{diffred}{sowie}}\textbf{ deren} \newline \textbf{\textcolor{diffred}{notwendige Ausgestaltung;}} & \textbf{4.} \newline \textbf{\textcolor{diffgreen}{Durchführung einer Hilfe- und  \newline Leistungsplankonferenz nach \S{} 36b zur  \newline Abstimmung}}\textbf{ der Art der Hilfe oder  \newline Leistung }\textbf{\textcolor{diffgreen}{und}}\textbf{ deren }\textbf{\textcolor{diffgreen}{notwendiger  \newline Ausgestaltung nach Inhalt, Umfang und  \newline Dauer unter Beteiligung betroffener  \newline Leistungsträger,}} \\
\hline
 & \textbf{5.} \newline \textbf{\textcolor{diffred}{Durchführung einer Hilfe- und  \newline Leistungsplankonferenz nach \S{} 36b zur  \newline Abstimmung}}\textbf{ der Art der Hilfe oder  \newline Leistung }\textbf{\textcolor{diffred}{und}}\textbf{ deren }\textbf{\textcolor{diffred}{notwendiger  \newline Ausgestaltung nach Inhalt, Umfang und  \newline Dauer unter Beteiligung betroffener  \newline Leistungsträger,}} & \textbf{5.} \newline \textbf{\textcolor{diffgreen}{Auswahl}}\textbf{ der }\textbf{\textcolor{diffgreen}{zu gewährenden}}\textbf{ Art der  \newline Hilfe oder Leistung }\textbf{\textcolor{diffgreen}{sowie}}\textbf{ deren} \newline \textbf{\textcolor{diffgreen}{notwendige Ausgestaltung;}} \\
\hline
... & ... & ... \\
\hline
(3) Werden bei der Durchführung der Hilfe  \newline andere Personen, Dienste oder  \newline Einrichtungen tätig, so sind sie oder deren  \newline Mitarbeiterinnen und Mitarbeiter an der  \newline Aufstellung des Hilfeplans und seiner  \newline Überprüfung zu beteiligen. Soweit dies zur  \newline Feststellung des Bedarfs, der zu  \newline gewährenden Art der Hilfe oder der  \newline notwendigen Leistungen nach Inhalt,  \newline Umfang und Dauer erforderlich ist, sollen  \newline öffentliche Stellen, insbesondere andere  \newline Sozialleistungsträger, Rehabilitationsträger  \newline oder die Schule beteiligt werden. Gewährt  \newline der Träger der öffentlichen Jugendhilfe  \newline Leistungen zur Teilhabe, sind die  \newline Vorschriften zum Verfahren bei einer  \newline Mehrheit von Rehabilitationsträgern nach  \newline dem Neunten Buch zu beachten. & \textbf{(3) Das Kind oder der Jugendliche und  \newline der Personensorgeberechtigte sind vor  \newline der Entscheidung über die  \newline Inanspruchnahme einer Hilfe oder  \newline Leistung nach diesem Abschnitt und vor  \newline einer notwendigen Änderung von Art  \newline und Umfang der gewährten Hilfe oder  \newline Leistung zu beraten und auf die  \newline möglichen Folgen einer Hilfe- oder  \newline Leistungsgewährung für die  \newline Entwicklung des Kindes oder des  \newline Jugendlichen sowie für sei}\textbf{\textcolor{diffred}{-}}\textbf{ne familiale  \newline Lebenssituation hinzuweisen.  \newline Beteiligung in allen Verfahrensschritten  \newline sowie Beratung und Aufklärung nach  \newline Satz 1 erfolgen in einer für den  \newline Personensorgeberechtigten und das  \newline Kind oder den Jugendlichen  \newline verständlichen, nachvollziehbaren und  \newline wahrnehmbaren Form zu. An der Hilfe-  \newline und Leistungsplanung wird auf  \newline Verlangen des Leistungsberechtigten  \newline eine Person seines Vertrauens beteiligt. } & \textbf{(3) Das Kind oder der Jugendliche und  \newline der Personensorgeberechtigte sind vor  \newline der Entscheidung über die  \newline Inanspruchnahme einer Hilfe oder  \newline Leistung nach diesem Abschnitt und vor  \newline einer notwendigen Änderung von Art  \newline und Umfang der gewährten Hilfe oder  \newline Leistung zu beraten und auf die  \newline möglichen Folgen einer Hilfe- oder  \newline Leistungsgewährung für die  \newline Entwicklung des Kindes oder des  \newline Jugendlichen sowie für seine familiale  \newline Lebenssituation hinzuweisen.  \newline Beteiligung in allen Verfahrensschritten  \newline sowie Beratung und Aufklärung nach  \newline Satz 1 erfolgen in einer für den  \newline Personensorgeberechtigten und das  \newline Kind oder den Jugendlichen  \newline verständlichen, nachvollziehbaren und  \newline wahrnehmbaren Form zu. An der Hilfe-  \newline und Leistungsplanung wird auf  \newline Verlangen des Leistungsberechtigten  \newline eine Person seines Vertrauens beteiligt. } \\
\hline
(4) Erscheinen Hilfen nach \S{} 35a  \newline erforderlich, so soll bei der Aufstellung und  \newline Änderung des Hilfeplans sowie bei der  \newline Durchführung der Hilfe die Person, die eine  \newline Stellungnahme nach \S{} 35a Absatz 1a  \newline abgegeben hat, beteiligt werden. & \textbf{(4) }\textbf{\textcolor{diffred}{Di}}\textbf{e Prinzipien der} & \textbf{(4) }\textbf{\textcolor{diffgreen}{Folg}}\textbf{e}\textbf{\textcolor{diffgreen}{nde}}\textbf{ Prinzipien }\textbf{\textcolor{diffgreen}{finden bei}}\textbf{ der} \newline \textbf{\textcolor{diffgreen}{Hilfe- und Leistungsplanung Beachtung:}} \\
\hline
... & ... & ... \\
\hline
 & \textbf{7.} \newline \textbf{\textcolor{diffred}{Zielorientierung finden Beachtung.}} & \textbf{7.} \newline \textbf{\textcolor{diffgreen}{Zielorientierung.}} \\
\hline
... & ... & ... \\
\hline
(1) Der Träger der öffentlichen Jugendhilfe  \newline trägt die Kosten der Hilfe grundsätzlich nur  \newline dann, wenn sie auf der Grundlage seiner  \newline Entscheidung nach Maßgabe des  \newline Hilfeplans unter Beachtung des Wunsch-  \newline und Wahlrechts erbracht wird; dies gilt  \newline auch in den Fällen, in denen Eltern durch  \newline das Familiengericht oder Jugendliche und  \newline junge Volljährige durch den Jugendrichter  \newline zur Inanspruchnahme von Hilfen  \newline verpflichtet werden. Die Vorschriften über  \newline die Heranziehung zu den Kosten der Hilfe  \newline bleiben unberührt. & \textbf{(1) Als Grundlage für den  \newline Verwaltungsakt über die ausgewählte  \newline Hilfe oder Leistung und für deren  \newline Ausgestaltung stellt der Träger der  \newline öffentlichen Jugendhilfe zusammen mit  \newline dem Personensorgeberechtigten und  \newline dem Kind oder Jugendlichen einen  \newline Hilfe- und Leistungsplan }\textbf{\textcolor{diffred}{auf, der  \newline Feststellungen über den Bedarf, die  \newline verfügbaren und aktivierbaren  \newline Selbsthilferessourcen des  \newline Leistungsberechtigten, die zu  \newline gewährende Art der Hilfe oder Leistung  \newline sowie deren notwendige Ausgestaltung  \newline nach Inhalt, Umfang und Dauer enthält.  \newline Die Berücksichtigung des Wunsch- und  \newline Wahlrechts nach \S{} 5 sowie notwendige  \newline Beteiligungen nach Absatz 4 und 5  \newline werden im}}\textbf{ Hilfe- und Leistungsplan} \newline \textbf{\textcolor{diffred}{dokumentiert.}}\textbf{ Ist eine Hilfe- und  \newline Leistungsplankonferenz nach \S{} 36b  \newline durchgeführt worden, sind deren  \newline Ergebnisse der Erstellung des Hilfe- und  \newline Leistungsplans zugrunde zu legen. } & \textbf{(1) Als Grundlage für den  \newline Verwaltungsakt über die ausgewählte  \newline Hilfe oder Leistung und für deren  \newline Ausgestaltung stellt der Träger der  \newline öffentlichen Jugendhilfe zusammen mit  \newline dem Personensorgeberechtigten und  \newline dem Kind oder Jugendlichen einen  \newline Hilfe- und Leistungsplan }\textbf{\textcolor{diffgreen}{auf. Der}}\textbf{ Hilfe-  \newline und Leistungsplan }\textbf{\textcolor{diffgreen}{dient der Steuerung,  \newline Wirkungskontrolle und Dokumentation  \newline des Hilfe- und Leistungsprozesses. Er  \newline soll regelmäßig, dem Bedarf im  \newline Einzelfall entsprechend, überprüft und  \newline fortgeschrieben werden.}}\textbf{ Ist eine Hilfe-  \newline und Leistungsplankonferenz nach \S{} 36b  \newline durchgeführt worden, sind deren  \newline Ergebnisse der Erstellung des Hilfe- und  \newline Leistungsplans zugrunde zu legen. } \\
\hline
(2) Abweichend von Absatz 1 soll der  \newline Träger der öffentlichen Jugendhilfe die  \newline niedrigschwellige unmittelbare  \newline Inanspruchnahme von ambulanten Hilfen,  \newline insbesondere der Erziehungsberatung  \newline nach \S{} 28, zulassen. Dazu soll der Träger  \newline der öffentlichen Jugendhilfe mit den  \newline Leistungserbringern Vereinbarungen  \newline schließen, in denen die Voraussetzungen  \newline und die Ausgestaltung der  \newline Leistungserbringung sowie die Übernahme  \newline der Kosten geregelt werden. Dabei finden  \newline der nach \S{} 80 Absatz 1 Nummer 2  \newline ermittelte Bedarf, die Planungen zur  \newline Sicherstellung des bedarfsgerechten  \newline Zusammenwirkens der Angebote von  \newline Jugendhilfeleistungen in den Lebens- und  \newline Wohnbereichen von jungen Menschen und  \newline Familien nach \S{} 80 Absatz 2 Nummer 3  \newline sowie die geplanten Maßnahmen zur  \newline Qualitätsgewährleistung der  \newline Leistungserbringung nach \S{} 80 Absatz 3  \newline Beachtung. & \textbf{(2) Der Hilfe- und Leistungsplan }\textbf{\textcolor{diffred}{dient}}\textbf{  \newline der }\textbf{\textcolor{diffred}{Steuerung, Wirkungskontrolle}}\textbf{ und} \newline \textbf{\textcolor{diffred}{Dokumentation}}\textbf{ des Hilfe- und} \newline \textbf{\textcolor{diffred}{Leistungsprozesses. Er soll regelmäßig,  \newline spätestens nach zwei Jahren, überprüft  \newline und fortgeschrieben werden.}}\textbf{  \newline Erreichbare und überprüfbare Ziele der  \newline Hilfe oder Leistung und deren  \newline Fortschreibung sowie Maßstäbe und  \newline Kriterien der Wirkungskontrolle  \newline einschließlich des  \newline Überprüfungszeitpunkts werden im  \newline Hilfe- und Leistungsplan festgehalten. } & \textbf{(2) Der Hilfe- und Leistungsplan }\textbf{\textcolor{diffgreen}{enthält  \newline Feststellungen über den Bedarf, die  \newline verfügbaren und aktivierbaren  \newline Selbsthilferessourcen des  \newline Leistungsberechtigten, die zu  \newline gewährende Art}}\textbf{ der }\textbf{\textcolor{diffgreen}{Hilfe oder Leistung  \newline sowie deren notwendige Ausgestaltung  \newline nach Inhalt, Umfang}}\textbf{ und }\textbf{\textcolor{diffgreen}{Dauer. Die  \newline Berücksichtigung}}\textbf{ des }\textbf{\textcolor{diffgreen}{Wunsch- und  \newline Wahlrechts nach \S{} 5 sowie notwendige  \newline Beteiligungen nach Absatz 4 und 5  \newline werden im}}\textbf{ Hilfe- und }\textbf{\textcolor{diffgreen}{Leistungsplan  \newline dokumentiert.}}\textbf{ Erreichbare und  \newline überprüfbare Ziele der Hilfe oder  \newline Leistung und deren Fortschreibung  \newline sowie Maßstäbe und Kriterien der  \newline Wirkungskontrolle einschließlich des  \newline Überprüfungszeitpunkts werden im  \newline Hilfe- und Leistungsplan festgehalten.} \\
\hline
(3) Werden Hilfen abweichend von den  \newline Absätzen 1 und 2 vom  \newline Leistungsberechtigten selbst beschafft, so  \newline ist der Träger der öffentlichen Jugendhilfe  \newline zur Übernahme der erforderlichen  \newline Aufwendungen nur verpflichtet, wenn & \textbf{(3) Hat das Kind oder der Jugendliche  \newline ein oder mehrere Geschwister, so soll  \newline der Geschwisterbeziehung bei der  \newline Aufstellung und Überprüfung des Hilfe-  \newline und Leistungs}\textbf{\textcolor{diffred}{-}}\textbf{plans sowie bei der  \newline Durchführung der Hilfe oder Leistung  \newline Rechnung getragen werden. } & \textbf{(3) Hat das Kind oder der Jugendliche  \newline ein oder mehrere Geschwister, so soll  \newline der Geschwisterbeziehung bei der  \newline Aufstellung und Überprüfung des Hilfe-  \newline und Leistungsplans sowie bei der  \newline Durchführung der Hilfe oder Leistung  \newline Rechnung getragen werden. } \\
\hline
... & ... & ... \\
\hline
 & \textbf{(4) Werden bei der Durchführung der  \newline Hilfe oder Leistung andere Personen,  \newline Dienste oder Einrichtungen tätig, so  \newline sind sie oder deren }\textbf{\textcolor{diffred}{Mitarbeiterinnen  \newline und Mitarbeiter}}\textbf{ an der Aufstellung des  \newline Hilfe- und Leistungsplans und seiner  \newline Überprüfung zu beteiligen. Soweit dies  \newline zur Feststellung des Bedarfs, der zu  \newline gewährenden Art der Hilfe oder Leistung  \newline oder von deren notwendigen Leistungen  \newline nach Inhalt, Umfang und Dauer  \newline erforderlich ist, sollen öffentliche  \newline Stellen, insbesondere andere  \newline Sozialleistungsträger,  \newline Rehabilitationsträger oder die Schule  \newline beteiligt werden. } & \textbf{(4) Werden bei der Durchführung der  \newline Hilfe oder Leistung andere Personen,  \newline Dienste oder Einrichtungen tätig, so  \newline sind sie oder deren }\textbf{\textcolor{diffgreen}{Mitarbeitende}}\textbf{ an der  \newline Aufstellung des Hilfe- und  \newline Leistungsplans und seiner Überprüfung  \newline zu beteiligen. Soweit dies zur  \newline Feststellung des Bedarfs, der zu  \newline gewährenden Art der Hilfe oder Leistung  \newline oder von deren notwendigen Leistungen  \newline nach Inhalt, Umfang und Dauer  \newline erforderlich ist, sollen öffentliche  \newline Stellen, insbesondere andere  \newline Sozialleistungsträger,  \newline Rehabilitationsträger oder die Schule  \newline beteiligt werden. } \\
\hline
... & ... & ... \\
\hline
 & \textbf{(6) Der Hilfe- und Leistungsplan bedarf  \newline der }\textbf{\textcolor{diffred}{Schrif}}\textbf{tform. Der Träger der  \newline öffentlichen Jugendhilfe stellt ihn dem  \newline Leistungsberechtigten zur Verfügung. } & \textbf{(6) Der Hilfe- und Leistungsplan bedarf  \newline der }\textbf{\textcolor{diffgreen}{Tex}}\textbf{tform. Der Träger der  \newline öffentlichen Jugendhilfe stellt ihn dem  \newline Leistungsberechtigten zur Verfügung. } \\
\hline
... & ... & ... \\
\hline
(1) Zur Sicherstellung von Kontinuität und  \newline Bedarfsgerechtigkeit der  \newline Leistungsgewährung sind von den  \newline zuständigen öffentlichen Stellen,  \newline insbesondere von Sozialleistungsträgern  \newline oder Rehabilitationsträgern rechtzeitig im  \newline Rahmen des Hilfeplans Vereinbarungen  \newline zur Durchführung des  \newline Zuständigkeitsübergangs zu treffen. Im  \newline Rahmen der Beratungen zum  \newline Zuständigkeitsübergang prüfen der Träger  \newline der öffentlichen Jugendhilfe und die andere  \newline öffentliche Stelle, insbesondere der andere  \newline Sozialleistungsträger oder  \newline Rehabilitationsträger gemeinsam, welche  \newline Leistung nach dem  \newline Zuständigkeitsübergang dem Bedarf des  \newline jungen Menschen entspricht. & \textbf{(1) Mit Zustimmung des  \newline Leistungsberechtigten }\textbf{\textcolor{diffred}{kann}}\textbf{ der Träger  \newline der öffentlichen Jugendhilfe zur  \newline Aufstellung oder Überprüfung des Hilfe-  \newline und Leistungsplans nach \S{} 36a eine  \newline Hilfe- und Leistungsplankonferenz unter  \newline Berücksichtigung der Willensäußerung  \newline und der Interessen des Kindes oder  \newline Jugendlichen durchführen. }\textbf{\textcolor{diffred}{Der  \newline Leistungsberechtigte oder}}\textbf{ die nach \S{}  \newline 36a Absatz 4 Satz 1 an der Aufstellung  \newline des Hilfe- und Leistungsplans  \newline Beteiligten }\textbf{\textcolor{diffred}{können}}\textbf{ dem Träger der  \newline öffentlichen Jugendhilfe die  \newline Durchführung einer Hilfe- und  \newline Leistungsplankonferenz }\textbf{\textcolor{diffred}{vorschlagen.  \newline Den Vorschlag auf}}\textbf{ Durchführung einer  \newline Hilfe- und Leistungsplankonferenz }\textbf{\textcolor{diffred}{kann  \newline der Träger der öffentlichen Jugendhilfe  \newline ablehnen, wenn der maßgebliche  \newline Sachverhalt schriftlich ermittelt werden  \newline kann, der Aufwand zur Durchführung  \newline nicht in einem angemessenen  \newline Verhältnis zum Umfang der beantragten  \newline Hilfe oder Leistung steht oder dadurch  \newline der Hilfezweck in Frage gestellt wird.}} & \textbf{(1) Mit Zustimmung }\textbf{\textcolor{diffgreen}{oder auf Vorschlag}}\textbf{  \newline des Leistungsberechtigten }\textbf{\textcolor{diffgreen}{soll}}\textbf{ der  \newline Träger der öffentlichen Jugendhilfe zur  \newline Aufstellung oder Überprüfung des Hilfe-  \newline und Leistungsplans nach \S{} 36a eine  \newline Hilfe- und Leistungsplankonferenz unter  \newline Berücksichtigung der Willensäußerung  \newline und der Interessen des Kindes oder  \newline Jugendlichen durchführen. }\textbf{\textcolor{diffgreen}{Dies gilt  \newline auch, wenn}}\textbf{ die nach \S{} 36a Absatz 4 Satz  \newline 1 an der Aufstellung des Hilfe- und  \newline Leistungsplans Beteiligten dem Träger  \newline der öffentlichen Jugendhilfe die  \newline Durchführung einer Hilfe- und  \newline Leistungsplankonferenz }\textbf{\textcolor{diffgreen}{vorschlagen  \newline und der Leistungsberechtigte der}}\textbf{  \newline Durchführung einer Hilfe- und  \newline Leistungsplankonferenz }\textbf{\textcolor{diffgreen}{zustimmt.}} \\
\hline
(2) Abweichend von Absatz 1 werden bei  \newline einem Zuständigkeitsübergang vom Träger  \newline der öffentlichen Jugendhilfe auf einen  \newline Träger der Eingliederungshilfe rechtzeitig  \newline im Rahmen eines Teilhabeplanverfahrens  \newline nach \S{} 19 des Neunten Buches die  \newline Voraussetzungen für die Sicherstellung  \newline einer nahtlosen und bedarfsgerechten  \newline Leistungsgewährung nach dem  \newline Zuständigkeitsübergang geklärt. Die  \newline Teilhabeplanung ist frühzeitig, in der Regel  \newline ein Jahr vor dem voraussichtlichen  \newline Zuständigkeitswechsel, vom Träger der  \newline Jugendhilfe einzuleiten. Mit Zustimmung  \newline des Leistungsberechtigten oder seines  \newline Personensorgeberechtigten ist eine  \newline Teilhabeplankonferenz nach \S{} 20 des  \newline Neunten Buches durchzuführen. Stellt der  \newline beteiligte Träger der Eingliederungshilfe  \newline fest, dass seine Zuständigkeit sowie die  \newline Leistungsberechtigung absehbar gegeben  \newline sind, soll er entsprechend \S{} 19 Absatz 5  \newline des Neunten Buches die Teilhabeplanung  \newline vom Träger der öffentlichen Jugendhilfe  \newline übernehmen. Dies beinhaltet gemäß \S{} 21  \newline des Neunten Buches auch die  \newline Durchführung des Verfahrens zur  \newline Gesamtplanung nach den \S{}\S{} 117 bis 122  \newline des Neunten Buches. & \textbf{(2) In einer Hilfe- und  \newline Leistungsplankonferenz beraten der  \newline Träger der öffentlichen Jugendhilfe, das  \newline Kind oder der Jugendliche und der  \newline Personensorgeberechtigte gemeinsam  \newline auf der Grundlage des festgestellten  \newline Bedarfs insbesondere über die Art der  \newline Hilfe oder Leistung und deren  \newline notwendiger Ausgestaltung nach Inhalt,  \newline Umfang und Dauer unter Einbeziehung  \newline der nach \S{} 36a Absatz 4 und 5 an der  \newline Aufstellung des Hilfe- und  \newline Leistungsplans Beteiligten. Auf  \newline Verlangen des }\textbf{\textcolor{diffred}{Leistungsberechtigten}}\textbf{  \newline wird eine Person seines Vertrauens an  \newline der Hilfe- und Leistungsplankonferenz  \newline beteiligt. } & \textbf{(2) In einer Hilfe- und  \newline Leistungsplankonferenz beraten der  \newline Träger der öffentlichen Jugendhilfe, das  \newline Kind oder der Jugendliche und der  \newline Personensorgeberechtigte gemeinsam  \newline auf der Grundlage des festgestellten  \newline Bedarfs insbesondere über die Art der  \newline Hilfe oder Leistung und deren  \newline notwendiger Ausgestaltung nach Inhalt,  \newline Umfang und Dauer unter Einbeziehung } \textbf{der nach \S{} 36a Absatz 4 und 5 an der  \newline Aufstellung des Hilfe- und  \newline Leistungsplans Beteiligten. Auf  \newline Verlangen des} \newline \textbf{\textcolor{diffgreen}{Personensorgeberechtigten oder des  \newline Kindes oder Jugendlichen}}\textbf{ wird eine  \newline Person seines Vertrauens an der Hilfe-  \newline und Leistungsplankonferenz beteiligt.} \\
\hline
... & ... & ... \\
\hline
 & \textbf{(2) Abweichend von Absatz 1 werden  \newline bei einem Zuständigkeitsübergang vom  \newline Träger der öffentlichen Jugendhilfe auf  \newline einen Träger der Eingliederungshilfe  \newline rechtzeitig im Rahmen eines  \newline Teilhabeplanverfahrens nach \S{} 19 des  \newline Neunten Buches die Voraussetzungen  \newline für die Sicherstellung einer nahtlosen  \newline und bedarfsgerechten  \newline Leistungsgewährung nach dem  \newline Zuständigkeitsübergang geklärt. Die  \newline Teilhabeplanung ist frühzeitig, in der  \newline Regel ein Jahr vor dem  \newline voraussichtlichen  \newline Zuständigkeitswechsel, vom Träger der  \newline Jugendhilfe einzuleiten. Mit Zustimmung  \newline des Leistungsberechtigten oder seines  \newline Personensorgeberechtigten ist eine  \newline Teilhabeplankonferenz nach \S{} 20 des  \newline Neunten Buches durchzuführen. Stellt  \newline der beteiligte Träger der  \newline Eingliederungshilfe fest, dass seine  \newline Zuständigkeit sowie die  \newline Leistungsberechtigung absehbar  \newline gegeben sind, soll er entsprechend \S{} 19  \newline Absatz 5 des Neunten Buches die  \newline Teilhabeplanung vom Träger der  \newline öffentlichen Jugendhilfe übernehmen.  \newline Dies beinhaltet gemäß \S{} 21 des Neunten  \newline Buches auch die Durchführung des  \newline Verfahrens zur Gesamtplanung nach  \newline den \S{}\S{} 117 bis 122 des Neunten Buches. } & \textbf{(2) Abweichend von Absatz 1 werden  \newline bei einem Zuständigkeitsübergang vom  \newline Träger der öffentlichen Jugendhilfe auf  \newline einen Träger der Eingliederungshilfe  \newline rechtzeitig im Rahmen eines  \newline Teilhabeplanverfahrens nach \S{} 19 des  \newline Neunten Buches die Voraussetzungen  \newline für die Sicherstellung einer nahtlosen  \newline und bedarfsgerechten  \newline Leistungsgewährung nach dem  \newline Zuständigkeitsübergang geklärt. Die  \newline Teilhabeplanung ist frühzeitig, in der  \newline Regel ein Jahr vor dem  \newline voraussichtlichen  \newline Zuständigkeitswechsel, vom Träger der  \newline Jugendhilfe einzuleiten. Mit Zustimmung  \newline des Leistungsberechtigten oder seines  \newline Personensorgeberechtigten ist eine  \newline Teilhabeplankonferenz nach \S{} 20 des  \newline Neunten Buches durchzuführen. Stellt  \newline der beteiligte Träger der  \newline Eingliederungshilfe fest, dass seine } \textbf{Zuständigkeit sowie die  \newline Leistungsberechtigung absehbar  \newline gegeben sind, soll er entsprechend \S{} 19  \newline Absatz 5 des Neunten Buches die  \newline Teilhabeplanung vom Träger der  \newline öffentlichen Jugendhilfe übernehmen.  \newline Dies beinhaltet gemäß \S{} 21 des Neunten  \newline Buches auch die Durchführung des  \newline Verfahrens zur Gesamtplanung nach  \newline den \S{}\S{} 117 bis 122 des Neunten Buches.} \newline \textbf{\textcolor{diffgreen}{\S{} 41a findet keine Anwendung.}} \\
\hline
... & ... & ... \\
\hline
(3) Sofern der Inhaber der elterlichen  \newline Sorge durch eine Erklärung nach \S{} 1688  \newline Absatz 3 Satz 1 des Bürgerlichen  \newline Gesetzbuchs die Entscheidungsbefugnisse  \newline der Pflegeperson so weit einschränkt, dass  \newline die Einschränkung eine dem Wohl des  \newline Kindes oder des Jugendlichen förderliche  \newline Entwicklung nicht mehr ermöglicht, sollen  \newline die Beteiligten das Jugendamt einschalten.  \newline Auch bei sonstigen  \newline Meinungsverschiedenheiten zwischen  \newline ihnen sollen die Beteiligten das Jugendamt  \newline einschalten. & \textbf{(3) Bei der Auswahl der Einrichtung  \newline oder der Pflegeperson sind der  \newline Personensorgeberechtigte und das Kind  \newline oder der Jugendliche oder bei Hilfen  \newline nach \S{} 41 der junge Volljährige zu  \newline beteiligen. Der Wahl und den Wünschen  \newline des Leistungsberechtigten ist zu  \newline entsprechen, sofern sie nicht mit  \newline unverhältnismäßigen Mehrkosten  \newline verbunden sind. Wünschen die in Satz 1  \newline genannten Personen die Erbringung  \newline einer in \S{} 78a genannten Leistung in  \newline einer Einrichtung, mit deren Träger  \newline keine Vereinbarungen nach \S{} 78b  \newline bestehen, so soll der Wahl nur  \newline entsprochen werden, wenn die  \newline Erbringung der Leistung in dieser  \newline Einrichtung nach Maßgabe des Hilfe-  \newline oder Leistungsplans geboten ist. Bei  \newline der }\textbf{\textcolor{diffred}{Auswahl einer Pflegeperson, die  \newline ihren gewöhnlichen Aufenthalt  \newline außerhalb des Bereichs des örtlich  \newline zuständigen Trägers hat, soll der  \newline örtliche Träger der öffentlichen  \newline Jugendhilfe beteiligt werden, in dessen  \newline Bereich die Pflegeperson ihren  \newline gewöhnlichen Aufenthalt hat. Bei der}}\textbf{  \newline Entscheidung nach Satz 2 und 3 hat  \newline zunächst die Prüfung nach \S{} 5 Absatz 3  \newline zu erfolgen. } & \textbf{(3) Bei der Auswahl der Einrichtung  \newline oder der Pflegeperson sind der  \newline Personensorgeberechtigte und das Kind  \newline oder der Jugendliche oder bei Hilfen  \newline nach \S{} 41 der junge Volljährige zu  \newline beteiligen. Der Wahl und den Wünschen  \newline des Leistungsberechtigten ist zu  \newline entsprechen, sofern sie nicht mit  \newline unverhältnismäßigen Mehrkosten  \newline verbunden sind. Wünschen die in Satz 1  \newline genannten Personen die Erbringung  \newline einer in \S{} 78a genannten Leistung in  \newline einer Einrichtung, mit deren Träger  \newline keine Vereinbarungen nach \S{} 78b  \newline bestehen, so soll der Wahl nur  \newline entsprochen werden, wenn die  \newline Erbringung der Leistung in dieser  \newline Einrichtung nach Maßgabe des Hilfe-  \newline oder Leistungsplans geboten ist. Bei  \newline der Entscheidung nach Satz 2 und 3 hat  \newline zunächst die Prüfung nach \S{} 5 Absatz 3  \newline zu erfolgen. }\textbf{\textcolor{diffgreen}{Bei der Auswahl einer  \newline Pflegeperson, die ihren gewöhnlichen  \newline Aufenthalt außerhalb des Bereichs des  \newline örtlich zuständigen Trägers hat, oder  \newline bei der Auswahl einer familienähnlichen  \newline Betreuungsform, die außerhalb des  \newline Bereichs des örtlich zuständigen  \newline Trägers gelegen ist, muss der örtliche  \newline Träger der öffentlichen Jugendhilfe  \newline hierüber schriftlich informiert werden  \newline und Gelegenheit zur Stellungnahme  \newline erhalten, in dessen Bereich die }} \textbf{\textcolor{diffgreen}{Pflegeperson ihren gewöhnlichen  \newline Aufenthalt hat oder die familienähnliche  \newline Betreuungsform gelegen ist. Über die  \newline Art der Zustellung der schriftlichen  \newline Information nach Satz 5 entscheidet der  \newline örtliche Träger der öffentlichen  \newline Jugendhilfe nach pflichtgemäßem  \newline Ermessen; es gelten die jeweiligen  \newline landesrechtlichen Vorschriften über das  \newline Zustellungsverfahren.}} \\
\hline
... & ... & ... \\
\hline
\textbf{Beratung und Unterstützung der  \newline Pflegeperson } & entfällt & entfällt \\
\hline
... & ... & ... \\
\hline
\textbf{Sicherung der Rechte von Kindern und  \newline Jugendlichen in Familienpflege } & entfällt & entfällt \\
\hline
... & ... & ... \\
\hline
\textbf{Ergänzende Bestimmungen zur  \newline Hilfeplanung bei Hilfen außerhalb der  \newline eigenen Familie } & entfällt & entfällt \\
\hline
... & ... & ... \\
\hline
 & \textbf{(2) Liegen keine als  \newline Entscheidungsgrundlage}\textbf{\textcolor{diffred}{n}}\textbf{  \newline ausreichenden Gutachten, ärztliche  \newline Stellungnahmen oder vergleichbare  \newline Bescheinigungen vor, prüft der Träger  \newline der öffentlichen Jugendhilfe, ob für die  \newline Feststellung des Rehabilitationsbedarfs  \newline eine }\textbf{\textcolor{diffred}{kürzere}}\textbf{ ärztliche Stellungnahme  \newline oder vergleichbare Bescheinigung  \newline insbesondere hinsichtlich des  \newline Vorliegens einer körperlichen,  \newline seelischen, geistigen oder  \newline Sinnesbeeinträchtigung nach \S{} 7 Absatz  \newline 2 Satz 2 und 3 erforderlich und  \newline ausreichend ist. Ist dies der Fall, holt er} \newline \textbf{\textcolor{diffred}{diese ein und legt sie seinen  \newline Entscheidungen zugrunde.}}\textbf{ \S{} 17 Absatz  \newline 2 Satz 1 des Neunten Buches gilt  \newline entsprechend. Diese Stellungnahme }\textbf{\textcolor{diffred}{soll}}\textbf{  \newline der Träger der öffentlichen Jugendhilfe  \newline als leistender Rehabilitationsträger} \newline \textbf{\textcolor{diffred}{seinen Entscheidungen}}\textbf{ zugrunde zu  \newline legen. Die gewährten Leistungen der  \newline Eingliederungshilfe sollen }\textbf{\textcolor{diffred}{nicht}}\textbf{ von der} \newline \textbf{\textcolor{diffred}{Person}}\textbf{ oder dem Dienst oder der  \newline Einrichtung erbracht werden, der die  \newline Person }\textbf{\textcolor{diffred}{angehört, die eine  \newline Stellungnahme nach Satz 1 abgibt.}} & \textbf{(2) Liegen keine als  \newline Entscheidungsgrundlage  \newline ausreichenden Gutachten, ärztliche}\textbf{\textcolor{diffgreen}{n}}\textbf{  \newline Stellungnahmen oder vergleichbare  \newline Bescheinigungen vor, prüft der Träger  \newline der öffentlichen Jugendhilfe, ob für die  \newline Feststellung des Rehabilitationsbedarfs  \newline eine ärztliche Stellungnahme oder  \newline vergleichbare Bescheinigung}\textbf{\textcolor{diffgreen}{,}}\textbf{  \newline insbesondere hinsichtlich des  \newline Vorliegens einer körperlichen,  \newline seelischen, geistigen oder  \newline Sinnesbeeinträchtigung nach \S{} 7 Absatz  \newline 2 Satz 2 und 3}\textbf{\textcolor{diffgreen}{,}}\textbf{ erforderlich und  \newline ausreichend ist. Ist dies der Fall, holt er} \newline \textbf{\textcolor{diffgreen}{eine ärztliche Stellungnahme oder eine  \newline vergleichbare Bescheinigung ein.}}\textbf{ \S{} 17  \newline Absatz 2 Satz 1 }\textbf{\textcolor{diffgreen}{zweiter Teilsatz}}\textbf{ des  \newline Neunten Buches gilt entsprechend.  \newline Diese }\textbf{\textcolor{diffgreen}{ärztliche}}\textbf{ Stellungnahme }\textbf{\textcolor{diffgreen}{oder  \newline vergleichbare Bescheinigung hat}}\textbf{ der  \newline Träger der öffentlichen Jugendhilfe als  \newline leistender Rehabilitationsträger }\textbf{\textcolor{diffgreen}{seiner  \newline Entscheidung}}\textbf{ zugrunde zu legen. Die  \newline gewährten Leistungen der  \newline Eingliederungshilfe sollen }\textbf{\textcolor{diffgreen}{weder}}\textbf{ von  \newline der }\textbf{\textcolor{diffgreen}{Person, die die Stellungnahme}}\textbf{ oder } \textbf{\textcolor{diffgreen}{die Bescheinigung nach Satz 2  \newline abgegeben oder ausstellt hat, noch von}}\textbf{  \newline dem Dienst oder der Einrichtung  \newline erbracht werden, der die}\textbf{\textcolor{diffgreen}{se}}\textbf{ Person} \newline \textbf{\textcolor{diffgreen}{angehört.}} \\
\hline
 & \textbf{(3) }\textbf{\textcolor{diffred}{Bei}}\textbf{ der }\textbf{\textcolor{diffred}{Entscheidung über}}\textbf{ die} \newline \textbf{\textcolor{diffred}{Erforderlichkeit eines Gutachtens}}\textbf{ nach  \newline \S{} 17 }\textbf{\textcolor{diffred}{Absatz 1}}\textbf{ des Neunten Buches }\textbf{\textcolor{diffred}{und  \newline bei den Prüfungen nach Absatz 1die  \newline Einholung einer Stellungnahme nach  \newline Satz 1 sind das Kind oder der  \newline Jugendliche und der  \newline Personensorgeberechtigte nach  \newline Maßgabe von \S{} 36 Absatz 3 Satz 2 zu  \newline beteiligen.}} & \textbf{(3)} \textbf{\textcolor{diffgreen}{Hält}}\textbf{ der }\textbf{\textcolor{diffgreen}{Träger der öffentlichen  \newline Jugendhilfe für}}\textbf{ die }\textbf{\textcolor{diffgreen}{Feststellung des  \newline Rehabilitationsbedarfs ein Gutachten für  \newline erforderlich, finden die Regelungen zur  \newline Begutachtung}}\textbf{ nach \S{} 17 des Neunten  \newline Buches }\textbf{\textcolor{diffgreen}{Anwendung.}} \\
\hline
 & \textbf{(4) }\textbf{\textcolor{diffred}{Hält}}\textbf{ der }\textbf{\textcolor{diffred}{Träger}}\textbf{ der }\textbf{\textcolor{diffred}{öffentlichen  \newline Jugendhilfe für die Feststellung des  \newline Rehabilitationsbedarfs ein Gutachten für  \newline erforderlich, finden die Regelungen zur  \newline Begutachtung}}\textbf{ nach \S{} }\textbf{\textcolor{diffred}{17 des Neunten  \newline Buches Anwendung.}} & \textbf{(4) }\textbf{\textcolor{diffgreen}{Bei}}\textbf{ der }\textbf{\textcolor{diffgreen}{Entscheidung über die  \newline Erforderlichkeit eines Gutachtens und  \newline bei den Prüfungen nach Absatz 1 oder  \newline nach Absatz 2 Satz 1 sind das Kind oder}}\textbf{  \newline der }\textbf{\textcolor{diffgreen}{Jugendliche und der  \newline Personensorgeberechtigte}}\textbf{ nach} \newline \textbf{\textcolor{diffgreen}{Maßgabe von}}\textbf{ \S{} }\textbf{\textcolor{diffgreen}{36 Absatz 3 Satz 2 zu  \newline beteiligen.}} \\
\hline
... & ... & ... \\
\hline
 & \textbf{(1) Bei Leistungen der  \newline Eingliederungshilfe nach \S{}\S{} 27 Absatz} \newline \textbf{\textcolor{diffred}{3,}}\textbf{ 35a hat der Träger der öffentlichen  \newline Jugendhilfe als Rehabilitationsträger die  \newline Regelungen zur Erkennung und  \newline Ermittlung des Rehabilitationsbedarfs  \newline des Kapitels 3 des Teils 1 des Neunten  \newline Buches anzuwenden. } & \textbf{(1) Bei Leistungen der  \newline Eingliederungshilfe nach \S{}\S{} 27 Absatz }\textbf{\textcolor{diffgreen}{3  \newline und}}\textbf{ 35a hat der Träger der öffentlichen  \newline Jugendhilfe als Rehabilitationsträger die  \newline Regelungen zur Erkennung und  \newline Ermittlung des Rehabilitationsbedarfs  \newline des Kapitels 3 des Teils 1 des Neunten  \newline Buches anzuwenden. } \\
\hline
... & ... & ... \\
\hline
 & \textbf{\textcolor{diffred}{(3) Das Bundesministerium für Familie,  \newline Senioren, Frauen und Jugend kann  \newline durch Rechtsverordnung mit  \newline Zustimmung des Bundesrates das  \newline Nähere über das Instrument zur  \newline Bedarfsermittlung bestimmen. }} &  \\
\hline
... & ... & ... \\
\hline
 & \textbf{(1) Bei Leistungen der  \newline Eingliederungshilfe nach \S{}\S{} 27 Absatz} \newline \textbf{\textcolor{diffred}{3,}}\textbf{ 35a enthält der Hilfe- und  \newline Leistungsplan die Inhalte nach \S{} 19  \newline Absatz 2 Satz 2 des Neunten Buches;  \newline dies gilt auch, wenn weder Leistungen  \newline verschiedener Leistungsgruppen nach \S{}  \newline 5 des Neunten Buches noch mehrerer  \newline Rehabilitationsträger nach \S{} 6 Absatz 1  \newline des Neunten Buches erforderlich sind  \newline und damit die Voraussetzungen nach \S{}  \newline 19 Absatz 1 nicht vorliegen, der  \newline Leistungsberechtigte aber die  \newline Erstellung eines Teilhabeplans wünscht.  \newline Daneben dokumentiert der Hilfe- und  \newline Leistungsplan mindestens } & \textbf{(1) Bei Leistungen der  \newline Eingliederungshilfe nach \S{}\S{} 27 Absatz }\textbf{\textcolor{diffgreen}{3  \newline und}}\textbf{ 35a enthält der Hilfe- und  \newline Leistungsplan die Inhalte nach \S{} 19  \newline Absatz 2 Satz 2 des Neunten Buches;  \newline dies gilt auch, wenn weder Leistungen  \newline verschiedener Leistungsgruppen nach \S{}  \newline 5 des Neunten Buches noch mehrerer  \newline Rehabilitationsträger nach \S{} 6 Absatz 1  \newline des Neunten Buches erforderlich sind  \newline und damit die Voraussetzungen nach \S{}  \newline 19 Absatz 1 nicht vorliegen, der } \textbf{Leistungsberechtigte aber die  \newline Erstellung eines Teilhabeplans wünscht.  \newline Daneben dokumentiert der Hilfe- und  \newline Leistungsplan mindestens } \\
\hline
... & ... & ... \\
\hline
 & \textbf{(2) Die im Hilfe- und Leistungsplan  \newline festgestellte Leistung sowie deren  \newline Ausgestaltung nach \S{} 36a Absatz 1 sind  \newline für die Entscheidung über die zu  \newline bewilligende und }\textbf{\textcolor{diffred}{er- bringenden}}\textbf{  \newline Leistungen nach \S{} 15 Absatz 3 Satz 1  \newline des Neunten Buches oder \S{} 38 Ab}\textbf{\textcolor{diffred}{-}}\textbf{satz 6  \newline bindend, wenn der Träger der  \newline öffentlichen Jugendhilfe  \newline Leistungsverantwortlicher nach \S{} 15  \newline des Neunten Buches ist. Wenn nach den  \newline Vorschriften zur Koordinierung der  \newline Leistungen nach Teil 1 Kapitel 4 des  \newline Neunten Buches ein anderer  \newline Rehabilitationsträger die  \newline Leistungsverantwortung trägt, bildet die  \newline im Rahmen des Hilfe- und  \newline Leistungsplans nach \S{} 36a Absatz 1  \newline festgestellte Leistung und deren  \newline Ausgestaltung die für den Teilhabeplan  \newline nach \S{} 19 des Neunten Buches  \newline erforderlichen Feststellungen des  \newline Trägers der öffentlichen Jugendhilfe  \newline nach \S{} 15 Absatz 2 des Neunten  \newline Buches. } & \textbf{(2) Die im Hilfe- und Leistungsplan  \newline festgestellte Leistung sowie deren  \newline Ausgestaltung nach \S{} 36a Absatz 1 sind  \newline für die Entscheidung über die zu  \newline bewilligende und }\textbf{\textcolor{diffgreen}{erbringenden}}\textbf{  \newline Leistungen nach \S{} 15 Absatz 3 Satz 1  \newline des Neunten Buches oder \S{} 38 Absatz 6  \newline bindend, wenn der Träger der  \newline öffentlichen Jugendhilfe  \newline Leistungsverantwortlicher nach \S{} 15  \newline des Neunten Buches ist. Wenn nach den  \newline Vorschriften zur Koordinierung der  \newline Leistungen nach Teil 1 Kapitel 4 des  \newline Neunten Buches ein anderer  \newline Rehabilitationsträger die  \newline Leistungsverantwortung trägt, bildet die  \newline im Rahmen des Hilfe- und  \newline Leistungsplans nach \S{} 36a Absatz 1  \newline festgestellte Leistung und deren  \newline Ausgestaltung die für den Teilhabeplan } \textbf{nach \S{} 19 des Neunten Buches  \newline erforderlichen Feststellungen des  \newline Trägers der öffentlichen Jugendhilfe  \newline nach \S{} 15 Absatz 2 des Neunten  \newline Buches. } \\
\hline
... & ... & ... \\
\hline
 & \textbf{(4) Bestehen im Einzelfall Anhaltspunkte  \newline für eine Pflegebedürftigkeit nach dem  \newline Elften Buch, wird die zuständige  \newline Pflegekasse mit Zustimmung des  \newline Leistungsberechtigten vom Träger der  \newline öffentlichen Jugendhilfe informiert und  \newline muss an der Aufstellung und  \newline Überprüfung des Hilfe- und  \newline Leistungsplans beratend teilnehmen,  \newline soweit dies zur Feststellung des  \newline Bedarfs, der zu gewährenden Art der  \newline Leistung oder von deren not}\textbf{\textcolor{diffred}{-}}\textbf{wendiger  \newline Ausgestaltung nach Inhalt, Umfang und  \newline Dauer erforderlich ist. Bestehen im  \newline Einzelfall Anhaltspunkte, dass  \newline Leistungen der Hilfe zur Pflege nach  \newline dem Siebten Kapitel des Zwölften  \newline Buches erforderlich sind, so soll der  \newline Träger dieser Leistungen mit Zu- \newline stimmung der Leistungsberechtigten  \newline informiert und an der Aufstellung und  \newline Überprüfung des Hilfe- und  \newline Leistungsplans beteiligt werden, soweit  \newline dies zur Feststellung des Bedarfs, der  \newline zu gewährenden Art der Leistung oder  \newline von deren notwendiger Ausgestaltung  \newline nach Inhalt, Umfang und Dauer  \newline erforderlich ist. } & \textbf{(4) Bestehen im Einzelfall Anhaltspunkte  \newline für eine Pflegebedürftigkeit nach dem  \newline Elften Buch, wird die zuständige  \newline Pflegekasse mit Zustimmung des  \newline Leistungsberechtigten vom Träger der  \newline öffentlichen Jugendhilfe informiert und  \newline muss an der Aufstellung und  \newline Überprüfung des Hilfe- und  \newline Leistungsplans beratend teilnehmen,  \newline soweit dies zur Feststellung des  \newline Bedarfs, der zu gewährenden Art der  \newline Leistung oder von deren notwendiger  \newline Ausgestaltung nach Inhalt, Umfang und  \newline Dauer erforderlich ist. Bestehen im  \newline Einzelfall Anhaltspunkte, dass  \newline Leistungen der Hilfe zur Pflege nach  \newline dem Siebten Kapitel des Zwölften  \newline Buches erforderlich sind, so soll der  \newline Träger dieser Leistungen mit Zu- \newline stimmung der Leistungsberechtigten  \newline informiert und an der Aufstellung und  \newline Überprüfung des Hilfe- und  \newline Leistungsplans beteiligt werden, soweit  \newline dies zur Feststellung des Bedarfs, der  \newline zu gewährenden Art der Leistung oder  \newline von deren notwendiger Ausgestaltung  \newline nach Inhalt, Umfang und Dauer  \newline erforderlich ist. } \\
\hline
 & \textbf{(5) Soweit Leistungen verschiedener  \newline Leistungsgruppen nach \S{} 5 des Neunten  \newline Buches oder mehrere  \newline Rehabilitationsträger nach \S{} 6 Absatz 1  \newline des Neunten Buches erforderlich sind,  \newline hat der Träger der öffentlichen  \newline Jugendhilfe als leistender  \newline Rehabilitationsträger die Regelungen  \newline zum Teilhabeplan nach \S{} 19 des  \newline Neunten Buches und legt i}\textbf{\textcolor{diffred}{h}}\textbf{n seiner  \newline Entscheidung über die Gewährung einer  \newline Leistung der Eingliederungshilfe  \newline zugrunde. Im Übrigen i}\textbf{\textcolor{diffred}{s}}\textbf{t \S{} 19 des  \newline Neunten }\textbf{\textcolor{diffred}{Buches anzuwenden.}} & \textbf{(5) Soweit Leistungen verschiedener  \newline Leistungsgruppen nach \S{} 5 des Neunten  \newline Buches oder mehrere  \newline Rehabilitationsträger nach \S{} 6 Absatz 1  \newline des Neunten Buches erforderlich sind,} \newline \textbf{\textcolor{diffgreen}{oder der Leistungsberechtigte oder der  \newline Personensorgeberechtigte dies  \newline wünscht,}}\textbf{ hat der Träger der öffentlichen  \newline Jugendhilfe als leistender  \newline Rehabilitationsträger die Regelungen } \textbf{zum Teilhabeplan nach \S{} 19 des  \newline Neunten Buches }\textbf{\textcolor{diffgreen}{anzuwenden}}\textbf{ und legt} \newline \textbf{\textcolor{diffgreen}{d}}\textbf{i}\textbf{\textcolor{diffgreen}{ese}}\textbf{n seiner Entscheidung über die  \newline Gewährung einer Leistung der  \newline Eingliederungshilfe zugrunde. Im  \newline Übrigen }\textbf{\textcolor{diffgreen}{g}}\textbf{i}\textbf{\textcolor{diffgreen}{l}}\textbf{t \S{} 19 des Neunten }\textbf{\textcolor{diffgreen}{Buches.}} \\
\hline
... & ... & ... \\
\hline
 & \textbf{\textcolor{diffred}{Ergänz}}\textbf{ende Bestimmungen zur Hilfe-  \newline und Leistungsplankonferenz bei  \newline Leistungen der Eingliederungshilfe für  \newline Kinder und Jugendliche mit  \newline Behinderungen } & \textbf{\textcolor{diffgreen}{B}}\textbf{e}\textbf{\textcolor{diffgreen}{so}}\textbf{nde}\textbf{\textcolor{diffgreen}{re}}\textbf{ Bestimmungen zur Hilfe- und  \newline Leistungsplankonferenz bei Leistungen  \newline der Eingliederungshilfe für Kinder und  \newline Jugendliche mit Behinderungen } \\
\hline
... & ... & ... \\
\hline
 & \textbf{(2) In eine Hilfe- und  \newline Leistungsplankonferenz }\textbf{\textcolor{diffred}{könn}}\textbf{en die  \newline nach \S{} 38c Absatz 3 und 4 an der  \newline Aufstellung und Überprüfung des Hilfe-  \newline und Leistungsplans Beteiligten  \newline einbezogen werden. } & \textbf{(2) In eine Hilfe- und  \newline Leistungsplankonferenz }\textbf{\textcolor{diffgreen}{soll}}\textbf{en die nach  \newline \S{} 38c Absatz 3 und 4 an der Aufstellung  \newline und Überprüfung des Hilfe- und  \newline Leistungsplans Beteiligten einbezogen  \newline werden. } \\
\hline
 & \textbf{(3) Ist der Träger der öffentlichen  \newline Jugendhilfe Leistungsverantwortlicher  \newline nach \S{} 15 des Neunten Buches, soll er  \newline die Hilfe- und Leistungsplankonferenz  \newline mit einer Teilhabeplankonferenz nach \S{}  \newline 20 des Neunten Buches verbinden. Ist  \newline der Träger der öffentlichen Jugendhilfe  \newline beteiligter Rehabilitationsträger nach \S{}  \newline 15 des Neunten Buches soll er dem  \newline Leistungsberechtigten und den anderen  \newline Rehabilitationsträgern anbieten, mit  \newline deren Einvernehmen das Verfahren  \newline anstelle des leistenden  \newline Rehabilitationsträgers durchzuführen;  \newline die Vorschriften über die  \newline Leistungsverantwortung der  \newline Rehabilitationsträger nach den \S{}\S{} 14  \newline und 15 des Neunten Buches bleiben  \newline hiervon unberührt. } & \textbf{(3) Ist der Träger der öffentlichen  \newline Jugendhilfe Leistungsverantwortlicher  \newline nach \S{} 15 des Neunten Buches, soll er  \newline die Hilfe- und Leistungsplankonferenz  \newline mit einer Teilhabeplankonferenz nach \S{}  \newline 20 des Neunten Buches verbinden. Ist  \newline der Träger der öffentlichen Jugendhilfe  \newline beteiligter Rehabilitationsträger nach \S{}  \newline 15 des Neunten Buches}\textbf{\textcolor{diffgreen}{,}}\textbf{ soll er dem  \newline Leistungsberechtigten und den anderen  \newline Rehabilitationsträgern anbieten, mit  \newline deren Einvernehmen das Verfahren  \newline anstelle des leistenden  \newline Rehabilitationsträgers durchzuführen;  \newline die Vorschriften über die  \newline Leistungsverantwortung der  \newline Rehabilitationsträger nach den \S{}\S{} 14  \newline und 15 des Neunten Buches bleiben  \newline hiervon unberührt. } \\
\hline
... & ... & ... \\
\hline
(1) Wird Hilfe nach den \S{}\S{} 32 bis 35 oder  \newline nach \S{} 35a Absatz 2 Nummer 2 bis 4  \newline gewährt, so ist auch der notwendige  \newline Unterhalt des Kindes oder Jugendlichen  \newline außerhalb des Elternhauses  \newline sicherzustellen. Er umfasst die Kosten für  \newline den Sachaufwand sowie für die Pflege und  \newline Erziehung des Kindes oder Jugendlichen. & \textbf{(1) Werden Hilfen oder Leistungen nach  \newline den \S{}\S{} 32 bis 34 und 35a Absatz 4  \newline Nummer 3 }\textbf{\textcolor{diffred}{und 4}}\textbf{ gewährt, haben die  \newline Eltern einen Anspruch auf Beratung und  \newline Unterstützung sowie Förderung der  \newline Beziehung zu ihrem Kind. Durch  \newline Beratung und Unterstützung sollen die  \newline Entwicklungs-, Teilhabe- oder  \newline Erziehungsbedingungen in der  \newline Herkunftsfamilie innerhalb eines im  \newline Hinblick auf die Entwicklung des Kindes  \newline oder Jugendlichen vertretbaren  \newline Zeitraums so weit verbessert werden,  \newline dass sie das Kind oder den  \newline Jugendlichen wieder selbst erziehen  \newline kann. Ist eine nachhaltige Verbesserung  \newline der Entwicklungs-, Teilhabe- oder  \newline Erziehungsbedingungen in der  \newline Herkunftsfamilie innerhalb dieses  \newline Zeitraums nicht erreichbar, so dienen  \newline die Beratung und Unterstützung der  \newline Eltern sowie die Förderung ihrer  \newline Beziehung zum Kind der Erarbeitung  \newline und Sicherung einer anderen, dem Wohl  \newline des Kindes oder Jugendlichen  \newline förderlichen und auf Dauer angelegten  \newline Lebensperspektive. } & \textbf{(1) Werden Hilfen oder Leistungen nach  \newline den \S{}\S{} 32 bis 34 und 35a Absatz 4  \newline Nummer 3 gewährt, haben die Eltern  \newline einen Anspruch auf Beratung und  \newline Unterstützung sowie Förderung der  \newline Beziehung zu ihrem Kind. Durch  \newline Beratung und Unterstützung sollen die  \newline Entwicklungs-, Teilhabe- oder  \newline Erziehungsbedingungen in der  \newline Herkunftsfamilie innerhalb eines im  \newline Hinblick auf die Entwicklung des Kindes  \newline oder Jugendlichen vertretbaren  \newline Zeitraums so weit verbessert werden,  \newline dass sie das Kind oder den  \newline Jugendlichen wieder selbst erziehen  \newline kann. Ist eine nachhaltige Verbesserung  \newline der Entwicklungs-, Teilhabe- oder  \newline Erziehungsbedingungen in der  \newline Herkunftsfamilie innerhalb dieses  \newline Zeitraums nicht erreichbar, so dienen  \newline die Beratung und Unterstützung der  \newline Eltern sowie die Förderung ihrer  \newline Beziehung zum Kind der Erarbeitung  \newline und Sicherung einer anderen, dem Wohl  \newline des Kindes oder Jugendlichen  \newline förderlichen und auf Dauer angelegten  \newline Lebensperspektive. } \\
\hline
... & ... & ... \\
\hline
 & Die Pflegeperson hat vor der Aufnahme  \newline des Kindes oder des Jugendlichen und  \newline während der Dauer des  \newline Pflegeverhältnisses Anspruch auf Beratung  \newline und Unterstützung. Dies gilt auch in den  \newline Fällen, in denen für das Kind oder den  \newline Jugendlichen weder Hilfe zur Erziehung  \newline noch Eingliederungshilfe gewährt w\textcolor{diffred}{i}rd, und  \newline in den Fällen, in denen die Pflegeperson  \newline nicht der Erlaubnis zur Vollzeitpflege nach  \newline \S{} 44 bedarf. Lebt das Kind oder der  \newline Jugendliche bei einer Pflegeperson  \newline außerhalb des Bereichs des zuständigen  \newline Trägers der öffentlichen Jugendhilfe, so  \newline sind ortsnahe Beratung und Unterstützung  \newline sicherzustellen. Der zuständige Träger der  \newline öffentlichen Jugendhilfe hat die  \newline aufgewendeten Kosten einschließlich der  \newline Verwaltungskosten auch in den Fällen zu  \newline erstatten, in denen die Beratung und  \newline Unterstützung im Wege der Amtshilfe  \newline geleistet werden. Zusammenschlüsse von  \newline Pflegepersonen sollen beraten, unterstützt  \newline und gefördert werden. & Die Pflegeperson hat vor der Aufnahme  \newline des Kindes oder des Jugendlichen und  \newline während der Dauer des  \newline Pflegeverhältnisses Anspruch auf Beratung  \newline und Unterstützung. Dies gilt auch in den  \newline Fällen, in denen für das Kind oder den  \newline Jugendlichen weder Hilfe zur Erziehung  \newline noch \textbf{\textcolor{diffgreen}{Leistungen der}} Eingliederungshilfe  \newline gewährt \textbf{w}\textbf{\textcolor{diffgreen}{e}}\textbf{rd}\textbf{\textcolor{diffgreen}{en}}, und in den Fällen, in  \newline denen die Pflegeperson nicht der Erlaubnis  \newline zur Vollzeitpflege nach \S{} 44 bedarf. Lebt  \newline das Kind oder der Jugendliche bei einer  \newline Pflegeperson außerhalb des Bereichs des  \newline zuständigen Trägers der öffentlichen  \newline Jugendhilfe, so sind ortsnahe Beratung und  \newline Unterstützung sicherzustellen. Der  \newline zuständige Träger der öffentlichen  \newline Jugendhilfe hat die aufgewendeten Kosten  \newline einschließlich der Verwaltungskosten auch  \newline in den Fällen zu erstatten, in denen die  \newline Beratung und Unterstützung im Wege der  \newline Amtshilfe geleistet werden.  \newline Zusammenschlüsse von Pflegepersonen  sollen beraten, unterstützt und gefördert  \newline werden. \\
\hline
... & ... & ... \\
\hline
 & \textbf{1.   \newline im Anwendungsbereich der Verordnung  \newline (EU) 2019/1111 des Rates vom 25. Juni  \newline 2019 über die Zuständigkeit, die  \newline Anerkennung und Vollstreckung von  \newline Entscheidungen in Ehesachen und in  \newline Verfahren betreffend die elterliche  \newline Verantwortung und über internationale  \newline Kindesentführungen (ABl. L 178 vom} \newline \textbf{\textcolor{diffred}{2.7.2019, S. 1) die Voraussetzungen des  \newline Artikels 82 oder}} & \textbf{1.   \newline im Anwendungsbereich der Verordnung  \newline (EU) 2019/1111 des Rates vom 25. Juni  \newline 2019 über die Zuständigkeit, die  \newline Anerkennung und Vollstreckung von  \newline Entscheidungen in Ehesachen und in  \newline Verfahren betreffend die elterliche  \newline Verantwortung und über internationale  \newline Kindesentführungen (ABl. L 178 vom } \\
\hline
 &  & \textbf{\textcolor{diffgreen}{2.7.2019, S. 1) die Voraussetzungen des  \newline Artikels 82 oder}} \\
\hline
... & ... & ... \\
\hline
 & \textbf{1.   \newline zur Feststellung einer seelischen  \newline Störung mit Krankheitswert die  \newline Stellungnahme einer }\textbf{\textcolor{diffred}{in \S{} 35a Absatz 1a  \newline Satz 1 genannten Person einholen,}} & \textbf{1.   \newline zur Feststellung einer seelischen  \newline Störung mit Krankheitswert die  \newline Stellungnahme }\textbf{\textcolor{diffgreen}{eines Arztes für Kinder-  \newline und Jugendpsychiatrie und - \newline psychotherapie, eines Kinder- oder  \newline Jugendlichenpsychotherapeuten, eines  \newline Fachpsychotherapeuten, eines  \newline Psychotherapeuten mit}}\textbf{ einer} \newline \textbf{\textcolor{diffgreen}{Weiterbildung für die Behandlung von  \newline Kindern und Jugendlichen oder eines  \newline Arztes oder eines psychologischen  \newline Psychotherapeuten, der über besondere  \newline Erfahrung auf dem Gebiet seelischer  \newline Störungen bei Kindern und  \newline Jugendlichen verfügt,}} \\
\hline
... & ... & ... \\
\hline
 & \textbf{a)   \newline über eine Betriebserlaubnis nach \S{} 45  \newline für eine Einrichtung im Inland verfügt, in  \newline der Hilfe zur Erziehung erbracht w}\textbf{\textcolor{diffred}{i}}\textbf{rd, } & \textbf{a)   \newline über eine Betriebserlaubnis nach \S{} 45  \newline für eine Einrichtung im Inland verfügt, in  \newline der Hilfe zur Erziehung }\textbf{\textcolor{diffgreen}{oder Leistungen  \newline der Eingliederungshilfe}}\textbf{ erbracht  \newline w}\textbf{\textcolor{diffgreen}{e}}\textbf{rd}\textbf{\textcolor{diffgreen}{en}}\textbf{, } \\
\hline
... & ... & ... \\
\hline
 & \textbf{c)   \newline mit der Erbringung der Hilfen nur  \newline Fachkräfte nach \S{} 72 Absatz 1 betraut, } & \textbf{c)   \newline mit der Erbringung der Hilfen }\textbf{\textcolor{diffgreen}{oder  \newline Leistungen}}\textbf{ nur Fachkräfte nach \S{} 72  \newline Absatz 1 betraut, } \\
\hline
... & ... & ... \\
\hline
 & \textbf{(3) Überprüfung und Fortschreibung des} \newline \textbf{\textcolor{diffred}{Hilfeplans}}\textbf{ sollen nach Maßgabe von \S{}  \newline 36 Absatz 2 Satz 2 am Ort der  \newline Leistungserbringung unter Beteiligung  \newline des Kindes oder des Jugendlichen  \newline erfolgen. Unabhängig von der  \newline Überprüfung und Fortschreibung des  \newline Hilfeplans nach Satz 1 soll der Träger  \newline der öffentlichen Jugendhilfe nach den  \newline Erfordernissen im Einzelfall an Ort und  \newline Stelle überprüfen, ob die Anforderungen  \newline nach Absatz 2 Nummer 2 Buchstabe b  \newline und c sowie Nummer 3 weiter erfüllt  \newline sind. } & \textbf{(3) Überprüfung und Fortschreibung des} \newline \textbf{\textcolor{diffgreen}{Hilfe- und Leistungsplans}}\textbf{ sollen nach  \newline Maßgabe von \S{} 36}\textbf{\textcolor{diffgreen}{a}}\textbf{ Absatz 2 Satz 2 am  \newline Ort der Leistungserbringung unter  \newline Beteiligung des Kindes oder des  \newline Jugendlichen erfolgen. Unabhängig von  \newline der Überprüfung und Fortschreibung  \newline des Hilfeplans nach Satz 1 soll der  \newline Träger der öffentlichen Jugendhilfe  \newline nach den Erfordernissen im Einzelfall an  \newline Ort und Stelle überprüfen, ob die  \newline Anforderungen nach Absatz 2 Nummer  \newline 2 Buchstabe b und c sowie Nummer 3  \newline weiter erfüllt sind. } \\
\hline
... & ... & ... \\
\hline
 & \textbf{1.   \newline den Beginn und das geplante Ende der  \newline Leistungserbringung im Ausland unter  \newline Angabe von Namen und Anschrift des  \newline Leistungserbringers, des  \newline Aufenthaltsorts des Kindes oder  \newline Jugendlichen sowie der Namen der mit  \newline der Erbringung der Hilfe betrauten  \newline Fachkräfte, } & \textbf{1.   \newline den Beginn und das geplante Ende der  \newline Leistungserbringung im Ausland unter  \newline Angabe von Namen und Anschrift des  \newline Leistungserbringers, des  \newline Aufenthaltsorts des Kindes oder  \newline Jugendlichen sowie der Namen der mit  \newline der Erbringung der Hilfe }\textbf{\textcolor{diffgreen}{oder Leistung}}\textbf{  \newline betrauten Fachkräfte, } \\
\hline
... & ... & ... \\
\hline
 & \textbf{b)   \newline des Haager Übereinkommens vom 19.  \newline Oktober 1996 über die Zuständigkeit,  \newline das anzuwendende Recht, die  \newline Anerkennung, Vollstreckung und  \newline Zusammenarbeit auf dem Gebiet der  \newline elterlichen Verantwortung und der  \newline Maßnahmen zum Schutz von Kindern  \newline zur Erfüllung der Maßgaben des Artikels} \newline \textbf{\textcolor{diffred}{33 zu übermitteln. Die  \newline erlaubniserteilende Behörde wirkt auf  \newline die unverzügliche Beendigung der  \newline Leistungserbringung im Ausland hin,  \newline wenn sich aus den Angaben nach Satz 1  \newline ergibt, dass die an die  \newline Leistungserbringung im Ausland  \newline gestellten gesetzlichen Anforderungen  \newline nicht erfüllt sind.}} & \textbf{b)   \newline des Haager Übereinkommens vom 19.  \newline Oktober 1996 über die Zuständigkeit,  \newline das anzuwendende Recht, die  \newline Anerkennung, Vollstreckung und  \newline Zusammenarbeit auf dem Gebiet der  \newline elterlichen Verantwortung und der  \newline Maßnahmen zum Schutz von Kindern  \newline zur Erfüllung der Maßgaben des Artikels } \\
\hline
 &  & \textbf{\textcolor{diffgreen}{33 zu übermitteln. Die  \newline erlaubniserteilende Behörde wirkt auf  \newline die unverzügliche Beendigung der  \newline Leistungserbringung im Ausland hin,  \newline wenn sich aus den Angaben nach Satz 1  \newline ergibt, dass die an die  \newline Leistungserbringung im Ausland  \newline gestellten gesetzlichen Anforderungen  \newline nicht erfüllt sind.}} \\
\hline
... & ... & ... \\
\hline
3.  \newline ein ausländisches Kind oder ein  \newline ausländischer Jugendlicher unbegleitet  \newline nach Deutschland kommt und sich weder  \newline Personensorge- noch  \newline Erziehungsberechtigte im Inland aufhalten. & 3. unverändert & 3. unverändert \\
\hline
... & ... & ... \\
\hline
Absatz 1 &  & unverändert \\
\hline
(2) Das Jugendamt hat während der  \newline Inobhutnahme unverzüglich das Kind oder  \newline den Jugendlichen umfassend und in einer  \newline verständlichen, nachvollziehbaren und  \newline wahrnehmbaren Form über diese  \newline Maßnahme aufzuklären, die Situation, die  \newline zur Inobhutnahme geführt hat, zusammen  \newline mit dem Kind oder dem Jugendlichen zu  \newline klären und Möglichkeiten der Hilfe und  \newline Unterstützung aufzuzeigen. Dem Kind oder  \newline dem Jugendlichen ist unverzüglich  \newline Gelegenheit zu geben, eine Person seines  \newline Vertrauens zu benachrichtigen. Das  \newline Jugendamt hat während der Inobhutnahme  für das Wohl des Kindes oder des  \newline Jugendlichen zu sorgen und dabei den  \newline notwendigen Unterhalt und die  \newline Krankenhilfe sicherzustellen; \S{} 39 Absatz 4  \newline Satz 2 gilt entsprechend. Das Jugendamt  \newline ist während der Inobhutnahme berechtigt,  \newline alle Rechtshandlungen vorzunehmen, die  \newline zum Wohl des Kindes oder Jugendlichen  \newline notwendig sind; der mutmaßliche Wille der  \newline Personensorge- oder der  \newline Erziehungsberechtigten ist dabei  \newline angemessen zu berücksichtigen. Im Fall  \newline des Absatzes 1 Satz 1 Nummer 3 gehört  \newline zu den Rechtshandlungen nach Satz 4, zu  \newline denen das Jugendamt verpflichtet ist,  \newline insbesondere die unverzügliche Stellung  \newline eines Asylantrags für das Kind oder den  \newline Jugendlichen in Fällen, in denen Tatsachen  \newline die Annahme rechtfertigen, dass das Kind  \newline oder der Jugendliche internationalen  \newline Schutz im Sinne des \S{} 1 Absatz 1 Nummer  \newline 2 des Asylgesetzes benötigt; dabei ist das  \newline Kind oder der Jugendliche zu beteiligen. &  & (2) Das Jugendamt hat während der  \newline Inobhutnahme unverzüglich das Kind oder  \newline den Jugendlichen umfassend und in einer  \newline verständlichen, nachvollziehbaren und  \newline wahrnehmbaren Form über diese  \newline Maßnahme aufzuklären, die Situation, die  \newline zur Inobhutnahme geführt hat, zusammen  \newline mit dem Kind oder dem Jugendlichen zu  \newline klären und Möglichkeiten der Hilfe und  \newline Unterstützung aufzuzeigen. Dem Kind oder  \newline dem Jugendlichen ist unverzüglich  \newline Gelegenheit zu geben, eine Person seines  \newline Vertrauens zu benachrichtigen. Das  \newline Jugendamt hat während der Inobhutnahme  für das Wohl des Kindes oder des  \newline Jugendlichen zu sorgen und dabei den  \newline notwendigen Unterhalt und die  \newline Krankenhilfe sicherzustellen; \S{} 39\textbf{\textcolor{diffgreen}{c}} Absatz  \newline 4 Satz 2 gilt entsprechend. Das Jugendamt  \newline ist während der Inobhutnahme berechtigt,  \newline alle Rechtshandlungen vorzunehmen, die  \newline zum Wohl des Kindes oder Jugendlichen  \newline notwendig sind; der mutmaßliche Wille der  \newline Personensorge- oder der  \newline Erziehungsberechtigten ist dabei  \newline angemessen zu berücksichtigen. Im Fall  \newline des Absatzes 1 Satz 1 Nummer 3 gehört  \newline zu den Rechtshandlungen nach Satz 4, zu  \newline denen das Jugendamt verpflichtet ist,  \newline insbesondere die unverzügliche Stellung  \newline eines Asylantrags für das Kind oder den  \newline Jugendlichen in Fällen, in denen Tatsachen  \newline die Annahme rechtfertigen, dass das Kind  \newline oder der Jugendliche internationalen  \newline Schutz im Sinne des \S{} 1 Absatz 1 Nummer  \newline 2 des Asylgesetzes benötigt; dabei ist das  \newline Kind oder der Jugendliche zu beteiligen. \\
\hline
\rowcolor{gray!25}
\textbf{\S{} 42a} &  & \textbf{\S{} 42a} \\
\hline
Vorläufige Inobhutnahme von  \newline ausländischen Kindern und Jugendlichen  \newline nach unbegleiteter Einreise &  & Vorläufige Inobhutnahme von  \newline ausländischen Kindern und Jugendlichen  \newline nach unbegleiteter Einreise \\
\hline
Absätze 1 bis 3a &  & unverändert \\
\hline
(4) Das Jugendamt hat der nach  \newline Landesrecht für die Verteilung von  \newline unbegleiteten ausländischen Kindern und  \newline Jugendlichen zuständigen Stelle die  \newline vorläufige Inobhutnahme des Kindes oder  \newline des Jugendlichen innerhalb von sieben  \newline Werktagen nach Beginn der Maßnahme  \newline zur Erfüllung der in \S{} 42b genannten  \newline Aufgaben mitzuteilen. Zu diesem Zweck  \newline sind auch die Ergebnisse der Einschätzung  \newline nach Absatz 2 Satz 1 mitzuteilen. Die nach  \newline Landesrecht zuständige Stelle hat  \newline gegenüber dem Bundesverwaltungsamt  \newline innerhalb von drei Werktagen das Kind  \newline oder den Jugendlichen zur Verteilung  \newline anzumelden oder den Ausschluss der  \newline Verteilung anzuzeigen. &  & (4) Das Jugendamt hat der nach  \newline Landesrecht für die Verteilung von  \newline unbegleiteten ausländischen Kindern und  \newline Jugendlichen zuständigen Stelle die  \newline vorläufige Inobhutnahme des Kindes oder  \newline des Jugendlichen innerhalb von \textcolor{diffgreen}{einem  \newline Monat} nach Beginn der Maßnahme zur  \newline Erfüllung der in \S{} 42b genannten Aufgaben  \newline mitzuteilen. Zu diesem Zweck sind auch  \newline die Ergebnisse der Einschätzung nach  \newline Absatz 2 Satz 1 mitzuteilen. Die nach  \newline Landesrecht zuständige Stelle hat  \newline gegenüber dem Bundesverwaltungsamt  \newline innerhalb von drei Werktagen das Kind  \newline oder den Jugendlichen zur Verteilung  \newline anzumelden oder den Ausschluss der  \newline Verteilung anzuzeigen. \\
\hline
Absätze 5 bis 6 &  & unverändert \\
\hline
\rowcolor{gray!25}
\textbf{\S{} 42b} &  & \textbf{\S{} 42b} \\
\hline
Verfahren zur Verteilung unbegleiteter  \newline ausländischer Kinder und Jugendlicher &  & Verfahren zur Verteilung unbegleiteter  \newline ausländischer Kinder und Jugendlicher \\
\hline
Absätze 1 bis 3 &  & unverändert \\
\hline
(4) Die Durchführung eines  \newline Verteilungsverfahrens ist bei einem  \newline unbegleiteten ausländischen Kind oder  \newline Jugendlichen ausgeschlossen, wenn   \newline 1. dadurch dessen Wohl gefährdet würde,   \newline 2. dessen Gesundheitszustand die  \newline Durchführung eines Verteilungsverfahrens  \newline innerhalb von 14 Werktagen nach Beginn  \newline der vorläufigen Inobhutnahme gemäß \S{}  \newline 42a nicht zulässt,  \newline 3.   dessen Zusammenführung mit einer  \newline verwandten Person kurzfristig erfolgen  \newline kann, zum Beispiel aufgrund der  \newline Verordnung (EU) Nr. 604/2013 des  \newline Europäischen Parlaments und des Rates  \newline vom 26. Juni 2013 zur Festlegung der  \newline Kriterien und Verfahren zur Bestimmung  \newline des Mitgliedstaats, der für die Prüfung  \newline eines von einem Drittstaatsangehörigen  \newline oder Staatenlosen in einem Mitgliedstaat  \newline gestellten Antrags auf internationalen  \newline Schutz zuständig ist (ABl. L 180 vom  \newline 29.6.2013, S. 31), und dies dem Wohl des  \newline Kindes entspricht oder   \newline 4.   die Durchführung des  \newline Verteilungsverfahrens nicht innerhalb von  \newline einem Monat nach Beginn der vorläufigen  \newline Inobhutnahme erfolgt. &  & (4) Die Durchführung eines  \newline Verteilungsverfahrens ist bei einem  \newline unbegleiteten ausländischen Kind oder  \newline Jugendlichen ausgeschlossen, wenn   \newline 1. dadurch dessen Wohl gefährdet würde,   \newline 2. dessen Gesundheitszustand die  \newline Durchführung eines Verteilungsverfahrens  \newline innerhalb von 14 Werktagen nach Beginn  \newline der vorläufigen Inobhutnahme gemäß \S{}  \newline 42a nicht zulässt,  \newline 3.   dessen Zusammenführung mit einer  \newline verwandten Person kurzfristig erfolgen  \newline kann, zum Beispiel aufgrund der  \newline Verordnung (EU) Nr. 604/2013 des  \newline Europäischen Parlaments und des Rates  \newline vom 26. Juni 2013 zur Festlegung der  \newline Kriterien und Verfahren zur Bestimmung  \newline des Mitgliedstaats, der für die Prüfung  \newline eines von einem Drittstaatsangehörigen  \newline oder Staatenlosen in einem Mitgliedstaat  \newline gestellten Antrags auf internationalen  \newline Schutz zuständig ist (ABl. L 180 vom  \newline 29.6.2013, S. 31), und dies dem Wohl des  \newline Kindes entspricht oder   \newline 4.   die Durchführung des  \newline Verteilungsverfahrens nicht innerhalb von  \newline \textcolor{diffgreen}{zwei Monaten} nach Beginn der vorläufigen  \newline Inobhutnahme erfolgt. \\
\hline
Absätze 5 bis 8 &  & unverändert \\
\hline
\rowcolor{gray!25}
\textbf{\S{} 42e} &  & \textbf{\S{} 42e} \\
\hline
Berichtspflicht &  & \textcolor{diffgreen}{Auf}e\textcolor{diffgreen}{nt}h\textcolor{diffgreen}{al}t \\
\hline
Die Bundesregierung hat dem Deutschen  \newline Bundestag jährlich einen Bericht über die  \newline Situation unbegleiteter ausländischer  \newline Minderjähriger in Deutschland vorzulegen. &  & \textcolor{diffgreen}{Unbegleitete ausländische Jugendliche  \newline sind verpflichtet, ihren gewöhnlichen  \newline Aufenthalt zu nehmen in} dem \textcolor{diffgreen}{Bereich des  \newline aufgrund der Zuweisungsentscheidung der  \newline zuständigen Landesbehörde nach \S{} 88a} \\
\hline
 &  & \textcolor{diffgreen}{Absatz 2 Satz 1 zuständigen Jugendamts.  \newline Über Folgen eines Verstoßes sind  \newline unbegleitete ausländische Jugendliche  \newline aufzuklären.} \\
\hline
\rowcolor{gray!25}
\textbf{\S{} 42f} &  & \textbf{\S{} 42f} \\
\hline
Behördliches Verfahren zur  \newline Altersfeststellung &  & Behördliches Verfahren zur  \newline Altersfeststellung \\
\hline
(1) Das Jugendamt hat im Rahmen der  \newline vorläufigen Inobhutnahme der  \newline ausländischen Person gemäß \S{} 42a deren  \newline Minderjährigkeit durch Einsichtnahme in  \newline deren Ausweispapiere festzustellen oder  \newline hilfsweise mittels einer qualifizierten  \newline Inaugenscheinnahme einzuschätzen und  \newline festzustellen. \S{} 8 Absatz 1 und \S{} 42 Absatz  \newline 2 Satz 2 sind entsprechend anzuwenden. &  & (1) Das Jugendamt hat im Rahmen der  \newline vorläufigen Inobhutnahme der  \newline ausländischen Person gemäß \S{} 42a deren  \newline Minderjährigkeit durch Einsichtnahme in  \newline deren Ausweispapiere festzustellen oder  \newline hilfsweise mittels einer qualifizierten  \newline Inaugenscheinnahme einzuschätzen und  \newline festzustellen. \textbf{\textcolor{diffgreen}{Die betroffene Person ist  \newline über die Folgen der Altersbestimmung  \newline und einer Verweigerung einer  \newline Mitwirkung aufzuklären.}} \S{} 8 Absatz 1 und  \newline \S{} 42 Absatz 2 Satz 2 sind entsprechend  \newline anzuwenden. \\
\hline
(2) Auf Antrag des Betroffenen oder seines  \newline Vertreters oder von Amts wegen hat das  \newline Jugendamt in Zweifelsfällen eine ärztliche  \newline Untersuchung zur Altersbestimmung zu  \newline veranlassen. Ist eine ärztliche  \newline Untersuchung durchzuführen, ist die  \newline betroffene Person durch das Jugendamt  \newline umfassend über die  \newline Untersuchungsmethode und über die  \newline möglichen Folgen der Altersbestimmung  \newline aufzuklären. Ist die ärztliche Untersuchung  \newline von Amts wegen durchzuführen, ist die  \newline betroffene Person zusätzlich über die  \newline Folgen einer Weigerung, sich der ärztlichen  \newline Untersuchung zu unterziehen, aufzuklären;  \newline die Untersuchung darf nur mit Einwilligung  \newline der betroffenen Person und ihres Vertreters  \newline durchgeführt werden. Die \S{}\S{} 60, 62 und 65  \newline bis 67 des Ersten Buches sind  \newline entsprechend anzuwenden. &  & (2) \textbf{Auf Antrag }\textbf{\textcolor{diffgreen}{der betroffenen Person}}\textbf{  \newline oder i}\textbf{\textcolor{diffgreen}{hr}}\textbf{es Vertreters oder von Amts  \newline wegen hat das Jugendamt eine ärztliche  \newline Untersuchung zur Altersbestimmung zu} \newline \textbf{\textcolor{diffgreen}{veranlassen, wenn nicht mit an  \newline Sicherheit grenzender  \newline Wahrscheinlichkeit nach Absatz 1}}\textbf{ die} \newline \textbf{\textcolor{diffgreen}{Volljährigkeit der ausländischen Person  \newline festgestellt wurde. Die}}\textbf{ betroffene  \newline Person }\textbf{\textcolor{diffgreen}{ist}}\textbf{ durch das Jugendamt  \newline umfassend über die  \newline Untersuchungsmethode aufzuklären.} \newline \textbf{\textcolor{diffgreen}{Die}}\textbf{ Untersuchung darf nur mit  \newline Einwilligung der betroffenen Person und  \newline ihres Vertreters durchgeführt werden.}  \newline Die \S{}\S{} 60, 62 und 65 bis 67 des Ersten  \newline Buches sind entsprechend anzuwenden. \\
\hline
(3) Widerspruch und Klage gegen die  \newline Entscheidung des Jugendamts, aufgrund  \newline der Altersfeststellung nach dieser Vorschrift  \newline die vorläufige Inobhutnahme nach \S{} 42a  \newline oder die Inobhutnahme nach \S{} 42 Absatz 1  \newline Satz 1 Nummer 3 abzulehnen oder zu  \newline beenden, haben keine aufschiebende  Wirkung. Landesrecht kann bestimmen,  \newline dass gegen diese Entscheidung Klage  \newline ohne Nachprüfung in einem Vorverfahren  \newline nach \S{} 68 der Verwaltungsgerichtsordnung  \newline erhoben werden kann. &  & \textbf{(3) }\textbf{\textcolor{diffgreen}{Eine}}\textbf{ Klage gegen die Entscheidung  \newline des Jugendamts, aufgrund der  \newline Altersfeststellung nach dieser Vorschrift  \newline die vorläufige Inobhutnahme nach \S{} 42a  \newline oder die Inobhutnahme nach \S{} 42  \newline Absatz 1 Satz 1 Nummer 3 abzulehnen  \newline oder zu beenden, ha}\textbf{\textcolor{diffgreen}{t}}\textbf{ keine } \textbf{aufschiebende Wirkung. }\textbf{\textcolor{diffgreen}{Das}}\textbf{  \newline Vorverfahren nach \S{} 68 der  \newline Verwaltungsgerichtsordnung }\textbf{\textcolor{diffgreen}{entfällt.}} \\
\hline
... & ... & ... \\
\hline
\textbf{Erlaubnis zur Vollzeitpflege } & \textbf{Erlaubnis zur Vollzeitpflege } & \textbf{Erlaubnis zur Vollzeitpflege } \\
\hline
... & ... & ... \\
\hline
\textbf{(3) Das Jugendamt soll den  \newline Erfordernissen des Einzelfalls  \newline entsprechend an Ort und Stelle  \newline überprüfen, ob die Voraussetzungen für  \newline die Erteilung der Erlaubnis weiter  \newline bestehen. Ist das Wohl des Kindes oder  \newline des Jugendlichen in der Pflegestelle  \newline gefährdet und ist die Pflegeperson nicht  \newline bereit oder in der Lage, die Gefährdung  \newline abzuwenden, so ist die Erlaubnis  \newline zurückzunehmen oder zu widerrufen. } & (3) unverändert & (3) unverändert \\
\hline
... & ... & ... \\
\hline
1.  \newline der Träger die für den Betrieb der  \newline Einrichtung erforderliche Zuverlässigkeit  \newline besitzt, & 1. unverändert & 1.   \newline d\textcolor{diffgreen}{i}e Träger die für den Betrieb der  \newline Einrichtung erforderliche Zuverlässigkeit  \newline \textcolor{diffgreen}{besitzt }\textbf{\textcolor{diffgreen}{und die Gewähr für eine den  \newline Zielen des Grundgesetzes förderliche  \newline Arbeit bietet,}} \\
\hline
2.  \newline die dem Zweck und der Konzeption der  \newline Einrichtung entsprechenden räumlichen,  \newline fachlichen, wirtschaftlichen und personellen  \newline Voraussetzungen für den Betrieb erfüllt  \newline sind und durch den Träger gewährleistet  \newline werden, & 2. unverändert & 2.   \newline die dem Zweck und der Konzeption der  \newline Einrichtung entsprechenden räumlichen,  \newline fachlichen, wirtschaftlichen und personellen  \newline Voraussetzungen für den Betrieb erfüllt  \newline sind und durch den Träger gewährleistet  \newline werden, \textcolor{diffgreen}{wobei sich die personellen  \newline Voraussetzungen auch nach den jeweils in  \newline der Einrichtung wahrzunehmenden  \newline Funktionen richten,} \\
\hline
... & ... & ... \\
\hline
\rowcolor{gray!25}
\textbf{\S{} 58a} &  & \textbf{\S{} 58a } \\
\hline
 &  & \textbf{\textcolor{diffgreen}{Bestätigung eines Todes eines  \newline sorgeberechtigten Elternteils}} \\
\hline
 &  & \textbf{\textcolor{diffgreen}{Ein Elternteil, dem die elterliche Sorge  \newline oder ein Teil der elterlichen Sorge  \newline bislang gemeinsam mit dem anderen  \newline Elternteil zustand und dem infolge des  \newline Todes des anderen Elternteils die  \newline elterliche Sorge oder ein Teil der  \newline elterlichen Sorge allein zusteht, kann  \newline bei dem nach \S{} 87c Absatz 6 Satz 2  \newline zuständigen Jugendamt hierüber eine  \newline Bestätigung beantragen. Das  \newline Jugendamt kann die Bestätigung  \newline ablehnen, wenn berechtigte Zweifel an  \newline der alleinigen Sorge des Elternteils oder  \newline der Anwendbarkeit deutschen Rechts  \newline bestehen.}} \\
\hline
... & ... & ... \\
\hline
(2) Soweit von der freien Jugendhilfe  \newline Einrichtungen, Dienste und  \newline Veranstaltungen geschaffen werden, um  \newline die Gewährung von Leistungen nach  \newline diesem Buch zu ermöglichen, kann die  \newline Förderung von der Bereitschaft abhängig  \newline gemacht werden, diese Einrichtungen,  \newline Dienste und Veranstaltungen nach  \newline Maßgabe der Jugendhilfeplanung und  \newline unter Beachtung der in \S{} 9 genannten  \newline Grundsätze anzubieten. \S{} 4 Absatz 1 bleibt  \newline unberührt. & (2) unverändert & (2) unverändert \\
\hline
... & ... & ... \\
\hline
(1) Werden Einrichtungen und Dienste der  \newline Träger der freien Jugendhilfe in Anspruch  \newline genommen, so sind Vereinbarungen über  \newline die Höhe der Kosten der Inanspruchnahme  \newline sowie über Inhalt, Umfang und Qualität der  \newline Leistung, über Grundsätze und Maßstäbe  \newline für die Bewertung der Qualität der Leistung  \newline und über geeignete Maßnahmen zu ihrer  \newline Gewährleistung zwischen der öffentlichen  \newline und der freien Jugendhilfe anzustreben. Zu  \newline den Grundsätzen und Maßstäben für die  \newline Bewertung der Qualität der Leistung nach  \newline Satz 1 zählen auch Qualitätsmerkmale für  \newline die inklusive Ausrichtung der  \newline Aufgabenwahrnehmung und die  \newline Berücksichtigung der spezifischen  \newline Bedürfnisse von jungen Menschen mit  \newline Behinderungen. Das Nähere regelt das  \newline Landesrecht. Die \S{}\S{} 78a bis 78g bleiben  \newline unberührt. & (1) unverändert & (1) Werden Einrichtungen und Dienste der  \newline Träger der freien Jugendhilfe in Anspruch  \newline genommen, so sind Vereinbarungen über  \newline die Höhe der Kosten der Inanspruchnahme  \newline sowie über Inhalt, Umfang und Qualität der  \newline Leistung, über Grundsätze und Maßstäbe  \newline für die Bewertung der Qualität der Leistung  \newline und über geeignete Maßnahmen zu ihrer  \newline Gewährleistung zwischen \textbf{\textcolor{diffgreen}{dem Träger}}\textbf{ der  \newline öffentlichen }\textbf{\textcolor{diffgreen}{Jugendhilfe}}\textbf{ und }\textbf{\textcolor{diffgreen}{unter  \newline Berücksichtigung}}\textbf{ der} \newline \textbf{\textcolor{diffgreen}{Leistungsfähigkeit, Wirtschaftlichkeit  \newline und Sparsamkeit geeigneten}}\textbf{ freien} \newline \textbf{\textcolor{diffgreen}{Trä}}\textbf{ge}\textbf{\textcolor{diffgreen}{r}}\textbf{n }anzustreben. Zu den Grundsätzen  \newline und Maßstäben für die Bewertung der  \newline Qualität der Leistung nach Satz 1 zählen  \newline auch Qualitätsmerkmale für die inklusive  \newline Ausrichtung der Aufgabenwahrnehmung  \newline und die Berücksichtigung der spezifischen  \newline Bedürfnisse von jungen Menschen mit  \newline Behinderungen. Das Nähere regelt das  \newline Landesrecht. Die \S{}\S{} 78a bis 78g bleiben  \newline unberührt. \\
\hline
... & ... & ... \\
\hline
2.  \newline Leistungen in gemeinsamen Wohnformen  \newline für Mütter/Väter und Kinder (\S{} 19), & 2. unverändert & 2. unverändert \\
\hline
... & ... & ... \\
\hline
c)  \newline in intensiver sozialpädagogischer  \newline Einzelbetreuung (\S{} 35), sofern sie  \newline außerhalb der eigenen Familie erfolgt, & c) unverändert & c) unverändert \\
\hline
d)  \newline in sonstiger teilstationärer oder stationärer  \newline Form (\S{} 27), & d) \textcolor{diffred}{unverändert} & d)  \textcolor{diffgreen}{in sonstiger teilstationärer oder stationärer  \newline Form (\S{} 27}\textbf{\textcolor{diffgreen}{a}}\textcolor{diffgreen}{),} \\
\hline
... & ... & ... \\
\hline
- Aus Synopsis 2024 - \newline 1.  \newline Inhalt, Umfang und Qualität der  \newline Leistungsangebote  \newline - Aus Synopsis 2026 - \newline 1.  \newline Inhalt, Umfang und Qualität der  \newline Leistungsangebote  \newline (Leistungsvereinbarung), & 1. unverändert & 1.  \newline Inhalt, Umfang und Qualität \textbf{\textcolor{diffgreen}{einschließlich  \newline der Wirksamkeit}} der Leistungsangebote  \newline (Leistungsvereinbarung), \\
\hline
... & ... & ... \\
\hline
(2) Die Vereinbarungen sind mit den  \newline Trägern abzuschließen, die unter  \newline Berücksichtigung der Grundsätze der  \newline Leistungsfähigkeit, Wirtschaftlichkeit und  \newline Sparsamkeit zur Erbringung der Leistung  \newline geeignet sind. Vereinbarungen über die  \newline Erbringung von Auslandsmaßnahmen  \newline dürfen nur mit solchen Trägern  \newline abgeschlossen werden, die die Maßgaben  \newline nach \S{} 38 Absatz 2 Nummer 2 Buchstabe  \newline a bis d erfüllen. & (2) Die Vereinbarungen sind mit \textcolor{diffred}{den  \newline Trägern abzuschließen, die} unter  \newline Berücksichtigung \textcolor{diffred}{der Grundsätze} der  \newline Leistungsfähigkeit, Wirtschaftlichkeit und  \newline Sparsamkeit \textcolor{diffred}{zur Erbringung der Leistung  \newline geeignet sind.} Vereinbarungen über die  \newline Erbringung von Auslandsmaßnahmen  \newline dürfen nur mit solchen Trägern  \newline abgeschlossen werden, die die Maßgaben  \newline nach \textbf{\S{} 40 Absatz 2 Nummer 2  \newline Buchstabe a bis d} erfüllen. & (2) Die Vereinbarungen sind mit unter  \newline Berücksichtigung der Leistungsfähigkeit,  \newline Wirtschaftlichkeit und Sparsamkeit  \newline \textcolor{diffgreen}{geeigneten Trägern abzuschließen, die  \newline eine bedarfsdeckende Leistungserbringung  \newline nach den Besonderheiten des Einzelfalls  \newline unter Berücksichtigung des Wunsch- und  \newline Wahlrechts des Leistungsberechtigten  \newline nach \S{} 5 sicherstellen.} Vereinbarungen  \newline über die Erbringung von  \newline Auslandsmaßnahmen dürfen nur mit  \newline solchen Trägern abgeschlossen werden,  \newline die die Maßgaben nach \textbf{\S{} 40 Absatz 2  \newline Nummer 2 Buchstabe a bis d} erfüllen. \\
\hline
 &  & \textbf{\textcolor{diffgreen}{(2a) Die Ergebnisse der Vereinbarungen  \newline sind den Leistungsberechtigten in einer  \newline für sie verständlichen,  \newline nachvollziehbaren und wahrnehmbaren  \newline Form zugänglich zu machen.}} \\
\hline
... & ... & ... \\
\hline
 &  & \textbf{\textcolor{diffgreen}{(4) Liegen die Voraussetzungen für die  \newline Übernahme des Leistungsentgelts nach  \newline Absatz 1 und 3 vor, hat der  \newline Leistungserbringer, der eine bewilligte  \newline Leistung gegenüber dem  \newline Leistungsberechtigten erbringt,  \newline Anspruch auf Vergütung dieser  \newline Leistung gegenüber dem Träger der  \newline öffentlichen Jugendhilfe.}} \\
\hline
\rowcolor{gray!25}
\textbf{\S{} 78c} &  & \textbf{\S{} 78c} \\
\hline
Inhalt der Leistungs- und  \newline Entgeltvereinbarungen &  & Inhalt der Leistungs- und  \newline Entgeltvereinbarungen \\
\hline
(1) Die Leistungsvereinbarung muss die  \newline wesentlichen Leistungsmerkmale,  \newline insbesondere   \newline 1.   Art, Ziel und Qualität des  \newline Leistungsangebots,   \newline 2. den in der Einrichtung zu betreuenden  \newline Personenkreis,   \newline 3. die erforderliche sächliche und  \newline personelle Ausstattung,   \newline 4. die Qualifikation des Personals sowie   \newline 5. die betriebsnotwendigen Anlagen der  \newline Einrichtung   \newline festlegen. In die Vereinbarung ist  \newline aufzunehmen, unter welchen  \newline Voraussetzungen der Träger der  \newline Einrichtung sich zur Erbringung von  \newline Leistungen verpflichtet. Der Träger muss  \newline gewährleisten, dass die Leistungsangebote  \newline zur Erbringung von Leistungen nach \S{} 78a  \newline Absatz 1 geeignet sowie ausreichend,  \newline zweckmäßig und wirtschaftlich sind. &  & (1) Die Leistungsvereinbarung muss die  \newline wesentlichen Leistungsmerkmale,  \newline insbesondere   \newline 1.   Art, Ziel und Qualität des  \newline Leistungsangebots,   \newline 2. den in der Einrichtung zu betreuenden  \newline Personenkreis,   \newline 3. die erforderliche sächliche und  \newline personelle Ausstattung,   \newline 4. die Qualifikation des \textcolor{diffgreen}{Personals, }\textbf{\textcolor{diffgreen}{die sich  \newline nach der Zweckbestimmung der  \newline Einrichtung und den jeweils in ihr  \newline wahrzunehmenden Funktionen richtet,}}  \newline sowie   \newline 5. die betriebsnotwendigen Anlagen der  \newline Einrichtung   \newline festlegen. In die Vereinbarung ist  \newline aufzunehmen, unter welchen  \newline Voraussetzungen der Träger der  \newline Einrichtung sich zur Erbringung von  \newline Leistungen verpflichtet. Der Träger muss  \newline gewährleisten, dass die Leistungsangebote  \newline zur Erbringung von Leistungen nach \S{} 78a  \newline Absatz 1 geeignet sowie ausreichend,  \newline zweckmäßig und wirtschaftlich sind. \\
\hline
Absatz 2 &  & unverändert \\
\hline
\rowcolor{gray!25}
\textbf{\S{} 78g} &  & \textbf{\S{} 78g} \\
\hline
Absatz 1 &  & unverändert \\
\hline
(2) Kommt eine Vereinbarung nach \S{} 78b  \newline Absatz 1 innerhalb von sechs Wochen  \newline nicht zustande, nachdem eine Partei  \newline schriftlich zu Verhandlungen aufgefordert  \newline hat, so entscheidet die Schiedsstelle auf  \newline Antrag einer Partei unverzüglich über die  \newline Gegenstände, über die keine Einigung  \newline erreicht werden konnte. Gegen die  \newline Entscheidung ist der Rechtsweg zu den  \newline Verwaltungsgerichten gegeben. Die Klage  \newline richtet sich gegen eine der beiden  \newline Vertragsparteien, nicht gegen die  \newline Schiedsstelle. Einer Nachprüfung der  \newline Entscheidung in einem Vorverfahren bedarf  \newline es nicht. &  & (2) Kommt eine Vereinbarung nach \S{} 78b  \newline Absatz 1 innerhalb von sechs Wochen  \newline nicht zustande, nachdem eine Partei  \newline schriftlich zu Verhandlungen aufgefordert  \newline hat, so entscheidet die Schiedsstelle auf  \newline Antrag einer Partei unverzüglich über die  \newline Gegenstände, über die keine Einigung  \newline erreicht werden konnte. Gegen die  \newline Entscheidung ist der Rechtsweg zu den  \newline \textbf{\textcolor{diffgreen}{Sozi}}\textbf{algerichten }gegeben. Die Klage  \newline richtet sich gegen eine der beiden  \newline Vertragsparteien, nicht gegen die  \newline Schiedsstelle. Einer Nachprüfung der  \newline Entscheidung in einem Vorverfahren bedarf  \newline es nicht. \\
\hline
Absätze 3 und 4 &  & unverändert \\
\hline
\rowcolor{gray!25}
\textbf{\S{} 79} &  & \textbf{\S{} 79} \\
\hline
Gesamtverantwortung &  & Gesamtverantwortung \\
\hline
Absätze 1 und 2 &  & unverändert \\
\hline
(3) Die Träger der öffentlichen Jugendhilfe  \newline haben für eine ausreichende Ausstattung  \newline der Jugendämter und der  \newline Landesjugendämter einschließlich der  \newline Möglichkeit der Nutzung digitaler Geräte zu  \newline sorgen; hierzu gehört auch eine dem  \newline Bedarf entsprechende Zahl von  \newline Fachkräften. Zur Planung und  \newline Bereitstellung einer bedarfsgerechten  \newline Personalausstattung ist ein Verfahren zur  \newline Personalbemessung zu nutzen. &  & (3) Die Träger der öffentlichen Jugendhilfe  \newline haben für eine ausreichende Ausstattung  \newline der Jugendämter und der  \newline Landesjugendämter einschließlich der  \newline Möglichkeit der Nutzung digitaler Geräte zu  \newline sorgen; hierzu gehört auch eine dem \textbf{\textcolor{diffgreen}{den  \newline jeweiligen Aufgabenbereichen und den  \newline darin jeweils wahrzunehmenden  \newline Funktionen}} entsprechende Zahl von  \newline Fachkräften. Zur Planung und  \newline Bereitstellung einer bedarfsgerechten  \newline Personalausstattung ist ein Verfahren zur  \newline Personalbemessung zu nutzen. \\
\hline
 &  & \textbf{\textcolor{diffgreen}{(4) Für die Bestimmung der örtlichen  \newline Zuständigkeit nach \S{}\S{} 86 bis 88a soll ein  \newline von der fachlich zuständigen obersten  \newline Bundesbehörde bereitgestelltes  \newline automatisiertes Prüfungssystem zur  \newline Anwendung kommen.}} \\
\hline
... & ... & ... \\
\hline
4.  \newline junge Menschen mit Behinderungen oder  \newline von Behinderung bedrohte junge  \newline Menschen mit jungen Menschen ohne  \newline Behinderung gemeinsam unter  \newline Berücksichtigung spezifischer  \newline Bedarfslagen gefördert werden können, & 4. unverändert & 4. unverändert \\
\hline
... & ... & ... \\
\hline
\rowcolor{gray!25}
 &  & \textbf{\textcolor{diffgreen}{\S{} 80a}} \\
\hline
 &  & \textbf{\textcolor{diffgreen}{Planung infrastruktureller  \newline Bildungsassistenz }} \textbf{\textcolor{diffgreen}{Infrastrukturelle Angebote zu einer  \newline wegen erzieherischen Bedarfs oder  \newline wegen einer Behinderung erforderlichen  \newline Anleitung und Begleitung von Kindern  \newline und Jugendlichen in  \newline Tageseinrichtungen nach \S{} 22 Absatz 1  \newline Satz 1 oder in Schulen oder  \newline Hochschulen sind Angebote, die nach  \newline Maßgabe von \S{} 80 unter Einbeziehung  \newline der nach Landesrecht für die Schulen  \newline und Hochschulen zuständigen  \newline Behörden geplant werden. Landesrecht  \newline regelt das Nähere, insbesondere auch  \newline die Beteiligung der in den Ländern für  \newline die Finanzierung der Schulen oder  \newline Hochschulen verantwortlichen Stellen  \newline an den Kosten der Angebote nach Satz  \newline 1.}} \\
\hline
\rowcolor{gray!25}
\textbf{\S{} 84} &  & \textbf{\S{} 84} \\
\hline
Jugendbericht &  & Jugendbericht \\
\hline
(1) Die Bundesregierung legt dem  \newline Deutschen Bundestag und dem Bundesrat  \newline in jeder Legislaturperiode einen Bericht  \newline über die Lage junger Menschen und die  \newline Bestrebungen und Leistungen der  \newline Jugendhilfe vor. Neben der  \newline Bestandsaufnahme und Analyse sollen die  \newline Berichte Vorschläge zur Weiterentwicklung  \newline der Jugendhilfe enthalten; jeder dritte  \newline Bericht soll einen Überblick über die  \newline Gesamtsituation der Jugendhilfe vermitteln. &  & (1) Die Bundesregierung legt dem  \newline Deutschen Bundestag und dem Bundesrat  \newline in jeder Legislaturperiode einen Bericht  \newline über die Lage junger Menschen und die  \newline Bestrebungen und Leistungen der  \newline Jugendhilfe vor. \textbf{Neben der  \newline Bestandsaufnahme und Analyse sollen  \newline die Berichte Vorschläge zur  \newline Weiterentwicklung der Jugendhilfe  \newline enthalten; jeder dritte Bericht soll einen  \newline Überblick über die Gesamtsituation der  \newline Jugendhilfe }\textbf{\textcolor{diffgreen}{vermitteln und auch die  \newline Situation unbegleiteter ausländischer  \newline Minderjähriger darstellen.}} \\
\hline
Absatz 2 &  & unverändert \\
\hline
... & ... & ... \\
\hline
1.  \newline die Beratung der örtlichen Träger und die  \newline Entwicklung von Empfehlungen zur  \newline Erfüllung der Aufgaben nach diesem Buch, & 1. unverändert & 1. unverändert \\
\hline
... & ... & ... \\
\hline
(5) Ist das Land überörtlicher Träger, so  \newline können durch Landesrecht bis zum 30.  \newline Juni 1993 einzelne seiner Aufgaben auf  \newline andere Körperschaften des öffentlichen  \newline Rechts, die nicht Träger der öffentlichen  \newline Jugendhilfe sind, übertragen werden. & (5) \textbf{Landesrecht kann }\textbf{\textcolor{diffred}{bis zum 1.1.2030}}\textbf{  \newline bestimmen, dass die Gewährung von  \newline Leistungen der Eingliederungshilfe für  \newline Kinder und Jugendliche mit  \newline Behinderungen im Sinne des \S{} 2 Absatz  \newline 2 Nummer 4 Buchstabe b auf den  \newline überörtlichen Träger der öffentlichen  \newline Jugendhilfe oder auf eine andere  \newline Körperschaft des öffentlichen Rechts  \newline übertragen wird. Im Falle einer  \newline Übertragung nach Satz 1 ist eine  \newline ortsnahe Wahrnehmung der Aufgaben  \newline nach \S{}\S{} 36 bis 39b unter Einbeziehung  \newline des örtlichen Trägers der öffentlichen  \newline Jugendhilfe sicherzustellen; \S{} 27 Absatz  \newline 5 bleibt unberührt.} & (5) \textbf{Landesrecht kann bestimmen, dass  \newline die Gewährung von Leistungen der  \newline Eingliederungshilfe für Kinder und  \newline Jugendliche mit Behinderungen im  \newline Sinne des \S{} 2 Absatz 2 Nummer 4  \newline Buchstabe b auf den überörtlichen  \newline Träger der öffentlichen Jugendhilfe oder  \newline auf eine andere Körperschaft des  \newline öffentlichen Rechts übertragen wird. Im } \textbf{Falle einer Übertragung nach Satz 1 ist  \newline eine ortsnahe Wahrnehmung der  \newline Aufgaben nach \S{}\S{} 36 bis 39b unter  \newline Einbeziehung des örtlichen Trägers der  \newline öffentlichen Jugendhilfe  \newline sicherzustellen; \S{} 27 Absatz 5 bleibt  \newline unberührt.} \\
\hline
 &  & \textbf{\textcolor{diffgreen}{(5) Aufgabenübertragungen, die durch  \newline ein Land als überörtlichem Träger bis  \newline zum 30. Juni 1993 auf andere  \newline Körperschaften des öffentlichen Rechts  \newline vorgenommen worden sind, bleiben  \newline wirksam.}} \\
\hline
\rowcolor{gray!25}
\textbf{\S{} 86} &  & \textbf{\S{} 86} \\
\hline
Absätze 1 bis 5 &  & unverändert \\
\hline
(6) Lebt ein Kind oder ein Jugendlicher  \newline zwei Jahre bei einer Pflegeperson und ist  \newline sein Verbleib bei dieser Pflegeperson auf  \newline Dauer zu erwarten, so ist oder wird  \newline abweichend von den Absätzen 1 bis 5 der  \newline örtliche Träger zuständig, in dessen  \newline Bereich die Pflegeperson ihren  \newline gewöhnlichen Aufenthalt hat. Er hat die  \newline Eltern und, falls den Eltern die  \newline Personensorge nicht oder nur teilweise  \newline zusteht, den Personensorgeberechtigten  \newline über den Wechsel der Zuständigkeit zu  \newline unterrichten. Endet der Aufenthalt bei der  \newline Pflegeperson, so endet die Zuständigkeit  \newline nach Satz 1. &  & (6) Lebt ein Kind oder ein Jugendlicher  \newline zwei Jahre bei einer Pflegeperson und ist  \newline sein Verbleib bei dieser Pflegeperson auf  \newline Dauer zu erwarten, so ist oder wird  \newline abweichend von den Absätzen 1 bis 5 der  \newline örtliche Träger zuständig, in dessen  \newline Bereich die Pflegeperson ihren  \newline gewöhnlichen Aufenthalt hat. Er hat die  \newline Eltern und, falls den Eltern die  \newline Personensorge nicht oder nur teilweise  \newline zusteht, den Personensorgeberechtigten  \newline über den Wechsel der Zuständigkeit zu  \newline unterrichten. Endet der Aufenthalt bei der  \newline Pflegeperson, so endet die Zuständigkeit  \newline nach Satz 1. \textbf{\textcolor{diffgreen}{Wurde bei der Auswahl der  \newline Pflegeperson der örtliche Träger, in  \newline dessen Bereich diese ihren  \newline gewöhnlichen Aufenthalt hat, nicht nach  \newline \S{} 37 Absatz 3 Satz 5 einbezogen, ist der  \newline Wechsel der Zuständigkeit nach Satz 1  \newline ausgeschlossen.}} \\
\hline
Absatz 7 &  & unverändert \\
\hline
... & ... & ... \\
\hline
\rowcolor{gray!25}
\textbf{\S{} 89d} &  & \textbf{\S{} 89d} \\
\hline
(1) Kosten, die ein örtlicher Träger  \newline aufwendet, sind vom Land zu erstatten,  \newline wenn   \newline 1.  innerhalb eines Monats nach der  \newline Einreise eines jungen Menschen oder  \newline eines Leistungsberechtigten nach \S{} 19  \newline Jugendhilfe gewährt wird und   \newline 2. sich die örtliche Zuständigkeit nach dem  \newline tatsächlichen Aufenthalt dieser Person  \newline oder nach der Zuweisungsentscheidung  \newline der zuständigen Landesbehörde richtet.   \newline Als Tag der Einreise gilt der Tag des  \newline Grenzübertritts, sofern dieser amtlich  \newline festgestellt wurde, oder der Tag, an dem  \newline der Aufenthalt im Inland erstmals  \newline festgestellt wurde, andernfalls der Tag der  \newline ersten Vorsprache bei einem Jugendamt.  \newline Die Erstattungspflicht nach Satz 1 bleibt  \newline unberührt, wenn die Person um Asyl  \newline nachsucht oder einen Asylantrag stellt. &  & (1) Kosten, die ein örtlicher Träger  \newline aufwendet, sind vom Land zu erstatten,  \newline wenn   \newline 1.  innerhalb eines Monats \textbf{nach }\textbf{\textcolor{diffgreen}{Kenntnis  \newline des tatsächlichen Aufenthalts}}\textbf{ eines} \newline \textbf{\textcolor{diffgreen}{eingereisten}}\textbf{ jungen Menschen }\textbf{\textcolor{diffgreen}{in  \newline dessen Bereich}} oder eines  \newline Leistungsberechtigten nach \S{} 19  \newline Jugendhilfe gewährt wird und   \newline 2. sich die örtliche Zuständigkeit nach dem  \newline tatsächlichen Aufenthalt dieser Person  \newline oder nach der Zuweisungsentscheidung  \newline der zuständigen Landesbehörde richtet.   \newline Als Tag der Einreise gilt der Tag des  \newline Grenzübertritts, sofern dieser amtlich  \newline festgestellt wurde, oder der Tag, an dem  \newline der Aufenthalt im Inland erstmals  \newline festgestellt wurde, andernfalls der Tag der  \newline ersten Vorsprache bei einem Jugendamt.  \newline Die Erstattungspflicht nach Satz 1 bleibt  \newline unberührt, wenn die Person um Asyl  \newline nachsucht oder einen Asylantrag stellt. \\
\hline
Absätze 2 bis 5 &  & unverändert \\
\hline
\rowcolor{gray!25}
\textbf{\S{} 89 f} &  & \textbf{\S{} 89 f} \\
\hline
Absatz 1 &  & unverändert \\
\hline
(2) Kosten unter 1 000 Euro werden nur bei  \newline vorläufigen Maßnahmen zum Schutz von  \newline Kindern und Jugendlichen (\S{} 89b), bei  \newline fortdauernder oder vorläufiger  \newline Leistungsverpflichtung (\S{} 89c) und bei  \newline Gewährung von Jugendhilfe nach der  \newline Einreise (\S{} 89d) erstattet. Verzugszinsen  \newline können nicht verlangt werden. &  & (2) Kosten unter 1 000 Euro werden nur bei  \newline vorläufigen Maßnahmen zum Schutz von  \newline Kindern und Jugendlichen (\S{} \textcolor{diffgreen}{89b) }\textbf{\textcolor{diffgreen}{und}} bei  \newline fortdauernder oder vorläufiger  \newline Leistungsverpflichtung (\S{} 89c) und bei  \newline Gewährung von Jugendhilfe nach der  \newline Einreise (\S{} 89d) erstattet. \textbf{\textcolor{diffgreen}{Bei Gewährung  \newline von Jugendhilfe an eingereiste junge  \newline Menschen (\S{} 89d) kann Landesrecht  \newline regeln, dass Kosten unter 1000 Euro  \newline erstattet werden. Landesrecht kann eine  \newline pauschale Abgeltung aufgewendeter  \newline Kosten vorsehen, soweit dies  \newline zweckmäßig ist.}} Verzugszinsen können  \newline nicht verlangt werden. \\
\hline
... & ... & ... \\
\hline
\textbf{Anwendungsbereich } & \textbf{Anwendungsbereich } & \textbf{Anwendungsbereich } \\
\hline
... & ... & ... \\
\hline
3.   \newline Eingliederungshilfe für seelisch behinderte  \newline Kinder und Jugendliche in  \newline Tageseinrichtungen und anderen  \newline teilstationären Einrichtungen nach \S{} 35a  \newline Absatz 2 Nummer 2 und & 3.   \newline Eingliederungshilfe \textbf{für Kinder und  \newline Jugendliche mit Behinderungen} in  \newline Tageseinrichtungen und anderen  \newline teilstationären Einrichtungen nach \S{} 35a  \newline Absatz \textcolor{diffred}{2} Nummer 2 und & 3.   \newline Eingliederungshilfe \textbf{für Kinder und  \newline Jugendliche mit Behinderungen} in  \newline Tageseinrichtungen und anderen  \newline teilstationären Einrichtungen nach \S{} 35a  \newline Absatz \textbf{\textcolor{diffgreen}{4}} Nummer 2 und \\
\hline
... & ... & ... \\
\hline
(3) Die Kosten umfassen auch die  \newline Aufwendungen für den notwendigen  \newline Unterhalt und die Krankenhilfe. & \textbf{(3) Ausgenommen von der  \newline Kostenbeitragspflicht nach den  \newline Absätzen 1 und 2 sind Leistungen zur} \newline \textbf{\textcolor{diffred}{Beschäftigung}}\textbf{ sowie Leistungen zum  \newline Erwerb und Erhalt praktischer  \newline Kenntnisse und Fähigkeiten nach \S{} 35f  \newline Absatz 2 Nummer 5, sowi}\textbf{\textcolor{diffred}{e}}\textbf{ diese der  \newline Vorbereitung auf die Teilhabe am  \newline Arbeitsleben nach \S{} 35e Absatz 1  \newline dienen. } & \textbf{(3) Ausgenommen von der  \newline Kostenbeitragspflicht nach den  \newline Absätzen 1 und 2 sind Leistungen zur} \newline \textbf{\textcolor{diffgreen}{Teilhabe am Arbeitsleben}}\textbf{ sowie  \newline Leistungen zum Erwerb und Erhalt  \newline praktischer Kenntnisse und Fähigkeiten  \newline nach \S{} 35f Absatz 2 Nummer 5, sow}\textbf{\textcolor{diffgreen}{e}}\textbf{i}\textbf{\textcolor{diffgreen}{t}}\textbf{  \newline diese der Vorbereitung auf die Teilhabe  \newline am Arbeitsleben nach \S{} 35e Absatz 1  \newline dienen. } \\
\hline
(4) Verwaltungskosten bleiben außer  \newline Betracht. & \textbf{(4) }\textbf{\textcolor{diffred}{U}}\textbf{n}\textbf{\textcolor{diffred}{ter}}\textbf{ den }\textbf{\textcolor{diffred}{w}}\textbf{eit}\textbf{\textcolor{diffred}{e}}\textbf{ren Leistungen  \newline werden Kostenbeiträge zu Leistungen  \newline zur Mobilität und Leistungen für  \newline Wohnraum (\S{} 35f Absatz 2 Nummer}\textbf{\textcolor{diffred}{n}}\textbf{ 1  \newline und 7) erhoben. } & \textbf{(4) }\textbf{\textcolor{diffgreen}{Nebe}}\textbf{n den }\textbf{\textcolor{diffgreen}{kostenb}}\textbf{eitr}\textbf{\textcolor{diffgreen}{agspflichtig}}\textbf{en  \newline Leistungen }\textbf{\textcolor{diffgreen}{nach Absatz 1 und 2}}\textbf{ werden  \newline Kostenbeiträge zu Leistungen zur  \newline Mobilität und Leistungen für Wohnraum  \newline (\S{} 35f Absatz 2 Nummer 1 und 7)  \newline erhoben. } \\
\hline
... & ... & ... \\
\hline
\textbf{Ausgestaltung der Heranziehung } & \textbf{Ausgestaltung der Heranziehung } & \textbf{Ausgestaltung der Heranziehung } \\
\hline
... & ... & ... \\
\hline
(2) Die Heranziehung erfolgt durch  \newline Erhebung eines Kostenbeitrags, der durch  \newline Leistungsbescheid festgesetzt wird;  \newline Elternteile werden getrennt herangezogen. & (2) Die Heranziehung erfolgt durch  \newline Erhebung eines Kostenbeitrags, der durch  \newline Leistungsbescheid festgesetzt \textcolor{diffred}{wird; }\textbf{\textcolor{diffred}{Eltern}}\textcolor{diffred}{  \newline werden }\textbf{\textcolor{diffred}{gemeinsam}}\textcolor{diffred}{ herangezogen, }\textbf{\textcolor{diffred}{wenn  \newline sie in einem gemeinsamen Haushalt  \newline leben, in dem auch der junge Mensch  \newline üblicherweise lebt, wenn er nicht im  \newline Rahmen einer Leistung über Tag und  \newline Nacht in einer betreuten Wohnform oder  \newline bei einer Pflegeperson untergebracht  \newline ist. Ist dies nicht der Fall, erfolgt eine  \newline getrennte Heranziehung.}} & (2) Die Heranziehung erfolgt durch  \newline Erhebung eines Kostenbeitrags, der durch  \newline Leistungsbescheid festgesetzt \textcolor{diffgreen}{wird.} \\
\hline
... & ... & ... \\
\hline
\rowcolor{gray!25}
 &  & \textbf{\textcolor{diffgreen}{\S{} 92a}} \\
\hline
 &  & \textbf{\textcolor{diffgreen}{Pauschale Heranziehung}} \\
\hline
 &  & \textbf{\textcolor{diffgreen}{(1) Zur Heranziehung zu den Kosten der  \newline Hilfen oder Leistungen nach \S{} 91 Absatz  \newline 1, bei denen Leistungen zum Unterhalt  \newline nach \S{} 39 umfassend gewährt werden,  \newline wird bei Elternteilen jeweils ein  \newline pauschaler Kostenbeitrag in Höhe von  \newline 50 Prozent der Regelbedarfsstufen im  \newline Sinne der Anlage des \S{} 28 des Zwölften  \newline Buches erhoben.}} \\
\hline
 &  & \textbf{\textcolor{diffgreen}{(2) Zur Heranziehung zu den Kosten der  \newline anderen Hilfen oder Leistungen nach \S{}  \newline 91 Absatz 1, bei denen keine Leistungen  \newline zum Unterhalt nach \S{} 39 gewährt  \newline werden, sowie zu den Kosten der Hilfen  \newline oder Leistungen nach \S{} 91 Absatz 2 }} \textbf{\textcolor{diffgreen}{wird bei den Elternteilen jeweils ein  \newline pauschaler Kostenbeitrag in Höhe von  \newline 20 Prozent der Regelbedarfsstufen im  \newline Sinne der Anlage des \S{} 28 des Zwölften  \newline Buches erhoben.}} \\
\hline
 &  & \textbf{\textcolor{diffgreen}{(3) Zur Heranziehung zu den Kosten der  \newline Leistungen nach \S{} 91 Absatz 4 wird bei  \newline den Elternteilen jeweils ein pauschaler  \newline Kostenbeitrag in Höhe von 10 Prozent  \newline der Regelbedarfsstufen im Sinne der  \newline Anlage des \S{} 28 des Zwölften Buches  \newline erhoben.}} \\
\hline
... & ... & ... \\
\hline
\textbf{Berechnung des Einkommens } & \textbf{Berechnung des Einkommens } & \textbf{Berechnung des Einkommens } \\
\hline
... & ... & ... \\
\hline
 & \textbf{(5) Kindergeld, das für den jungen  \newline Menschen, der die Leistung erhält,  \newline geleistet wird, wird dem maßgeblichen  \newline Einkommen im Sinne des Absatz 3 des} \newline \textbf{\textcolor{diffred}{Elternteils hinzugerechnet,}}\textbf{ der das  \newline Kindergeld }\textbf{\textcolor{diffred}{erhält.}}\textbf{ Erhält der junge  \newline Mensch selbst das }\textbf{\textcolor{diffred}{Kindergeld,}}\textbf{ so gilt }\textbf{\textcolor{diffred}{es  \newline unabhängig vom Einkommen als  \newline Einnahme des jungen Menschen.}} & \textbf{(5) Kindergeld, das für den jungen  \newline Menschen, der die Leistung erhält,  \newline geleistet wird, wird dem maßgeblichen  \newline Einkommen im Sinne des Absatz}\textbf{\textcolor{diffgreen}{es}}\textbf{ 3  \newline des }\textbf{\textcolor{diffgreen}{Elternteils,}}\textbf{ der das Kindergeld} \newline \textbf{\textcolor{diffgreen}{erhält, hinzugerechnet.}}\textbf{ Erhält der junge  \newline Mensch }\textbf{\textcolor{diffgreen}{für sich}}\textbf{ selbst das }\textbf{\textcolor{diffgreen}{Kindergeld  \newline nach \S{} 1 Absatz 2 des  \newline Bundeskindergeldgesetzes oder durch  \newline Abzweigung nach \S{} 74 Absatz 1 des  \newline Einkommensteuergesetzes,}}\textbf{ so gilt }\textbf{\textcolor{diffgreen}{\S{} 94  \newline Absatz 3.}} \\
\hline
... & ... & ... \\
\hline
\textbf{Umfang der Heranziehung } & \textbf{Umfang der Heranziehung } & \textbf{Umfang der Heranziehung } \\
\hline
... & ... & ... \\
\hline
(3) Werden Leistungen über Tag und Nacht  \newline außerhalb des Elternhauses erbracht und  \newline bezieht einer der Elternteile Kindergeld für  \newline den jungen Menschen, so hat dieser  \newline unabhängig von einer Heranziehung nach  \newline Absatz 1 Satz 1 und 2 einen Kostenbeitrag  \newline in Höhe des Kindergeldes zu zahlen. Zahlt  \newline der Elternteil den Kostenbeitrag nach Satz  \newline 1 nicht, so sind die Träger der öffentlichen  \newline Jugendhilfe insoweit berechtigt, das auf  \newline dieses Kind entfallende Kindergeld durch  \newline Geltendmachung eines  \newline Erstattungsanspruchs nach \S{} 74 Absatz 2  \newline des Einkommensteuergesetzes in  \newline Anspruch zu nehmen. Bezieht der Elternteil  \newline Kindergeld nach \S{} 1 Absatz 1 des  \newline Bundeskindergeldgesetzes, gilt Satz 2  \newline entsprechend. Bezieht der junge Mensch  \newline das Kindergeld selbst, gelten die Sätze 1  \newline und 2 entsprechend. Die Heranziehung der  \newline Elternteile erfolgt nachrangig zu der  \newline Heranziehung der jungen Menschen zu  \newline einem Kostenbeitrag in Höhe des  \newline Kindergeldes. & (3) Werden Leistungen über Tag und Nacht  \newline außerhalb des Elternhauses erbracht und  \newline \textbf{erhält der junge Mensch das Kindergeld} \newline \textbf{\textcolor{diffred}{selbst,}}\textbf{ hat er einen Kostenbeitrag in  \newline Höhe des }\textbf{\textcolor{diffred}{Kindegeldes als  \newline Kostenbeitrag}}\textbf{ zu zahlen.} Zahlt der \textbf{junge  \newline Mensch} den Kostenbeitrag nach Satz 1  \newline nicht, so sind die Träger der öffentlichen  \newline Jugendhilfe insoweit berechtigt, das auf  \newline \textcolor{diffred}{dieses Kind} entfallende Kindergeld durch  \newline Geltendmachung eines  \newline Erstattungsanspruchs nach \S{} 74 Absatz 2  \newline des Einkommensteuergesetzes in  \newline Anspruch zu nehmen. & \textbf{(3) Werden Leistungen über Tag und  \newline Nacht außerhalb des Elternhauses  \newline erbracht und erhält der junge Mensch  \newline das Kindergeld }\textbf{\textcolor{diffgreen}{für sich selbst nach \S{} 1  \newline Absatz 2 des  \newline Bundeskindergeldgesetzes oder durch  \newline Abzweigung nach \S{} 74 Absatz 1 des  \newline Einkommensteuergesetzes,}}\textbf{ hat er einen  \newline Kostenbeitrag in Höhe des }\textbf{\textcolor{diffgreen}{Kindergeldes}}\textbf{  \newline zu zahlen. Zahlt der junge Mensch den  \newline Kostenbeitrag nach Satz 1 nicht, so sind  \newline die Träger der öffentlichen Jugendhilfe  \newline insoweit berechtigt, das }\textbf{\textcolor{diffgreen}{nach Satz 1}}\textbf{ auf} \newline \textbf{\textcolor{diffgreen}{diesen jungen Menschen}}\textbf{ entfallende  \newline Kindergeld durch Geltendmachung  \newline eines Erstattungsanspruchs nach \S{} 74  \newline Absatz 2 des Einkommensteuergesetzes  \newline in Anspruch zu nehmen. }\textbf{\textcolor{diffgreen}{Der Einsatz  \newline von Geldleistungen nach \S{} 93 Absatz 1  \newline Satz 3 geht der Heranziehung nach Satz  \newline 1 vor. Kommt sowohl eine Heranziehung  \newline nach Satz 1 als auch nach \S{} 93 Absatz 1  \newline Satz 3 in Betracht, darf die Summe der  \newline Heranziehung den Höchstbetrag nach \S{}  \newline 93 Absatz 1 Satz 4 nicht überschreiten.}} \\
\hline
... & ... & ... \\
\hline
(5) Für die Festsetzung der Kostenbeiträge  \newline von Eltern werden nach  \newline Einkommensgruppen gestaffelte  \newline Pauschalbeträge durch Rechtsverordnung  \newline des zuständigen Bundesministeriums mit  \newline Zustimmung des Bundesrates bestimmt. & (5) Für die Festsetzung der Kostenbeiträge  \newline von Eltern \textbf{für Leistungen nach \S{} 91  \newline Absatz 1 und 2} werden nach  \newline Einkommensgruppen gestaffelte  \newline Pauschalbeträge durch Rechtsverordnung  \newline des zuständigen Bundesministeriums mit  \newline Zustimmung des Bundesrates bestimmt.  \newline \textbf{\textcolor{diffred}{Die Pauschalbeträge}}\textbf{ orientieren sich an} \newline \textbf{\textcolor{diffred}{der Höhe der}}\textbf{ Regelbedarfsstufen im  \newline Sinne der Anlage des \S{} 28 Zwölften  \newline Buches in einer Spanne von 0 bis }\textbf{\textcolor{diffred}{150  \newline Prozent. Bei getrennt herangezogenen  \newline Elternteilen}}\textbf{ darf die Summe beider  \newline Kostenbeiträge den Höchstbetrag der  \newline Kostenbeiträge aus der Anlage zur  \newline Rechtsverordnung nicht überschreiten.  \newline Bei Leistungen nach \S{} 41 steht die  \newline Heranziehung unter der Bedingung,  \newline dass Kindergeld für den jungen  \newline Menschen }\textbf{\textcolor{diffred}{bezogen wird;}}\textbf{ die  \newline Heranziehung ist auf die Höhe des  \newline Kindergeldes begrenzt. Für die  \newline Festsetzung der Kostenbeiträge von  \newline Eltern für Leistungen nach \S{} 91 Absatz }\textbf{\textcolor{diffred}{3}}\textbf{  \newline werden in der Rechtsverordnung die  \newline Anteile der }\textbf{\textcolor{diffred}{Kostenb}}\textbf{eteiligung an den  \newline Kosten der Leistung bestimmt.} & (5) Für die Festsetzung der Kostenbeiträge  \newline von Eltern \textbf{für Leistungen nach \S{} 91  \newline Absatz 1 und 2} werden nach  \newline Einkommensgruppen gestaffelte  \newline Pauschalbeträge durch Rechtsverordnung  \newline des zuständigen Bundesministeriums mit  \newline Zustimmung des Bundesrates bestimmt.  \newline \textbf{\textcolor{diffgreen}{Für Hilfen oder Leistungen nach \S{} 91  \newline Absatz 1, bei denen Leistungen zum  \newline Unterhalt nach \S{} 39 umfassend gewährt  \newline werden,}}\textbf{ orientieren sich }\textbf{\textcolor{diffgreen}{die  \newline Pauschalbeträge}}\textbf{ an }\textbf{\textcolor{diffgreen}{den}}\textbf{   \newline Regelbedarfsstufen im Sinne der Anlage  \newline des \S{} 28 }\textbf{\textcolor{diffgreen}{des}}\textbf{ Zwölften Buches in einer  \newline Spanne von 0 bis }\textbf{\textcolor{diffgreen}{100 Prozent; für  \newline andere Hilfen oder Leistungen nach \S{} 91  \newline Absatz 1 sowie Hilfen oder Leistungen  \newline nach \S{} 91 Absatz 2 orientieren sich die  \newline Pauschalbeträge an den für den  \newline häuslichen Lebensunterhalt vermuteten  \newline ersparten Aufwendungen. Werden beide  \newline Elternteile zu den Kosten herangezogen,}}\textbf{  \newline darf die Summe beider Kostenbeiträge  \newline den Höchstbetrag der Kostenbeiträge  \newline aus der Anlage zur Rechtsverordnung  \newline nicht überschreiten. Bei Leistungen  \newline nach \S{} 41 steht die Heranziehung }\textbf{\textcolor{diffgreen}{der  \newline Elternteile}}\textbf{ unter der Bedingung, dass} \newline \textbf{\textcolor{diffgreen}{die Elternteile}}\textbf{ Kindergeld für den jungen  \newline Menschen }\textbf{\textcolor{diffgreen}{erhalten;}}\textbf{ die Heranziehung  \newline ist auf die Höhe des Kindergeldes  \newline begrenzt. Für die Festsetzung der  \newline Kostenbeiträge von Eltern für  \newline Leistungen nach \S{} 91 Absatz }\textbf{\textcolor{diffgreen}{4}}\textbf{ werden  \newline in der Rechtsverordnung die Anteile der} \newline \textbf{\textcolor{diffgreen}{B}}\textbf{eteiligung an den Kosten der Leistung  \newline bestimmt. }\textbf{\textcolor{diffgreen}{Die Rechtsverordnung  \newline benennt für Kostenbeiträge nach \S{} 93  \newline Absatz 1 Satz 3 einen Höchstbetrag  \newline nach \S{} 93 Absatz 1 Satz 4, der sich an  \newline der Höhe des geleisteten  \newline Lebensunterhalts orientiert}}\textcolor{diffgreen}{.} \\
\hline
... & ... & ... \\
\hline
\rowcolor{gray!25}
\textbf{\S{} 97a} &  & \textbf{\S{} 97a} \\
\hline
Pflicht zur Auskunft &  & Pflicht zur Auskunft \\
\hline
Absätze 1 bis 3 &  & unverändert \\
\hline
(4) Kommt eine der nach den Absätzen 1  \newline und 2 zur Auskunft verpflichteten Personen  \newline ihrer Pflicht nicht nach oder bestehen  \newline tatsächliche Anhaltspunkte für die  \newline Unrichtigkeit ihrer Auskunft, so ist der  \newline Arbeitgeber dieser Person verpflichtet, dem  \newline örtlichen Träger über die Art des  \newline Beschäftigungsverhältnisses und den  \newline Arbeitsverdienst dieser Person Auskunft zu  \newline geben; Absatz 3 Satz 2 gilt entsprechend.  \newline Der zur Auskunft verpflichteten Person ist  \newline vor einer Nachfrage beim Arbeitgeber eine  \newline angemessene Frist zur Erteilung der  \newline Auskunft zu setzen. Sie ist darauf  \newline hinzuweisen, dass nach Fristablauf die  \newline erforderlichen Auskünfte beim Arbeitgeber  \newline eingeholt werden. &  & (4) \textbf{\textcolor{diffgreen}{Bestehen tatsächliche Anhaltspunkte  \newline für die Unrichtigkeit}}\textbf{ der }\textbf{\textcolor{diffgreen}{Auskunft einer}}\textbf{  \newline nach den Absätzen 1 und 2 zur Auskunft  \newline verpflichteten }\textbf{\textcolor{diffgreen}{Person,}}\textbf{ so ist der  \newline Arbeitgeber dieser Person verpflichtet,  \newline dem örtlichen Träger über die Art des  \newline Beschäftigungsverhältnisses und den  \newline Arbeitsverdienst dieser Person  \newline Auskunft zu geben; Absatz 3 Satz 2 gilt  \newline entsprechend.} Der zur Auskunft  \newline verpflichteten Person ist vor einer  \newline Nachfrage beim Arbeitgeber eine  \newline angemessene Frist zur Erteilung der  \newline Auskunft zu setzen. Sie ist darauf  \newline hinzuweisen, dass nach Fristablauf die  \newline erforderlichen Auskünfte beim Arbeitgeber  \newline eingeholt werden. \\
\hline
Absatz 5 &  & unverändert \\
\hline
\rowcolor{gray!25}
\textbf{\S{} 98} & \textbf{\textcolor{diffred}{\S{} 98 }} &  \\
\hline
\textbf{Zweck und Umfang der Erhebung } & \textbf{\textcolor{diffred}{Zweck und Umfang der Erhebung }} &  \\
\hline
(1) Zur Beurteilung der Auswirkungen der  \newline Bestimmungen dieses Buches und zu  \newline seiner Fortentwicklung sind laufende  \newline Erhebungen über & (1) unverändert &  \\
\hline
1.   \newline Kinder und tätige Personen in  \newline Tageseinrichtungen, & 1. unverändert &  \\
\hline
1a.   \newline Kinder in den Klassenstufen eins bis vier, & \textcolor{diffred}{1a. unverändert } &  \\
\hline
2.   \newline Kinder und tätige Personen in öffentlich  \newline geförderter Kindertagespflege, & 2. unverändert &  \\
\hline
3.   \newline Personen, die mit öffentlichen Mitteln  \newline geförderte Kindertagespflege gemeinsam  \newline oder auf Grund einer Erlaubnis nach \S{} 43  \newline Absatz 3 Satz 3 in Pflegestellen  \newline durchführen, und die von diesen betreuten  \newline Kinder, & 3. unverändert &  \\
\hline
4.   \newline die Empfänger & 4. unverändert &  \\
\hline
a)   \newline der Hilfe zur Erziehung, & \textcolor{diffred}{a) unverändert } &  \\
\hline
b)   \newline der Hilfe für junge Volljährige und & \textcolor{diffred}{b) unverändert } &  \\
\hline
c)   \newline der Eingliederungshilfe für seelisch  \newline behinderte Kinder und Jugendliche, & \textcolor{diffred}{c)   \newline der Eingliederungshilfe für Kinder und  \newline Jugendliche }\textbf{\textcolor{diffred}{mit Behinderungen oder  \newline drohenden Behinderungen}}\textcolor{diffred}{, } &  \\
\hline
5.   \newline Kinder und Jugendliche, zu deren Schutz  \newline vorläufige Maßnahmen getroffen worden  \newline sind, & 5. unverändert &  \\
\hline
6.   \newline Kinder und Jugendliche, die als Kind  \newline angenommen worden sind, & 6. unverändert &  \\
\hline
7.   \newline Kinder und Jugendliche, die unter  \newline Amtspflegschaft, Amtsvormundschaft oder  \newline Beistandschaft des Jugendamts stehen, & 7. unverändert &  \\
\hline
8.   \newline Kinder und Jugendliche, für die eine  \newline Pflegeerlaubnis erteilt worden ist, & 8. unverändert &  \\
\hline
9.   \newline Maßnahmen des Familiengerichts, & 9. unverändert &  \\
\hline
10.   \newline Angebote der Jugendarbeit nach \S{} 11  \newline sowie Fortbildungsmaßnahmen für  \newline ehrenamtliche Mitarbeiter anerkannter  \newline Träger der Jugendhilfe nach \S{} 74 Absatz 6, & 10. unverändert &  \\
\hline
11.   \newline die Träger der Jugendhilfe, die dort tätigen  \newline Personen und deren Einrichtungen mit  \newline Ausnahme der Tageseinrichtungen, & 11. unverändert &  \\
\hline
12.   \newline die Ausgaben und Einnahmen der  \newline öffentlichen Jugendhilfe sowie & 12. unverändert &  \\
\hline
13.   \newline Gefährdungseinschätzungen nach \S{} 8a als  \newline Bundesstatistik durchzuführen. & 13. unverändert &  \\
\hline
(2) Zur Verfolgung der gesellschaftlichen  \newline Entwicklung im Bereich der elterlichen  \newline Sorge sind im Rahmen der Kinder- und  \newline Jugendhilfestatistik auch laufende  \newline Erhebungen über Sorgeerklärungen  \newline durchzuführen. & (2) unverändert &  \\
\hline
\rowcolor{gray!25}
\textbf{\S{} 99} & \textbf{\textcolor{diffred}{\S{} 99 }} &  \\
\hline
\textbf{Erhebungsmerkmale } & \textbf{\textcolor{diffred}{Erhebungsmerkmale }} &  \\
\hline
(1) Erhebungsmerkmale bei den  \newline Erhebungen über Hilfe zur Erziehung nach  \newline den \S{}\S{} 27 bis 35, Eingliederungshilfe für  \newline seelisch behinderte Kinder und  \newline Jugendliche nach \S{} 35a und Hilfe für junge  \newline Volljährige nach \S{} 41 sind & \textcolor{diffred}{(1) Erhebungsmerkmale bei den  \newline Erhebungen über Hilfe zur Erziehung nach  \newline den \S{}\S{} 27 }\textbf{\textcolor{diffred}{Absatz 2, 27a}}\textcolor{diffred}{ bis 35,  \newline Eingliederungshilfe für Kinder und  \newline Jugendliche }\textbf{\textcolor{diffred}{mit Behinderungen oder  \newline drohenden Behinderungen}}\textcolor{diffred}{ nach \S{} 35a} \newline \textbf{\textcolor{diffred}{bis \S{} 35i}}\textcolor{diffred}{ und Hilfe für junge Volljährige  \newline nach \S{} 41 sind } &  \\
\hline
1. im Hinblick auf die Hilfe & \textcolor{diffred}{1. im Hinblick auf die Hilfe }\textbf{\textcolor{diffred}{oder Leistung}} &  \\
\hline
a)   \newline Art des Trägers des Hilfe durchführenden  \newline Dienstes oder der Hilfe durchführenden  \newline Einrichtung sowie bei Trägern der freien  \newline Jugendhilfe deren Verbandszugehörigkeit, & \textcolor{diffred}{a)   \newline Art des Trägers des }\textbf{\textcolor{diffred}{die}}\textcolor{diffred}{ Hilfe }\textbf{\textcolor{diffred}{oder die  \newline Leistung}}\textcolor{diffred}{ durchführenden Dienstes oder  \newline der }\textbf{\textcolor{diffred}{die}}\textcolor{diffred}{ Hilfe }\textbf{\textcolor{diffred}{oder die Leistung}}\textcolor{diffred}{  \newline durchführenden Einrichtung sowie bei  \newline Trägern der freien Jugendhilfe deren  \newline Verbandszugehörigkeit, } &  \\
\hline
b)   \newline Art der Hilfe, & \textcolor{diffred}{b)   \newline Art der Hilfe }\textbf{\textcolor{diffred}{oder Leistung}}\textcolor{diffred}{, } &  \\
\hline
c)   \newline Ort der Durchführung der Hilfe, & \textcolor{diffred}{c)   \newline Ort der Durchführung der Hilfe }\textbf{\textcolor{diffred}{oder  \newline Leistung}}\textcolor{diffred}{, } &  \\
\hline
d)   \newline Monat und Jahr des Beginns und Endes  \newline sowie Fortdauer der Hilfe, & \textcolor{diffred}{d)   \newline Monat und Jahr des Beginns und Endes  \newline sowie Fortdauer der Hilfe }\textbf{\textcolor{diffred}{oder Leistung}}\textcolor{diffred}{, } &  \\
\hline
e)   \newline familienrichterliche Entscheidungen zu  \newline Beginn der Hilfe, & \textcolor{diffred}{e)   \newline familienrichterliche Entscheidungen zu  \newline Beginn der Hilfe }\textbf{\textcolor{diffred}{oder Leistung}}\textcolor{diffred}{, } &  \\
\hline
f)   \newline Intensität der Hilfe, & \textcolor{diffred}{f)   \newline Intensität der Hilfe }\textbf{\textcolor{diffred}{oder Leistung}}\textcolor{diffred}{, } &  \\
\hline
g)   \newline Hilfe anregende Institutionen oder  \newline Personen, & \textcolor{diffred}{g)   \newline Hilfe }\textbf{\textcolor{diffred}{oder Leistung}}\textcolor{diffred}{ anregende  \newline Institutionen oder Personen, } &  \\
\hline
h)   \newline Gründe für die Hilfegewährung, & \textcolor{diffred}{h)   \newline Gründe für die Hilfe- }\textbf{\textcolor{diffred}{oder  \newline Leistungs}}\textcolor{diffred}{gewährung, } &  \\
\hline
i)   \newline Grund für die Beendigung der Hilfe, & \textcolor{diffred}{i)   \newline Grund für die Beendigung der Hilfe }\textbf{\textcolor{diffred}{oder  \newline Leistung}}\textcolor{diffred}{, } &  \\
\hline
j)   \newline vorangegangene  \newline Gefährdungseinschätzung nach \S{} 8a  \newline Absatz 1, & \textcolor{diffred}{j) unverändert } &  \\
\hline
k)   \newline Einleitung der Hilfe im Anschluss an eine  \newline vorläufige Maßnahme zum Schutz von  \newline Kindern und Jugendlichen im Fall des \S{} 42  \newline Absatz 1 Satz 1, & \textcolor{diffred}{k) unverändert } &  \\
\hline
l)   \newline gleichzeitige Inanspruchnahme einer  \newline weiteren Hilfe zur Erziehung, Hilfe für junge  \newline Volljährige oder Eingliederungshilfe bei  \newline einer seelischen Behinderung oder einer  \newline drohenden seelischen Behinderung sowie & \textcolor{diffred}{l)   \newline gleichzeitige Inanspruchnahme einer  \newline weiteren Hilfe zur Erziehung, Hilfe für junge  \newline Volljährige oder }\textbf{\textcolor{diffred}{Leistung der}}\textcolor{diffred}{  \newline Eingliederungshilfe sowie  } &  \\
\hline
2. im Hinblick auf junge Menschen & 2. unverändert &  \\
\hline
a)   \newline Geschlecht, & \textcolor{diffred}{a) unverändert } &  \\
\hline
b)   \newline Geburtsmonat und Geburtsjahr, & \textcolor{diffred}{b) unverändert } &  \\
\hline
c)   \newline Lebenssituation bei Beginn der Hilfe, & \textcolor{diffred}{c) unverändert } &  \\
\hline
d)   \newline ausländische Herkunft mindestens eines  \newline Elternteils, & \textcolor{diffred}{d) unverändert } &  \\
\hline
e)   \newline Deutsch als in der Familie vorrangig  \newline gesprochene Sprache, & \textcolor{diffred}{e) unverändert } &  \\
\hline
f)   \newline anschließender Aufenthalt, & \textcolor{diffred}{f) unverändert } &  \\
\hline
g)   \newline nachfolgende Hilfe; & \textcolor{diffred}{g)   \newline nachfolgende Hilfe }\textbf{\textcolor{diffred}{oder Leistung}}\textcolor{diffred}{; } &  \\
\hline
3.   \newline bei sozialpädagogischer Familienhilfe nach  \newline \S{} 31 und anderen familienorientierten  \newline Hilfen nach \S{} 27 zusätzlich zu den unter  \newline den Nummern 1 und 2 genannten  \newline Merkmalen & 3. unverändert &  \\
\hline
a)   \newline Geschlecht, Geburtsmonat und  \newline Geburtsjahr der in der Familie lebenden  \newline jungen Menschen sowie & \textcolor{diffred}{a) unverändert } &  \\
\hline
b)   \newline Zahl der außerhalb der Familie lebenden  \newline Kinder und Jugendlichen; & \textcolor{diffred}{b) unverändert } &  \\
\hline
4.   \newline für Hilfen außerhalb des Elternhauses nach  \newline \S{} 27 Absatz 1, 3 und 4, den \S{}\S{} 29 und 30,  \newline 32 bis 35a und 41 zusätzlich zu den unter  \newline den Nummern 1 und 2 genannten  \newline Merkmalen der Schulbesuch sowie das  \newline Ausbildungsverhältnis. & \textcolor{diffred}{4.   \newline für Hilfen }\textbf{\textcolor{diffred}{und Leistungen }}\textcolor{diffred}{außerhalb des  \newline Elternhauses nach \S{}}\textbf{\textcolor{diffred}{\S{}}}\textcolor{diffred}{ 27, }\textbf{\textcolor{diffred}{27a}}\textcolor{diffred}{ Absatz 1, 3  \newline und 4, den \S{}\S{} 29 und 30, 32 bis 35a und  \newline 41 zusätzlich zu den unter den Nummern 1  \newline und 2 genannten Merkmalen der  \newline Schulbesuch sowie das  \newline Ausbildungsverhältnis. } &  \\
\hline
(2) Erhebungsmerkmale bei den  \newline Erhebungen über vorläufige Maßnahmen  \newline zum Schutz von Kindern und Jugendlichen  \newline sind Kinder und Jugendliche, zu deren  \newline Schutz Maßnahmen nach \S{} 42 oder \S{} 42a  \newline getroffen worden sind, gegliedert nach & (2) unverändert &  \\
\hline
1.   \newline Art der Maßnahme, Art des Trägers der  \newline Maßnahme, Form der Unterbringung  \newline während der Maßnahme, hinweisgebender  \newline Institution oder Person, Zeitpunkt des  \newline Beginns und Dauer der Maßnahme,  \newline Durchführung aufgrund einer  \newline vorangegangenen  \newline Gefährdungseinschätzung nach \S{} 8a  \newline Absatz 1, Maßnahmeanlass, im  \newline Kalenderjahr bereits wiederholt  \newline stattfindende Inobhutnahme, Widerspruch  \newline der Personensorge- oder  \newline Erziehungsberechtigten gegen die  \newline Maßnahme, im Fall des Widerspruchs  \newline gegen die Maßnahme Herbeiführung einer  \newline Entscheidung des Familiengerichts nach \S{}  \newline 42 Absatz 3 Satz 2 Nummer 2, Grund für  \newline die Beendigung der Maßnahme,  \newline anschließendem Aufenthalt, Art der  \newline anschließenden Hilfe, & 1. unverändert &  \\
\hline
2.   \newline bei Kindern und Jugendlichen zusätzlich zu  \newline den unter Nummer 1 genannten  \newline Merkmalen nach Geschlecht, Altersgruppe  \newline zu Beginn der Maßnahme, ausländischer  \newline Herkunft mindestens eines Elternteils,  \newline Deutsch als in der Familie vorrangig  \newline gesprochene Sprache, Art des Aufenthalts  \newline vor Beginn der Maßnahme. & 2. unverändert &  \\
\hline
(3) Erhebungsmerkmale bei den  \newline Erhebungen über die Annahme als Kind  \newline sind & (3) unverändert &  \\
\hline
1.   \newline angenommene Kinder und Jugendliche,  \newline gegliedert & 1. unverändert &  \\
\hline
a)   \newline nach nationaler Adoption und  \newline internationaler Adoption nach \S{} 2a des  \newline Adoptionsvermittlungsgesetzes, & \textcolor{diffred}{a) unverändert } &  \\
\hline
b)   \newline nach Geschlecht, Geburtsdatum,  \newline Staatsangehörigkeit und Art des Trägers  \newline des Adoptionsvermittlungsdienstes, Datum  \newline des Adoptionsbeschlusses, & \textcolor{diffred}{b) unverändert } &  \\
\hline
c)   \newline nach Herkunft des angenommenen Kindes,  \newline Art der Unterbringung vor der  \newline Adoptionspflege, Geschlecht und  \newline Familienstand der Eltern oder des  \newline sorgeberechtigten Elternteils oder Tod der  \newline Eltern zu Beginn der Adoptionspflege  \newline sowie Ersetzung der Einwilligung zur  \newline Annahme als Kind, & \textcolor{diffred}{c) unverändert } &  \\
\hline
d)   \newline zusätzlich bei nationalen Adoptionen nach  \newline Datum des Beginns und Endes der  \newline Adoptionspflege und bei Unterbringung vor  \newline der Adoptionspflege in Pflegefamilien nach  \newline Datum des Beginns und Endes dieser  \newline Unterbringung sowie bei Annahme durch  \newline die vorherige Pflegefamilie nach Datum  \newline des Beginns und Endes dieser  \newline Unterbringung, & \textcolor{diffred}{d) unverändert } &  \\
\hline
e)   \newline zusätzlich bei der internationalen Adoption  \newline (\S{} 2a des Adoptionsvermittlungsgesetzes)  \newline nach Staatsangehörigkeit vor Ausspruch  \newline der Adoption, nach Herkunftsland und  \newline gewöhnlichem Aufenthalt vor der Adoption  \newline sowie nach Ausspruch der Adoption im  \newline Ausland oder Inland, & \textcolor{diffred}{e) unverändert } &  \\
\hline
f)   \newline nach Staatsangehörigkeit, Geschlecht und  \newline Familienstand der oder des Annehmenden  \newline sowie nach dem Verwandtschaftsverhältnis  \newline zu dem Kind, & \textcolor{diffred}{f) unverändert } &  \\
\hline
2.   \newline die Zahl der & 2. unverändert &  \\
\hline
a)   \newline ausgesprochenen und aufgehobenen  \newline Annahmen sowie der abgebrochenen  \newline Adoptionspflegen, gegliedert nach Art des  \newline Trägers des  \newline Adoptionsvermittlungsdienstes, & \textcolor{diffred}{a) unverändert } &  \\
\hline
b)   \newline vorgemerkten Adoptionsbewerber, die zur  \newline Annahme als Kind vorgemerkten und in  \newline Adoptionspflege untergebrachten Kinder  \newline und Jugendlichen zusätzlich nach ihrem  \newline Geschlecht, gegliedert nach Art des  \newline Trägers des  \newline Adoptionsvermittlungsdienstes, & \textcolor{diffred}{b) unverändert } &  \\
\hline
3.  \newline bei Anerkennungs- und  \newline Wirkungsfeststellung einer ausländischen  \newline Adoptionsentscheidung nach \S{} 2 des  \newline Adoptionswirkungsgesetzes sowie eines  \newline Umwandlungsausspruchs nach \S{} 3 des  \newline Adoptionswirkungsgesetzes die Zahl der & 3. unverändert &  \\
\hline
a)   \newline eingeleiteten Verfahren nach den \S{}\S{} 2 und  \newline 3 des Adoptionswirkungsgesetzes, & \textcolor{diffred}{a) unverändert } &  \\
\hline
b)   \newline beendeten Verfahren nach den \S{}\S{} 2 und 3  \newline des Adoptionswirkungsgesetzes, die  \newline ausländische Adoptionen nach \S{} 2a des  \newline Adoptionsvermittlungsgesetzes zum  \newline Gegenstand haben, gegliedert nach & \textcolor{diffred}{b) unverändert } &  \\
\hline
aa)   \newline dem Ergebnis des Verfahrens im Hinblick  \newline auf eine erfolgte und nicht erfolgte  \newline Vermittlung nach \S{} 2a Absatz 2 des  \newline Adoptionsvermittlungsgesetzes, & \textcolor{diffred}{aa) unverändert } &  \\
\hline
bb)   \newline dem Vorliegen einer Bescheinigung nach  \newline Artikel 23 des Haager Übereinkommens  \newline vom 29. Mai 1993 über den Schutz von  \newline Kindern und die Zusammenarbeit auf dem  \newline Gebiet der internationalen Adoption und & \textcolor{diffred}{bb) unverändert } &  \\
\hline
cc)   \newline der Verfahrensdauer. & \textcolor{diffred}{cc) unverändert } &  \\
\hline
(4) Erhebungsmerkmal bei den Erhebungen  \newline über die Amtspflegschaft und die  \newline Amtsvormundschaft sowie die  \newline Beistandschaft ist die Zahl der Kinder und  \newline Jugendlichen unter & (4) unverändert &  \\
\hline
1.   \newline gesetzlicher Amtsvormundschaft, & 1. unverändert &  \\
\hline
2.   \newline bestellter Amtsvormundschaft, & 2. unverändert &  \\
\hline
3.   \newline bestellter Amtspflegschaft sowie & 3. unverändert &  \\
\hline
4.   \newline Beistandschaft, & 4. unverändert &  \\
\hline
gegliedert nach Geschlecht, Art des  \newline Tätigwerdens des Jugendamts sowie nach  \newline deutscher und ausländischer  \newline Staatsangehörigkeit (Deutsche/Ausländer). & unverändert &  \\
\hline
(5) Erhebungsmerkmal bei den Erhebungen  \newline über & (5) unverändert &  \\
\hline
1.   \newline die Pflegeerlaubnis nach \S{} 43 ist die Zahl  \newline der Tagespflegepersonen, & 1. unverändert &  \\
\hline
2.   \newline die Pflegeerlaubnis nach \S{} 44 ist die Zahl  \newline der Kinder und Jugendlichen, gegliedert  \newline nach Geschlecht und Art der Pflege. & 2. unverändert &  \\
\hline
(6) Erhebungsmerkmale bei der Erhebung  \newline zum Schutzauftrag bei  \newline Kindeswohlgefährdung nach \S{} 8a sind  \newline Kinder und Jugendliche, bei denen eine  \newline Gefährdungseinschätzung nach Absatz 1  \newline vorgenommen worden ist, gegliedert & (6) unverändert &  \\
\hline
1.   \newline nach der hinweisgebenden Institution oder  \newline Person, der Art der Kindeswohlgefährdung,  \newline der Person, von der die Gefährdung  \newline ausgeht, dem Ergebnis der  \newline Gefährdungseinschätzung sowie  \newline wiederholter Meldung zu demselben Kind  \newline oder Jugendlichen im jeweiligen  \newline Kalenderjahr, & 1. unverändert &  \\
\hline
2.   \newline bei Kindern und Jugendlichen zusätzlich zu  \newline den in Nummer 1 genannten Merkmalen  \newline nach Geschlecht, Geburtsmonat,  \newline Geburtsjahr, ausländischer Herkunft  \newline mindestens eines Elternteils, Deutsch als in  \newline der Familie vorrangig gesprochene  \newline Sprache, Eingliederungshilfe und  \newline Aufenthaltsort des Kindes oder  \newline Jugendlichen zum Zeitpunkt der Meldung  \newline sowie den Altersgruppen der Eltern und der  \newline Inanspruchnahme einer Leistung gemäß  \newline den \S{}\S{} 16 bis 19 sowie 27 bis 35a und der  \newline Durchführung einer Maßnahme nach \S{} 42. & \textcolor{diffred}{2.   \newline bei Kindern und Jugendlichen zusätzlich zu  \newline den in Nummer 1 genannten Merkmalen  \newline nach Geschlecht, Geburtsmonat,  \newline Geburtsjahr, ausländischer Herkunft  \newline mindestens eines Elternteils, Deutsch als in  \newline der Familie vorrangig gesprochene  \newline Sprache, }\textbf{\textcolor{diffred}{Leistung der}}\textcolor{diffred}{ Eingliederungshilfe  \newline und Aufenthaltsort des Kindes oder  \newline Jugendlichen zum Zeitpunkt der Meldung  \newline sowie den Altersgruppen der Eltern und der  \newline Inanspruchnahme einer Leistung gemäß  \newline den \S{}\S{} 16 bis 19 sowie 27, }\textbf{\textcolor{diffred}{27a}}\textcolor{diffred}{ bis 35}\textbf{\textcolor{diffred}{i}}\textcolor{diffred}{ und  \newline der Durchführung einer Maßnahme nach \S{}  \newline 42. } &  \\
\hline
(6a) Erhebungsmerkmal bei den  \newline Erhebungen über Sorgeerklärungen und  \newline die gerichtliche Übertragung der  \newline gemeinsamen elterlichen Sorge nach \S{}  \newline 1626a Absatz 1 Nummer 3 des  \newline Bürgerlichen Gesetzbuchs ist die  \newline gemeinsame elterliche Sorge nicht  \newline miteinander verheirateter Eltern, gegliedert  \newline danach, ob Sorgeerklärungen beider Eltern  \newline vorliegen oder den Eltern die elterliche  \newline Sorge aufgrund einer gerichtlichen  \newline Entscheidung ganz oder zum Teil  \newline gemeinsam übertragen worden ist. & \textcolor{diffred}{(6a) unverändert } &  \\
\hline
(6b) Erhebungsmerkmal bei den  \newline Erhebungen über Maßnahmen des  \newline Familiengerichts ist die Zahl der Kinder und  \newline Jugendlichen, bei denen wegen einer  \newline Gefährdung ihres Wohls das  \newline familiengerichtliche Verfahren auf Grund  \newline einer Anrufung durch das Jugendamt nach  \newline \S{} 8a Absatz 2 Satz 1 oder \S{} 42 Absatz 3  \newline Satz 2 Nummer 2 oder auf andere Weise  \newline eingeleitet worden ist und & \textcolor{diffred}{(6b) unverändert } &  \\
\hline
1.   \newline den Personensorgeberechtigten auferlegt  \newline worden ist, Leistungen nach diesem Buch  \newline in Anspruch zu nehmen, & 1. unverändert &  \\
\hline
2.   \newline andere Gebote oder Verbote gegenüber  \newline den Personensorgeberechtigten oder  \newline Dritten ausgesprochen worden sind, & 2. unverändert &  \\
\hline
3.   \newline Erklärungen der  \newline Personensorgeberechtigten ersetzt worden  \newline sind, & 3. unverändert &  \\
\hline
4.  \newline die elterliche Sorge ganz oder teilweise  \newline entzogen und auf das Jugendamt oder  \newline einen Dritten als Vormund oder Pfleger  \newline übertragen worden ist, & 4. unverändert &  \\
\hline
gegliedert nach Geschlecht, Altersgruppen  \newline und zusätzlich bei Nummer 4 nach dem  \newline Umfang der übertragenen Angelegenheit.  \newline Zusätzlich sind die Fälle nach Geschlecht  \newline und Altersgruppen zu melden, in denen  \newline das Jugendamt insbesondere nach \S{} 8a  \newline Absatz 2 Satz 1 oder \S{} 42 Absatz 3 Satz 2  \newline Nummer 2 das Familiengericht anruft, weil  \newline es dessen Tätigwerden für erforderlich hält. & unverändert &  \\
\hline
(7) Erhebungsmerkmale bei den  \newline Erhebungen über Kinder und tätige  \newline Personen in Tageseinrichtungen sind & (7) unverändert &  \\
\hline
1.   \newline die Einrichtungen, gegliedert nach & 1. unverändert &  \\
\hline
a)   \newline der Art und Rechtsform des Trägers sowie  \newline bei Trägern der freien Jugendhilfe deren  \newline Verbandszugehörigkeit sowie besonderen  \newline Merkmalen, & \textcolor{diffred}{a) unverändert } &  \\
\hline
b)   \newline der Zahl der genehmigten Plätze, & \textcolor{diffred}{b) unverändert } &  \\
\hline
c)   \newline der Art und Anzahl der Gruppen, & \textcolor{diffred}{c) unverändert } &  \\
\hline
d)   \newline die Anzahl der Kinder insgesamt, & \textcolor{diffred}{d) unverändert } &  \\
\hline
e)   \newline Anzahl der Schließtage an regulären  \newline Öffnungstagen im vorangegangenen Jahr  \newline sowie & \textcolor{diffred}{e) unverändert } &  \\
\hline
f)   \newline Öffnungszeiten, & \textcolor{diffred}{f) unverändert } &  \\
\hline
2.   \newline für jede dort tätige Person & 2. unverändert &  \\
\hline
a)   \newline Geschlecht und Beschäftigungsumfang, & \textcolor{diffred}{a) unverändert } &  \\
\hline
b)   \newline für das pädagogisch und in der Verwaltung  \newline tätige Personal zusätzlich Geburtsmonat  \newline und Geburtsjahr, die Art des  \newline Berufsausbildungsabschlusses, Stellung im  \newline Beruf, Art der Beschäftigung und  \newline Arbeitsbereiche einschließlich  \newline Gruppenzugehörigkeit, Monat und Jahr des  \newline Beginns der Tätigkeit in der derzeitigen  \newline Einrichtung, & \textcolor{diffred}{b) unverändert } &  \\
\hline
3.   \newline für die dort geförderten Kinder & 3. unverändert &  \\
\hline
a)   \newline Geschlecht, Geburtsmonat und  \newline Geburtsjahr sowie Schulbesuch und  \newline Klassenstufe, & \textcolor{diffred}{a) unverändert } &  \\
\hline
b)   \newline ausländische Herkunft mindestens eines  \newline Elternteils, & \textcolor{diffred}{b) unverändert } &  \\
\hline
c)   \newline Deutsch als in der Familie vorrangig  \newline gesprochene Sprache, & \textcolor{diffred}{c) unverändert } &  \\
\hline
d)   \newline Betreuungszeit und Mittagsverpflegung, & \textcolor{diffred}{d) unverändert } &  \\
\hline
e)   \newline Eingliederungshilfe, & \textcolor{diffred}{e) unverändert } &  \\
\hline
f)   \newline Gruppenzugehörigkeit, & \textcolor{diffred}{f) unverändert } &  \\
\hline
g)   \newline Monat und Jahr der Aufnahme in der  \newline Tageseinrichtung. & \textcolor{diffred}{g) unverändert } &  \\
\hline
(7a)   \newline Erhebungsmerkmale bei den Erhebungen  \newline über Kinder in mit öffentlichen Mitteln  \newline geförderter Kindertagespflege sowie die die  \newline Kindertagespflege durchführenden  \newline Personen sind: & \textcolor{diffred}{(7a) unverändert } &  \\
\hline
1.   \newline für jede tätige Person & 1. unverändert &  \\
\hline
a)   \newline Geschlecht, Geburtsmonat und  \newline Geburtsjahr, & \textcolor{diffred}{a) unverändert } &  \\
\hline
b)   \newline Art und Umfang der Qualifikation, höchster  \newline allgemeinbildender Schulabschluss,  \newline höchster beruflicher Ausbildungs- und  \newline Hochschulabschluss, Anzahl der betreuten  \newline Kinder (Betreuungsverhältnisse am  \newline Stichtag) insgesamt und nach dem Ort der  \newline Betreuung, & \textcolor{diffred}{b) unverändert } &  \\
\hline
2.   \newline für die dort geförderten Kinder & 2. unverändert &  \\
\hline
a)   \newline Geschlecht, Geburtsmonat und  \newline Geburtsjahr sowie Schulbesuch, & \textcolor{diffred}{a) unverändert } &  \\
\hline
b)   \newline ausländische Herkunft mindestens eines  \newline Elternteils, & \textcolor{diffred}{b) unverändert } &  \\
\hline
c)   \newline Deutsch als in der Familie vorrangig  \newline gesprochene Sprache, & \textcolor{diffred}{c) unverändert } &  \\
\hline
d)   \newline Betreuungszeit und Mittagsverpflegung, & \textcolor{diffred}{d) unverändert } &  \\
\hline
e)   \newline Art und Umfang der öffentlichen  \newline Finanzierung und Förderung, & \textcolor{diffred}{e) unverändert } &  \\
\hline
f)   \newline Eingliederungshilfe, & \textcolor{diffred}{f) unverändert } &  \\
\hline
g)   \newline Verwandtschaftsverhältnis zur  \newline Pflegeperson, & \textcolor{diffred}{g) unverändert } &  \\
\hline
h)   \newline gleichzeitig bestehende andere  \newline Betreuungsarrangements, & \textcolor{diffred}{h) unverändert } &  \\
\hline
i)   \newline Monat und Jahr der Aufnahme in  \newline Kindertagespflege. & \textcolor{diffred}{i) unverändert } &  \\
\hline
(7b) Erhebungsmerkmale bei den  \newline Erhebungen über Personen, die mit  \newline öffentlichen Mitteln geförderte  \newline Kindertagespflege gemeinsam oder auf  \newline Grund einer Erlaubnis nach \S{} 43 Absatz 3  \newline Satz 3 durchführen und die von diesen  \newline betreuten Kinder sind die Zahl der  \newline Kindertagespflegepersonen und die Zahl  \newline der von diesen betreuten Kinder jeweils  \newline gegliedert nach Pflegestellen. & \textcolor{diffred}{(7b) unverändert } &  \\
\hline
(7c) Erhebungsmerkmale bei den  \newline Erhebungen über Kinder in den  \newline Klassenstufen eins bis vier sind & \textcolor{diffred}{(7c) unverändert } &  \\
\hline
1.   \newline Klassenstufe, & 1. unverändert &  \\
\hline
2.   \newline Anzahl der Wochenstunden, die das Kind in  \newline Angeboten nach \S{} 24 Absatz 4 verbringt, & 2. unverändert &  \\
\hline
3.   \newline Art der Angebote nach \S{} 24 Absatz 4. & 3. unverändert &  \\
\hline
(8) Erhebungsmerkmale bei den  \newline Erhebungen über die Angebote der  \newline Jugendarbeit nach \S{} 11 sowie bei den  \newline Erhebungen über Fortbildungsmaßnahmen  \newline für ehrenamtliche Mitarbeiter anerkannter  \newline Träger der Jugendhilfe nach \S{} 74 Absatz 6  \newline sind offene und Gruppenangebote sowie  \newline Veranstaltungen und Projekte der  \newline Jugendarbeit, soweit diese mit öffentlichen  \newline Mitteln pauschal oder maßnahmenbezogen  \newline gefördert werden oder der Träger eine  \newline öffentliche Förderung erhält, gegliedert  \newline nach & (8) unverändert &  \\
\hline
1.   \newline Art und Rechtsform des Trägers sowie bei  \newline Trägern der freien Jugendhilfe deren  \newline Verbandszugehörigkeit, & 1. unverändert &  \\
\hline
2.   \newline Dauer, Häufigkeit, Durchführungsort und  \newline Art des Angebots; zusätzlich bei  \newline schulbezogenen Angeboten die Art der  \newline kooperierenden Schule, & 2. unverändert &  \\
\hline
3.   \newline Art der Beschäftigung und Tätigkeit der bei  \newline der Durchführung des Angebots tätigen  \newline Personen sowie, mit Ausnahme der  \newline sonstigen pädagogisch tätigen Personen,  \newline deren Altersgruppe und Geschlecht, & 3. unverändert &  \\
\hline
4.   \newline Zahl der Teilnehmenden und der Besucher  \newline sowie, mit Ausnahme von Festen, Feiern,  \newline Konzerten, Sportveranstaltungen und  \newline sonstigen Veranstaltungen, deren  \newline Geschlecht und Altersgruppe, & 4. unverändert &  \\
\hline
5.   \newline Partnerländer und Veranstaltungen im In-  \newline oder Ausland bei Veranstaltungen und  \newline Projekten der internationalen Jugendarbeit. & 5. unverändert &  \\
\hline
(9) Erhebungsmerkmale bei den  \newline Erhebungen über die Träger der  \newline Jugendhilfe, die dort tätigen Personen und  \newline deren Einrichtungen, soweit diese nicht in  \newline Absatz 7 erfasst werden, sind & (9) unverändert &  \\
\hline
1.   \newline die Träger gegliedert nach & 1. unverändert &  \\
\hline
a)   \newline Art und Rechtsform des Trägers sowie bei  \newline Trägern der freien Jugendhilfe deren  \newline Verbandszugehörigkeit, & \textcolor{diffred}{a) unverändert } &  \\
\hline
b)   \newline den Betätigungsfeldern nach  \newline Aufgabenbereichen & \textcolor{diffred}{b) unverändert } &  \\
\hline
c)   \newline deren Personalausstattung sowie & \textcolor{diffred}{c) unverändert } &  \\
\hline
d)   \newline Anzahl der Einrichtungen, & \textcolor{diffred}{d) unverändert } &  \\
\hline
2. die Einrichtungen des Trägers mit  \newline Betriebserlaubnis nach \S{} 45 und  \newline Betreuungsformen nach diesem Gesetz,  \newline soweit diese nicht in Absatz 7 erfasst  \newline werden, gegliedert nach & 2. unverändert &  \\
\hline
a)   \newline Postleitzahl des Standorts, & \textcolor{diffred}{a) unverändert } &  \\
\hline
b)   \newline für jede vorhandene Gruppe und jede  \newline sonstige Betreuungsform nach diesem  \newline Gesetz, die von der Betriebserlaubnis  \newline umfasst ist, Angaben über die Art der  \newline Unterbringung oder Betreuung, deren  \newline Rechtsgrundlagen, Anzahl der  \newline genehmigten und belegten Plätze, Anzahl  \newline der Sollstellen des Personals und  \newline Hauptstelle der Einrichtung, & \textcolor{diffred}{b) unverändert } &  \\
\hline
3.   \newline für jede im Bereich der Jugendhilfe  \newline pädagogisch und in der Verwaltung tätige  \newline Person des Trägers & 3. unverändert &  \\
\hline
a)   \newline Geschlecht, Geburtsmonat und  \newline Geburtsjahr, & \textcolor{diffred}{a) unverändert } &  \\
\hline
b)   \newline Art des höchsten  \newline Berufsausbildungsabschlusses, Stellung im  \newline Beruf, Art der Beschäftigung,  \newline Beschäftigungsumfang und  \newline Arbeitsbereiche, & \textcolor{diffred}{b) unverändert } &  \\
\hline
c)   \newline Bundesland des überwiegenden  \newline Einsatzortes. & \textcolor{diffred}{c) unverändert } &  \\
\hline
(10) Erhebungsmerkmale bei der Erhebung  \newline der Ausgaben und Einnahmen der  \newline öffentlichen Jugendhilfe sind & (10) unverändert &  \\
\hline
1.   \newline die Art des Trägers, & 1. unverändert &  \\
\hline
2.   \newline die Ausgaben für Einzel- und  \newline Gruppenhilfen, gegliedert nach Ausgabe-  \newline und Hilfeart sowie die Einnahmen nach  \newline Einnahmeart, & \textcolor{diffred}{2.   \newline die Ausgaben für Einzel- und  \newline Gruppenhilfen }\textbf{\textcolor{diffred}{oder Einzel- und  \newline Gruppenleistungen}}\textcolor{diffred}{, gegliedert nach  \newline Ausgabe- und Hilfe}\textbf{\textcolor{diffred}{- oder}} \textbf{\textcolor{diffred}{Leistungs}}\textcolor{diffred}{art  \newline sowie die Einnahmen nach Einnahmeart, } &  \\
\hline
3.   \newline die Ausgaben und Einnahmen für  \newline Einrichtungen nach Arten gegliedert nach  \newline der Einrichtungsart, & 3.unverändert &  \\
\hline
4.   \newline die Ausgaben für das Personal, das bei  \newline den örtlichen und den überörtlichen  \newline Trägern sowie den kreisangehörigen  \newline Gemeinden und Gemeindeverbänden, die  \newline nicht örtliche Träger sind, Aufgaben der  \newline Jugendhilfe wahrnimmt. & 4.unverändert &  \\
\hline
\rowcolor{gray!25}
\textbf{\S{} 101} & \textbf{\textcolor{diffred}{\S{} 101 }} &  \\
\hline
\textbf{Periodizität und Berichtszeitraum} & \textbf{\textcolor{diffred}{Periodizität und Berichtszeitraum}} &  \\
\hline
(1) Die Erhebungen nach \S{} 99 Absatz 1 bis  \newline 5, 6a bis 7c und 10 sind jährlich  \newline durchzuführen, die Erhebungen nach \S{} 99  \newline Absatz 3 Nummer 3 erstmalig für das Jahr  \newline 2022; die Erhebungen nach \S{} 99 Absatz 1,  \newline soweit sie die Eingliederungshilfe für  \newline Kinder und Jugendliche mit seelischer  \newline Behinderung betreffen, sind 2007  \newline beginnend jährlich durchzuführen. Die  \newline Erhebung nach \S{} 99 Absatz 6 erfolgt  \newline laufend. Die übrigen Erhebungen nach \S{}  \newline 99 sind alle zwei Jahre durchzuführen, die  \newline Erhebungen nach \S{} 99 Absatz 8 erstmalig  \newline für das Jahr 2015 und die Erhebungen  \newline nach \S{} 99 Absatz 9 erstmalig für das Jahr  \newline 2014. & \textcolor{diffred}{(1) Die Erhebungen nach \S{} 99 Absatz 1 bis  \newline 5, 6a bis 7c und 10 sind jährlich  \newline durchzuführen, die Erhebungen nach \S{} 99  \newline Absatz 3 Nummer 3 erstmalig für das Jahr  \newline 2022; die Erhebungen nach \S{} 99 Absatz 1,  \newline soweit sie die Eingliederungshilfe für  \newline Kinder und Jugendliche mit seelischer  \newline Behinderung betreffen, sind 2007  \newline beginnend jährlich durchzuführen. Die  \newline Erhebung nach \S{} 99 Absatz 6 erfolgt  \newline laufend. Die übrigen Erhebungen nach \S{}  \newline 99 sind alle zwei Jahre durchzuführen, die  \newline Erhebungen nach \S{} 99 Absatz 8 erstmalig  \newline für das Jahr 2015 und die Erhebungen  \newline nach \S{} 99 Absatz 9 erstmalig für das Jahr  \newline 2014. }\textbf{\textcolor{diffred}{Die Erhebungen nach \S{} 99 Absatz  \newline 1, soweit sie die Eingliederungshilfe für  \newline Kinder und Jugendliche mit (drohenden)  \newline geistigen oder körperlichen  \newline Behinderungen betreffen, sind 2029  \newline beginnend jährlich durchzuführen.}} &  \\
\hline
(2) Die Angaben für die Erhebung nach & (2) unverändert &  \\
\hline
1.   \newline \S{} 99 Absatz 1 sind zu dem Zeitpunkt, zu  \newline dem die Hilfe endet, bei fortdauernder Hilfe  \newline zum 31. Dezember, & 1.unverändert &  \\
\hline
2.   \newline bis 5. (weggefallen) & 2.unverändert &  \\
\hline
6.   \newline \S{} 99 Absatz 2 sind zum Zeitpunkt des  \newline Endes einer vorläufigen Maßnahme, & 6.unverändert &  \\
\hline
7.   \newline \S{} 99 Absatz 3 Nummer 1 sind zum  \newline Zeitpunkt der rechtskräftigen gerichtlichen  \newline Entscheidung über die Annahme als Kind, & 7.unverändert &  \\
\hline
8.   \newline \S{} 99 Absatz 3 Nummer 2 Buchstabe a,  \newline Nummer 3 und Absatz 6a, 6b und 10 sind  \newline für das abgelaufene Kalenderjahr, & 8.unverändert &  \\
\hline
9.   \newline \S{} 99 Absatz 3 Nummer 2 Buchstabe b und  \newline Absatz 4 und 5 sind zum 31. Dezember, & 9.unverändert &  \\
\hline
10.   \newline \S{} 99 Absatz 7, 7a bis 7c sind zum 1. März, & 10.unverändert &  \\
\hline
11.   \newline \S{} 99 Absatz 6 sind zum Zeitpunkt des  \newline Abschlusses der  \newline Gefährdungseinschätzung, & 11.unverändert &  \\
\hline
12.   \newline \S{} 99 Absatz 8 sind für das abgelaufene  \newline Kalenderjahr, & 12.unverändert &  \\
\hline
13.   \newline \S{} 99 Absatz 9 sind zum 15. Dezember. zu  \newline erteilen. & 13.unverändert &  \\
\hline
\rowcolor{gray!25}
\textbf{\S{} 102} & \textbf{\textcolor{diffred}{\S{} 102 }} &  \\
\hline
\textbf{Auskunftspflicht } & \textbf{\textcolor{diffred}{Auskunftspflicht }} &  \\
\hline
(1) Für die Erhebungen besteht  \newline Auskunftspflicht. Die Angaben zu \S{} 100  \newline Nummer 4 sind freiwillig. & (1) unverändert &  \\
\hline
(2) Auskunftspflichtig sind & (2) unverändert &  \\
\hline
1.   \newline die örtlichen Träger der Jugendhilfe für die  \newline Erhebungen nach \S{} 99 Absatz 1 bis 10,  \newline nach Absatz 8 nur, soweit eigene Angebote  \newline gemacht wurden, & 1.unverändert &  \\
\hline
2.   \newline die überörtlichen Träger der Jugendhilfe für  \newline die Erhebungen nach \S{} 99 Absatz 3 und 7  \newline und 8 bis 10, nach Absatz 8 nur, soweit  \newline eigene Angebote gemacht wurden, & \textcolor{diffred}{2.   \newline die überörtlichen Träger der Jugendhilfe für  \newline die Erhebungen nach \S{} 99 Absatz }\textbf{\textcolor{diffred}{1,}}\textcolor{diffred}{ 3 und  \newline 7 und 8 bis 10, nach Absatz 8 nur, soweit  \newline eigene Angebote gemacht wurden, } &  \\
\hline
3.   \newline die obersten Landesjugendbehörden für  \newline die Erhebungen nach \S{} 99 Absatz 7 und 8  \newline bis 10, & 3.unverändert &  \\
\hline
4.   \newline die fachlich zuständige oberste  \newline Bundesbehörde für die Erhebung nach \S{}  \newline 99 Absatz 10, & 4.unverändert &  \\
\hline
5.   \newline die kreisangehörigen Gemeinden und die  \newline Gemeindeverbände, soweit sie Aufgaben  \newline der Jugendhilfe wahrnehmen, für die  \newline Erhebungen nach \S{} 99 Absatz 7 bis 10, & 5.unverändert &  \\
\hline
6.   \newline die Träger der freien Jugendhilfe für  \newline Erhebungen nach \S{} 99 Absatz 1, soweit sie  \newline eine Beratung nach \S{} 28 oder \S{} 41  \newline betreffen, nach \S{} 99 Absatz 8, soweit sie  \newline anerkannte Träger der freien Jugendhilfe  \newline nach \S{} 75 Absatz 1 oder Absatz 3 sind,  \newline und nach \S{} 99 Absatz 3, 7 und 9, & 6.unverändert &  \\
\hline
7.   \newline Adoptionsvermittlungsstellen nach \S{} 2  \newline Absatz 3 des   \newline Adoptionsvermittlungsgesetzes aufgrund  \newline ihrer Tätigkeit nach \S{} 1 des  \newline Adoptionsvermittlungsgesetzes sowie  \newline anerkannte Auslandsvermittlungsstellen  \newline nach \S{} 4 Absatz 2 Satz 3 des  \newline Adoptionsvermittlungsgesetzes aufgrund  \newline ihrer Tätigkeit nach \S{} 2a Absatz 4 Nummer  \newline 2 des Adoptionsvermittlungsgesetzes  \newline gemäß \S{} 99 Absatz 3 Nummer 1 sowie  \newline gemäß \S{} 99 Absatz 3 Nummer 2a für die  \newline Zahl der ausgesprochenen Annahmen und  \newline gemäß \S{} 99 Absatz 3 Nummer 2b für die  \newline Zahl der vorgemerkten Adoptionsbewerber, & 7.unverändert &  \\
\hline
8.   \newline die Leiter der Einrichtungen, Behörden und  \newline Geschäftsstellen in der Jugendhilfe für die  \newline Erhebungen nach \S{} 99 Absatz 7. Die  \newline Auskunftspflichtigen für Erhebungen nach  \newline \S{} 99 Absatz 7c werden durch Landesrecht  \newline bestimmt. & 8.unverändert &  \\
\hline
(3) Zur Durchführung der Erhebungen nach  \newline \S{} 99 Absatz 1, 3, 7, 8 und 9 übermitteln die  \newline Träger der öffentlichen Jugendhilfe den  \newline statistischen Ämtern der Länder auf  \newline Anforderung die erforderlichen Anschriften  \newline der übrigen Auskunftspflichtigen. & (3) unverändert &  \\
\hline
\rowcolor{gray!25}
\textbf{\S{} 104} &  & \textbf{\S{} 104} \\
\hline
Bußgeldvorschriften &  & Bußgeldvorschriften \\
\hline
(1) Ordnungswidrig handelt, wer   \newline 1.   ohne Erlaubnis nach \S{} 43 Absatz 1 oder  \newline \S{} 44 Absatz 1 Satz 1 ein Kind oder einen  \newline Jugendlichen betreut oder ihm Unterkunft  \newline gewährt,   \newline 2. entgegen \S{} 45 Absatz 1 Satz 1, auch in  \newline Verbindung mit \S{} 48a Absatz 1, ohne  \newline Erlaubnis eine Einrichtung oder eine  \newline sonstige Wohnform betreibt oder   \newline 3. entgegen \S{} 47 eine Anzeige nicht, nicht  \newline richtig, nicht vollständig oder nicht  \newline rechtzeitig erstattet oder eine Meldung nicht,  \newline nicht richtig, nicht vollständig oder nicht  \newline rechtzeitig macht oder vorsätzlich oder  \newline fahrlässig seiner Verpflichtung zur  \newline Dokumentation oder Aufbewahrung  derselben oder zum Nachweis der  \newline ordnungsgemäßen Buchführung auf  \newline entsprechendes Verlangen nicht  \newline nachkommt oder   \newline 4. entgegen \S{} 97a Absatz 4 vorsätzlich oder  \newline fahrlässig als Arbeitgeber eine Auskunft  \newline nicht, nicht richtig oder nicht vollständig  \newline erteilt. &  & (1) Ordnungswidrig handelt, wer   \newline 1.   ohne Erlaubnis nach \S{} 43 Absatz 1 oder  \newline \S{} 44 Absatz 1 Satz 1 ein Kind oder einen  \newline Jugendlichen betreut oder ihm Unterkunft  \newline gewährt,   \newline 2. entgegen \S{} 45 Absatz 1 Satz 1, auch in  \newline Verbindung mit \S{} 48a Absatz 1, ohne  \newline Erlaubnis eine Einrichtung oder eine  \newline sonstige Wohnform betreibt oder   \newline 3. entgegen \S{} 47 eine Anzeige nicht, nicht  \newline richtig, nicht vollständig oder nicht  \newline rechtzeitig erstattet oder eine Meldung nicht,  \newline nicht richtig, nicht vollständig oder nicht  \newline rechtzeitig macht oder vorsätzlich oder  \newline fahrlässig seiner Verpflichtung zur  \newline Dokumentation oder Aufbewahrung  derselben oder zum Nachweis der  \newline ordnungsgemäßen Buchführung auf  \newline entsprechendes Verlangen nicht  \newline nachkommt oder   \newline 4. entgegen \S{} 97a Absatz 4 vorsätzlich oder  \newline fahrlässig als Arbeitgeber eine Auskunft  \newline nicht, nicht richtig oder nicht vollständig  \newline \textcolor{diffgreen}{erteilt }\textbf{\textcolor{diffgreen}{oder  \newline 5. entgegen \S{} 42e vorsätzlich einen  \newline anderen gewöhnlichen Aufenthalt  \newline nimmt.}} \\
\hline
(2) Die Ordnungswidrigkeiten nach Absatz 1  \newline Nummer 1, 3 und 4 können mit einer  \newline Geldbuße bis zu fünfhundert Euro, die  \newline Ordnungswidrigkeit nach Absatz 1 Nummer  \newline 2 kann mit einer Geldbuße bis zu  \newline fünfzehntausend Euro geahndet werden. &  & (2) Die Ordnungswidrigkeiten nach Absatz 1  \newline Nummer 1, \textcolor{diffgreen}{3, }\textbf{\textcolor{diffgreen}{4}}\textbf{ und }\textbf{\textcolor{diffgreen}{5}} können mit einer  \newline Geldbuße bis zu fünfhundert Euro, die  \newline Ordnungswidrigkeit nach Absatz 1 Nummer  \newline 2 kann mit einer Geldbuße bis zu  \newline fünfzehntausend Euro geahndet werden. \\
\hline
... & ... & ... \\
\hline
\textbf{Übergangsregelung } & \textbf{Evaluation } & \textbf{Evaluation } \\
\hline
(1) Das Bundesministerium für Familie,  \newline Senioren, Frauen und Jugend begleitet und  \newline untersucht & (1) \textcolor{diffred}{entfällt} & \textbf{(1) }\textbf{\textcolor{diffgreen}{Das Bundesministerium für Familie,  \newline Senioren, Frauen und Jugend  \newline untersucht unter Beteiligung der Länder  \newline das Kinder- und  \newline Jugendstärkungsgesetz vom 3. Juni  \newline 2024 (BGBl. 9. Juni 2021, 29)  \newline einschließlich der Regelungen des  \newline Bundesgesetzes im Sinne des Artikel 1  \newline Nummer 12 \S{} 10 Absatz 3 Satz 3 des  \newline Kinder- und Jugendstärkungsgesetzes  \newline auf seine Wirkungen. Es wird  \newline untersucht, inwiefern die Regelungen  \newline das Ziel, gesellschaftliche Teilhabe und  \newline Chancengleichheit für alle jungen  \newline Menschen zu sichern oder herzustellen,  \newline erreicht werden konnte. Zudem wird  \newline untersucht, welche finanziellen  \newline Auswirkungen die Regelungen auf  \newline Länder und Kommunen haben. Als  \newline Kriterien für die Evaluation dienen  \newline insbesondere die Vollständigkeit der  \newline Umsetzung der Regelungen sowie die  \newline Inanspruchnahme von Leistungen auch  \newline unter Berücksichtigung der Perspektive  \newline der Normadressatinnen und - \newline adressaten. Als Grundlage dienen die }} \textbf{\textcolor{diffgreen}{Daten der Kinder- und  \newline Jugendhilfestatistik. Das  \newline Bundesministerium für Familien,  \newline Senioren, Frauen und Jugend berichtet  \newline dem Deutschen Bundestag und dem  \newline Bundesrat über die Ergebnisse dieser  \newline Untersuchung}}\textcolor{diffgreen}{.} \\
\hline
... & ... & ... \\
\hline
(2) Das Bundesministerium für Familie,  \newline Senioren, Frauen und Jugend untersucht in  \newline den Jahren 2022 bis 2024 die rechtlichen  \newline Wirkungen von \S{} 10 Absatz 4 und legt dem  \newline Bundestag und dem Bundesrat bis zum 31.  \newline Dezember 2024 einen Bericht über das  \newline Ergebnis der Untersuchung vor. Dabei  \newline sollen insbesondere die gesetzlichen  \newline Festlegungen des Achten und Neunten  \newline Buches & (2) \textcolor{diffred}{entfällt} & \textbf{(2) }\textbf{\textcolor{diffgreen}{Es wird ein Konzept zur künftigen  \newline inhaltlichen Ausgestaltung der Kinder-  \newline und Jugendhilfestatistik entwickelt, auf  \newline dessen Grundlage in einem  \newline Bundesgesetz die zur Beurteilung der  \newline Auswirkungen der Bestimmungen  \newline dieses Buches und zu seiner  \newline Fortentwicklung notwendigen laufenden  \newline Erhebungen auch im Hinblick auf die ab  \newline dem 1. Januar 2028 vorrangige  \newline Zuständigkeit der Kinder- und  \newline Jugendhilfe für Leistungen der  \newline Eingliederungshilfe für junge Menschen  \newline mit körperlichen oder geistigen  \newline Behinderungen oder hiervon bedrohte  \newline junge Menschen geregelt und 2030  \newline beginnend durchgeführt werden  \newline können.}} \\
\hline
... & ... & ... \\
\hline
(3) Soweit das Bundesministerium für  \newline Familie, Senioren, Frauen und Jugend  \newline Dritte in die Durchführung der  \newline Untersuchungen nach den Absätzen 1 und  \newline 2 einbezieht, beteiligt es hierzu vorab die  \newline Länder. & (3) entfällt  \textcolor{diffred}{Das Bundesministerium für Familie,  \newline Senioren, Frauen und Jugend untersucht  \newline unter Beteiligung der Länder die Wirkungen} \newline \textbf{\textcolor{diffred}{des Kinder- und  \newline Jugendstärkungsgesetzes vom 3. Juni  \newline 2024 (BGBl. 9. Juni 2021, 29)}}\textcolor{diffred}{ im Übrigen  \newline einschließlich seiner finanziellen  \newline Auswirkungen auf Länder und Kommunen  \newline und berichtet dem Deutschen Bundestag  \newline und dem Bundesrat über die Ergebnisse  \newline dieser Untersuchung.} & (3) entfällt \\
\hline
... & ... & ... \\
\hline
 & \textbf{(1) Die Leistungs- und  \newline Vergütungsvereinbarungen nach Kapitel  \newline 8 des Teils 2 des Neunten Buches  \newline gelten für die in \S{} 78a benannten  \newline Leistungen als Vereinbarungen nach \S{}  \newline 78b und bei ambulanten Leistungen als  \newline Vereinbarungen nach \S{} 77 }\textbf{\textcolor{diffred}{SGB VIII}}\textbf{ bis  \newline zum 31. Dezember 2032 fort. Die  \newline Vereinbarungen umfassen die  \newline Leistungen für minderjährige  \newline Leistungsberechtigte, auf die sich die  \newline Leistungs- und  \newline Vergütungsvereinbarungen im Sinne  \newline des Satz 1 bisher bezogen haben, sowie  \newline Leistungen nach \S{} 41, die inhaltlich den  \newline bisher vereinbarten Leistungen  \newline entsprechen. } & \textbf{(1) Die Leistungs- und  \newline Vergütungsvereinbarungen nach Kapitel  \newline 8 des Teils 2 des Neunten Buches  \newline gelten für die in \S{} 78a benannten  \newline Leistungen als Vereinbarungen nach \S{}  \newline 78b und bei ambulanten Leistungen als  \newline Vereinbarungen nach \S{} 77 bis zum 31.  \newline Dezember 2032 fort. Die} \newline \textbf{\textcolor{diffgreen}{Vereinbarungen, die als}}\textbf{ Vereinbarungen} \newline \textbf{\textcolor{diffgreen}{nach \S{} 78b oder \S{} 77 fortgelten,}}\textbf{  \newline umfassen die Leistungen für  \newline minderjährige Leistungsberechtigte, auf  \newline die sich die Leistungs- und  \newline Vergütungsvereinbarungen im Sinne  \newline des Satz}\textbf{\textcolor{diffgreen}{es}}\textbf{ 1 bisher bezogen haben,  \newline sowie Leistungen nach \S{} 41, die  \newline inhaltlich den bisher vereinbarten  \newline Leistungen entsprechen. } \\
\hline
... & ... & ... \\
\hline
 & \textbf{(3) Die Leistungsbescheide auf  \newline Grundlage des \S{} 99 des Neunten  \newline Buches gelten als Bescheide nach \S{} 27  \newline Absatz 3 fort. } & \textbf{(3) Die Leistungsbescheide }\textbf{\textcolor{diffgreen}{für  \newline minderjährige Leistungsberechtigte}}\textbf{ auf  \newline Grundlage des \S{} 99 des Neunten  \newline Buches gelten als Bescheide nach \S{} 27  \newline Absatz 3 fort. } \\
\hline
 & \textbf{(4) Jede Vertragspartei der Leistungs-  \newline und Vergütungsvereinbarungen nach  \newline Ab}\textbf{\textcolor{diffred}{-}}\textbf{satz 1 hat zum 1. Januar 2028  \newline unbeschadet der Laufzeit der nach  \newline Absatz 1 und 2 fort}\textbf{\textcolor{diffred}{-}}\textbf{geltenden Verträge  \newline und Leistungsbestandteile einen  \newline Anspruch auf Neuverhandlung einer  \newline Vereinbarung nach \S{} 78b SGB VIII. } & \textbf{(4) Jede Vertragspartei der Leistungs-  \newline und Vergütungsvereinbarungen nach  \newline Absatz 1 hat zum 1. Januar 2028  \newline unbeschadet der Laufzeit der nach  \newline Absatz 1 und 2 fortgeltenden Verträge  \newline und Leistungsbestandteile einen  \newline Anspruch auf Neuverhandlung einer  \newline Vereinbarung nach \S{} 78b SGB VIII.} \textbf{\textcolor{diffgreen}{Die  \newline Frist nach \S{} 78g Absatz 2 Satz1 wird im  \newline Falle von Neuverhandlungen nach Satz  \newline 1 um 12 Wochen verlängert.}} \\
\hline
 & \textbf{(5) Die Bescheide zur Festsetzung des  \newline Beitrags aus Einkommen zu den  \newline Aufwendungen nach \S{} 136 des Neunten  \newline Buches sowie die Festsetzung der  \newline Aufbringung der Mittel für die Kosten  \newline des Lebensunterhaltes nach \S{} 142 des  \newline Neunten Buches gelten bis zu ihrer  \newline Aufhebung fort, sofern die Leistung  \newline nach \S{} 99 des Neunten Buches  \newline entsprechend des }\textbf{\textcolor{diffred}{Absatz 2}}\textbf{ ab dem 1.  \newline Januar 2028 als Leistung nach \S{} 27  \newline Absatz 3 fortgesetzt wird. }\textbf{\textcolor{diffred}{Für den  \newline Träger}}\textbf{ der Eingliederungshilfe tritt der  \newline zuständige Träger der öffentlichen} \newline \textbf{\textcolor{diffred}{Jugendhilfe ein.}}\textbf{ Die Aufhebung der  \newline Bescheide muss rückwirkend zum 1.  \newline Januar 2028 erfolgen; die Aufhebung  \newline s}\textbf{\textcolor{diffred}{oll}}\textbf{ spätestens am 31. Dezember 2028} \newline \textbf{\textcolor{diffred}{bei}}\textbf{ den durch die Bescheide  \newline Verpflichteten }\textbf{\textcolor{diffred}{ein}}\textbf{gehen}. & \textbf{(5) Die Bescheide zur Festsetzung des  \newline Beitrags aus Einkommen zu den } \textbf{Aufwendungen nach \S{} 136 des Neunten  \newline Buches sowie die Festsetzung der  \newline Aufbringung der Mittel für die Kosten  \newline des Lebensunterhaltes nach \S{} 142 des  \newline Neunten Buches gelten bis zu ihrer  \newline Aufhebung fort, sofern die Leistung  \newline nach \S{} 99 des Neunten Buches  \newline entsprechend des }\textbf{\textcolor{diffgreen}{Absatzes 3}}\textbf{ ab dem 1.  \newline Januar 2028 als Leistung nach \S{} 27  \newline Absatz 3 fortgesetzt wird. }\textbf{\textcolor{diffgreen}{An die Stelle  \newline des Trägers}}\textbf{ der Eingliederungshilfe tritt  \newline der zuständige Träger der öffentlichen} \newline \textbf{\textcolor{diffgreen}{Jugendhilfe.}}\textbf{ Die Aufhebung der  \newline Bescheide muss rückwirkend zum 1.  \newline Januar 2028 erfolgen; die Aufhebung} \newline \textbf{\textcolor{diffgreen}{mu}}\textbf{s}\textbf{\textcolor{diffgreen}{s}}\textbf{ spätestens am 31. Dezember 2028  \newline den durch die Bescheide Verpflichteten} \newline \textbf{\textcolor{diffgreen}{zu}}\textbf{gehen. } \\
\hline
... & ... & ... \\
\hline
 & \textbf{1.   \newline Wurde mindestens von einem Elternteil  \newline des Leistungsberechtigten der Einsatz  \newline des Einkommens nach \S{} 136 des  \newline Neunten Buches oder die Aufbringung  \newline der Mittel für die Kosten des  \newline Lebensunterhalts in Höhe der für den  \newline häuslichen Lebensunterhalt ersparten  \newline Aufwendungen nach \S{} 142 des Neunten  \newline Buches gefordert und ist der nach den  \newline \S{}\S{} 91 bis 94 aufzubringende Be}\textbf{\textcolor{diffred}{i}}\textbf{trag  \newline höher als der Einkommenseinsatz oder  \newline die }\textbf{\textcolor{diffred}{Höhe der}}\textbf{ aufzubringenden ersparten  \newline Aufwendungen } & \textbf{1.   \newline Wurde mindestens von einem Elternteil  \newline des Leistungsberechtigten der Einsatz  \newline des Einkommens nach \S{} 136 des  \newline Neunten Buches oder die Aufbringung  \newline der Mittel für die Kosten des  \newline Lebensunterhalts in Höhe der für den  \newline häuslichen Lebensunterhalt ersparten  \newline Aufwendungen nach \S{} 142 des Neunten  \newline Buches gefordert und ist der nach den  \newline \S{}\S{} 91 bis 94 aufzubringende Betrag  \newline höher als der Einkommenseinsatz oder} \newline \textbf{\textcolor{diffgreen}{als}}\textbf{ die aufzubringenden ersparten  \newline Aufwendungen }\textbf{\textcolor{diffgreen}{nach Kapitel 9 des Teils  \newline 2 des Neunten Buches mit Gültigkeit  \newline vom 31. Dezember 2027, so ist der  \newline Kostenbeitrag nach den \S{}\S{} 91 bis 94 auf  \newline diesen Betrag begrenzt. Die Begrenzung  \newline gilt auch dann, wenn bis zum 31.  \newline Dezember 2027 nur ein Elternteil zu den  \newline Kosten herangezogen wurde und nach  \newline den \S{}\S{} 91 bis 94 beide Elternteile  \newline getrennt zu den Kosten herangezogen  \newline werden. Der bisher von einem Elternteil  \newline aufgebrachte Betrag gilt dann als  \newline Höchstbetrag für die Summe der  \newline Kostenbeiträge beider Elternteile nach }} \textbf{\textcolor{diffgreen}{den \S{}\S{} 91 bis 94. Die Sätze 1 bis 3 gelten  \newline für die Kostenbeiträge von Elternteilen  \newline entsprechend, die auf der Grundlage der  \newline \S{}\S{} 91 bis 94 mit Gültigkeit vom 31.  \newline Dezember 2027 zu den Kosten  \newline herangezogen wurden.}} \\
\hline
 &  & \textbf{\textcolor{diffgreen}{2.   \newline Wurde bisher von keinem Elternteil der  \newline Einsatz des Einkommens nach \S{} 136  \newline des Neunten Buches oder die  \newline Aufbringung der Mittel für die Kosten  \newline des Lebensunterhalts nach \S{} 142 des  \newline Neunten Buches gefordert und wird die  \newline bis zum 31. Dezember 2027 erbrachte  \newline Leistung nach \S{} 99 des Neunten Buchs  \newline auf der Grundlage des \S{} 27 Absatz 3  \newline fortgesetzt oder neu bewilligt, so wird  \newline für diese Leistung kein Kostenbeitrag  \newline nach den \S{}\S{} 91 bis 94 erhoben.}} \\
\hline
- Aus Synopsis 2024 - \newline \textbf{Neunten Buch Sozialgesetzbuch  } \newline - Aus Synopsis 2026 - \newline \textbf{Neuntes Buch Sozialgesetzbuch  } & \textbf{\textcolor{diffred}{2.   \newline Wurde bisher}}\textbf{ von }\textbf{\textcolor{diffred}{keinem Elternteil}}\textbf{ der} \newline \textbf{\textcolor{diffred}{Einsatz des Einkommens}}\textbf{ nach }\textbf{\textcolor{diffred}{\S{} 136  \newline des Neunten Buches oder die  \newline Aufbringung}}\textbf{ der }\textbf{\textcolor{diffred}{Mittel für die Kosten  \newline des Lebensunter-halts nach \S{} 142 des  \newline Neunten Buches gefordert und wird die}}\textbf{  \newline bis zum 31. Dezember 2027 }\textbf{\textcolor{diffred}{erbrachte  \newline Leistung nach \S{} 99 des Neunten Buchs  \newline auf der Grundlage des \S{} 27 Absatz 3  \newline fortgesetzt oder neu bewilligt, so wird  \newline für diese Leistung kein Kostenbeitrag  \newline nach den \S{}\S{} 91}}\textbf{ bis }\textbf{\textcolor{diffred}{94 erhoben.}}  \textbf{Neuntes Buch Sozialgesetzbuch } & \textbf{\textcolor{diffgreen}{(7) Für Leistungen auf der Grundlage}}\textbf{  \newline von }\textbf{\textcolor{diffgreen}{Bescheiden nach Absatz 3 richtet  \newline sich die örtliche Zuständigkeit nach \S{}\S{}  \newline 86, 86c, 86d und 88. Spätestens bis zum  \newline 31. Oktober 2028 ist die örtliche  \newline Zuständigkeit für Leistungen nach Satz  \newline 1 zu prüfen. Fand die Übergabe eines  \newline Falles}}\textbf{ der }\textbf{\textcolor{diffgreen}{Gewährung von Leistungen}}\textbf{  \newline nach }\textbf{\textcolor{diffgreen}{Satz 1 an den nach Satz 1  \newline zuständigen Träger}}\textbf{ der }\textbf{\textcolor{diffgreen}{Jugendhilfe  \newline nicht zum 1. Januar 2028 statt, sind dem  \newline Träger, der}}\textbf{ bis zum 31. Dezember 2027} \newline \textbf{\textcolor{diffgreen}{zuständig war, die Kosten}}\textbf{ bis }\textbf{\textcolor{diffgreen}{zur  \newline tatsächlichen Übergabe zu erstatten.}}  \textbf{Neuntes Buch Sozialgesetzbuch } \\
\hline
- Aus Synopsis 2024 - \newline \textbf{( - SGB IX)  \newline vom: 23.12.2016 - zuletzt geändert durch  \newline Art. 6 v. 22.12.2023 I Nr. 412 } \newline - Aus Synopsis 2026 - \newline \textbf{( - SGB IX)  \newline vom: 23.12.2016 - zuletzt geändert durch  \newline Art.13 v. 16.1.2026 I Nr. 14 } & \textbf{( - SGB IX)  \newline vom: 23.12.2016 - zuletzt geändert durch} \newline \textbf{\textcolor{diffred}{Art. 6}}\textbf{ v. }\textbf{\textcolor{diffred}{22}}\textbf{.1}\textbf{\textcolor{diffred}{2}}\textbf{.202}\textbf{\textcolor{diffred}{3}}\textbf{ I Nr. 4}\textbf{\textcolor{diffred}{12}} & \textbf{( - SGB IX)  \newline vom: 23.12.2016 - zuletzt geändert durch} \newline \textbf{\textcolor{diffgreen}{Art.13}}\textbf{ v. }\textbf{\textcolor{diffgreen}{16}}\textbf{.1.202}\textbf{\textcolor{diffgreen}{6}}\textbf{ I Nr. }\textbf{\textcolor{diffgreen}{1}}\textbf{4 } \\
\hline
... & ... & ... \\
\hline
\rowcolor{gray!25}
\textbf{\S{} 51} & \textbf{\textcolor{diffred}{\S{} 51 }} &  \\
\hline
\textbf{Rechtsweg und Zuständigkeit } & \textbf{\textcolor{diffred}{Rechtsweg und Zuständigkeit }} &  \\
\hline
(1) Die Gerichte der Sozialgerichtsbarkeit  \newline entscheiden über öffentlich-rechtliche  \newline Streitigkeiten & (1) unverändert &  \\
\hline
1.   \newline in Angelegenheiten der gesetzlichen  \newline Rentenversicherung einschließlich der  \newline Alterssicherung der Landwirte, & 1. unverändert &  \\
\hline
2.  \newline in Angelegenheiten der gesetzlichen  \newline Krankenversicherung, der sozialen  \newline Pflegeversicherung und der privaten  \newline Pflegeversicherung (Elftes Buch  \newline Sozialgesetzbuch), auch soweit durch  \newline diese Angelegenheiten Dritte betroffen  \newline werden; dies gilt nicht für Streitigkeiten in  \newline Angelegenheiten nach \S{} 110 des Fünften  \newline Buches Sozialgesetzbuch aufgrund einer  \newline Kündigung von Versorgungsverträgen, die  \newline für Hochschulkliniken oder  \newline Plankrankenhäuser (\S{} 108 Nr. 1 und 2 des  \newline Fünften Buches Sozialgesetzbuch) gelten, & 2. unverändert &  \\
\hline
3.  \newline in Angelegenheiten der gesetzlichen  \newline Unfallversicherung mit Ausnahme der  \newline Streitigkeiten aufgrund der Überwachung  \newline der Maßnahmen zur Prävention durch die  \newline Träger der gesetzlichen  \newline Unfallversicherung, & 3. unverändert &  \\
\hline
4.  \newline in Angelegenheiten der Arbeitsförderung  \newline einschließlich der übrigen Aufgaben der  \newline Bundesagentur für Arbeit, & 4. unverändert &  \\
\hline
5.  \newline in sonstigen Angelegenheiten der  \newline Sozialversicherung, & 5. unverändert &  \\
\hline
6.  \newline in Angelegenheiten des Sozialen  \newline Entschädigungsrechts, & 6. unverändert &  \\
\hline
6a.  \newline in Angelegenheiten der Sozialhilfe  \newline einschließlich der Angelegenheiten nach  \newline Teil 2 des Neunten Buches  \newline Sozialgesetzbuch und des  \newline Asylbewerberleistungsgesetzes, & \textcolor{diffred}{6a. unverändert } &  \\
\hline
 & \textbf{\textcolor{diffred}{6b.  \newline In Angelegenheiten der Kinder- und  \newline Jugendhilfe, soweit sie Leistungen der  \newline Eingliederungshilfe betreffen, }} &  \\
\hline
7.  \newline bei der Feststellung von Behinderungen  \newline und ihrem Grad sowie weiterer  \newline gesundheitlicher Merkmale, ferner der  \newline Ausstellung, Verlängerung, Berichtigung  \newline und Einziehung von Ausweisen nach \S{} 152  \newline des Neunten Buches Sozialgesetzbuch, & 7. unverändert &  \\
\hline
8.  \newline die aufgrund des  \newline Aufwendungsausgleichsgesetzes  \newline entstehen, & 8. unverändert &  \\
\hline
9.  \newline (weggefallen) & 9. unverändert &  \\
\hline
10.  \newline für die durch Gesetz der Rechtsweg vor  \newline diesen Gerichten eröffnet wird. & 10. unverändert &  \\
\hline
(2) Die Gerichte der Sozialgerichtsbarkeit  \newline entscheiden auch über privatrechtliche  \newline Streitigkeiten in Angelegenheiten der  \newline Zulassung von Trägern und Maßnahmen  \newline durch fachkundige Stellen nach dem  \newline Fünften Kapitel des Dritten Buches  \newline Sozialgesetzbuch und in Angelegenheiten  \newline der gesetzlichen Krankenversicherung,  \newline auch soweit durch diese Angelegenheiten  \newline Dritte betroffen werden. Satz 1 gilt für die  \newline soziale Pflegeversicherung und die private  \newline Pflegeversicherung (Elftes Buch  \newline Sozialgesetzbuch) entsprechend. & (2) unverändert &  \\
\hline
(3) Von der Zuständigkeit der Gerichte der  \newline Sozialgerichtsbarkeit nach den Absätzen 1  \newline und 2 ausgenommen sind Streitigkeiten in  \newline Verfahren nach dem Gesetz gegen  \newline Wettbewerbsbeschränkungen, die  \newline Rechtsbeziehungen nach \S{} 69 des Fünften  \newline Buches Sozialgesetzbuch betreffen  \textbf{Erstes Buch Sozialgesetzbuch } & \textcolor{diffred}{(3) unverändert} \newline \textbf{\textcolor{diffred}{Erstes Buch Sozialgesetzbuch }} &  \\
\hline
... & ... & ... \\
\hline
\rowcolor{gray!25}
\textbf{\S{} 85} &  & \textbf{\S{} 85} \\
\hline
Klagerecht der Verbände &  & Klagerecht der Verbände \\
\hline
Werden Menschen mit Behinderungen in  \newline ihren Rechten nach diesem Buch verletzt,  \newline können an ihrer Stelle und mit ihrem  \newline Einverständnis Verbände klagen, die nach  \newline ihrer Satzung Menschen mit  \newline Behinderungen auf Bundes- oder  \newline Landesebene vertreten und nicht selbst am  \newline Prozess beteiligt sind. In diesem Fall  \newline müssen alle Verfahrensvoraussetzungen  \newline wie bei einem Rechtsschutzersuchen durch  \newline den Menschen mit Behinderungen selbst  \newline vorliegen. &  & Werden Menschen mit Behinderungen in  \newline ihren Rechten nach diesem Buch \textbf{\textcolor{diffgreen}{oder  \newline nach dem zweiten Kapitel, Vierter  \newline Abschnitt, Dritter Unterabschnitt des  \newline Achten Buches}} verletzt, können an ihrer  \newline Stelle und mit ihrem Einverständnis  \newline Verbände klagen, die nach ihrer Satzung  \newline Menschen mit Behinderungen auf Bundes-  \newline oder Landesebene vertreten und nicht  \newline selbst am Prozess beteiligt sind. In diesem  \newline Fall müssen alle  \newline Verfahrensvoraussetzungen wie bei einem  \newline Rechtsschutzersuchen durch den  \newline Menschen mit Behinderungen selbst  \newline vorliegen. \\
\hline
... & ... & ... \\
\hline
- Aus Synopsis 2024 - \newline (5) Bei Personen, die am 31. Dezember  \newline 2019 Leistungen nach dem Sechsten  \newline Kapitel des Zwölften Buches in der am 31.  \newline Dezember 2019 geltenden Fassung  \newline bezogen haben und auch ab dem 1.  \newline Januar 2020 Leistungen nach Teil 2 dieses  \newline Buches erhalten, ist der Träger der  \newline Eingliederungshilfe örtlich zuständig,  \newline dessen örtliche Zuständigkeit sich am 1.  \newline Januar 2020 im Einzelfall in  \newline entsprechender Anwendung von \S{} 98  \newline Absatz 1 Satz 1 oder Absatz 5 des  \newline Zwölften Buches oder in entsprechender  \newline Anwendung von \S{} 98 Absatz 2 Satz 1 und  \newline 2 in Verbindung mit \S{} 107 des Zwölften  \newline Buches ergeben würde. Absatz 1 Satz 3  \newline bis 5 gilt entsprechend. Im Übrigen bleiben  \newline die Absätze 2 bis 4 unberührt.  \textbf{Sozialgerichtsgesetz } \newline - Aus Synopsis 2026 - \newline (5) Bei Personen, die am 31. Dezember  \newline 2019 Leistungen nach dem Sechsten  \newline Kapitel des Zwölften Buches in der am 31.  \newline Dezember 2019 geltenden Fassung  \newline bezogen haben und auch ab dem 1.  \newline Januar 2020 Leistungen nach Teil 2 dieses  \newline Buches erhalten, ist der Träger der  \newline Eingliederungshilfe örtlich zuständig,  \newline dessen örtliche Zuständigkeit sich am 1.  \newline Januar 2020 im Einzelfall in  \newline entsprechender Anwendung von \S{} 98  \newline Absatz 1 Satz 1 oder Absatz 5 des  \newline Zwölften Buches oder in entsprechender  \newline Anwendung von \S{} 98 Absatz 2 Satz 1 und  \newline 2 in Verbindung mit \S{} 107 des Zwölften  \newline Buches ergeben würde. Absatz 1 Satz 3  \newline bis 5 gilt entsprechend. Im Übrigen bleiben  \newline die Absätze 2 bis 4 unberührt. & \textcolor{diffred}{(5) unverändert  }\textbf{\textcolor{diffred}{Sozialgerichtsgesetz }} & (5) unverändert \\
\hline
\textbf{( - SGG)  \newline vom: 23.09.1975 - zuletzt geändert durch  \newline Art. 3 v. 22.03.2024 I Nr. 105 } & \textbf{\textcolor{diffred}{( - SGG)  \newline vom: 23.09.1975 - zuletzt geändert durch  \newline Art. 3 v. 22.03.2024 I Nr. 105 }} &  \\
\hline
\rowcolor{gray!25}
\textbf{\S{} 117} &  & \textbf{\S{} 117} \\
\hline
Gesamtplanverfahren &  & Gesamtplanverfahren \\
\hline
Absätze 1 bis 5 &  & unverändert \\
\hline
(6) Bei minderjährigen  \newline Leistungsberechtigten wird der nach \S{} 86  \newline des Achten Buches zuständige örtliche  \newline Träger der öffentlichen Jugendhilfe vom  \newline Träger der Eingliederungshilfe mit  \newline Zustimmung des  \newline Personensorgeberechtigten informiert und  \newline nimmt am Gesamtplanverfahren beratend  \newline teil, soweit dies zur Feststellung der  \newline Leistungen der Eingliederungshilfe nach  \newline den Kapiteln 3 bis 6 erforderlich ist.  \newline Hiervon kann in begründeten  \newline Ausnahmefällen abgesehen werden,  \newline insbesondere, wenn durch die Teilnahme  \newline des zuständigen örtlichen Trägers der  \newline öffentlichen Jugendhilfe das  \newline Gesamtplanverfahren verzögert würde. &  & (6) \textcolor{diffgreen}{entfällt} \\
\hline
\rowcolor{gray!25}
\textbf{\S{} 119} &  & \textbf{\S{} 119} \\
\hline
Gesamtplankonferenz &  & Gesamtplankonferenz \\
\hline
(1) Mit Zustimmung des  \newline Leistungsberechtigten kann der Träger der  \newline Eingliederungshilfe eine  \newline Gesamtplankonferenz durchführen, um die  \newline Leistungen für den Leistungsberechtigten  \newline nach den Kapiteln 3 bis 6 sicherzustellen.  \newline Die Leistungsberechtigten, die beteiligten  \newline Rehabilitationsträger und bei  \newline minderjährigen Leistungsberechtigten der  \newline nach \S{} 86 des Achten Buches zuständige  \newline örtliche Träger der öffentlichen Jugendhilfe  \newline können dem nach \S{} 15 verantwortlichen  \newline Träger der Eingliederungshilfe die  \newline Durchführung einer Gesamtplankonferenz  \newline vorschlagen. Den Vorschlag auf  \newline Durchführung einer Gesamtplankonferenz  \newline kann der Träger der Eingliederungshilfe  \newline ablehnen, wenn der maßgebliche  \newline Sachverhalt schriftlich ermittelt werden  \newline kann oder der Aufwand zur Durchführung  \newline nicht in einem angemessenen Verhältnis  \newline zum Umfang der beantragten Leistung  \newline steht. &  & (1) Mit Zustimmung des  \newline Leistungsberechtigten kann der Träger der  \newline Eingliederungshilfe eine  \newline Gesamtplankonferenz durchführen, um die  \newline Leistungen für den Leistungsberechtigten  \newline nach den Kapiteln 3 bis 6 sicherzustellen.  \newline Die \textcolor{diffgreen}{Leistungsberechtigten }\textbf{\textcolor{diffgreen}{und}}\textbf{ die  \newline beteiligten Rehabilitationsträger können  \newline dem nach \S{} 15 verantwortlichen Träger  \newline der Eingliederungshilfe die  \newline Durchführung einer  \newline Gesamtplankonferenz vorschlagen. }Den  \newline Vorschlag auf Durchführung einer  \newline Gesamtplankonferenz kann der Träger der  \newline Eingliederungshilfe ablehnen, wenn der  \newline maßgebliche Sachverhalt schriftlich  \newline ermittelt werden kann oder der Aufwand  \newline zur Durchführung nicht in einem  \newline angemessenen Verhältnis zum Umfang der  \newline beantragten Leistung steht. \\
\hline
Absätze 2 bis 4 &  & unverändert \\
\hline
\rowcolor{gray!25}
\textbf{\S{} 134} &  & \textbf{\S{} 134} \\
\hline
Sonderregelung zum Inhalt der  \newline Vereinbarungen zur Erbringung von  \newline Leistungen für minderjährige  \newline Leistungsberechtigte und in Sonderfällen &  & Sonderregelung zum Inhalt der  \newline Vereinbarungen zur Erbringung von  \newline Leistungen für minderjährige  \newline Leistungsberechtigte und in Sonderfällen \\
\hline
(1) In der schriftlichen Vereinbarung zur  \newline Erbringung von Leistungen für  \newline minderjährige Leistungsberechtigte  \newline zwischen dem Träger der  \newline Eingliederungshilfe und dem  \newline Leistungserbringer sind zu regeln:  \newline 1.  \newline Inhalt, Umfang und Qualität einschließlich  \newline der Wirksamkeit der Leistungen  \newline (Leistungsvereinbarung) sowie  \newline 2.  \newline die Vergütung der Leistung  \newline (Vergütungsvereinbarung). &  & (1) \textcolor{diffgreen}{entfällt} \\
\hline
(2) In die Leistungsvereinbarung sind als  \newline wesentliche Leistungsmerkmale  \newline insbesondere aufzunehmen:  \newline 1.  \newline die betriebsnotwendigen Anlagen des  \newline Leistungserbringers,  \newline 2.  \newline der zu betreuende Personenkreis,  \newline 3.  \newline Art, Ziel und Qualität der Leistung,  \newline 4.  \newline die Festlegung der personellen  \newline Ausstattung,  \newline 5.  \newline die Qualifikation des Personals sowie  \newline 6.  \newline die erforderliche sächliche Ausstattung. &  & (2) \textcolor{diffgreen}{entfällt} \\
\hline
(3) Die Vergütungsvereinbarung besteht  \newline mindestens aus  \newline 1.  \newline der Grundpauschale für Unterkunft und  \newline Verpflegung,  \newline 2.  \newline der Maßnahmepauschale sowie  \newline 3.  \newline einem Betrag für betriebsnotwendige  \newline Anlagen einschließlich ihrer Ausstattung  \newline (Investitionsbetrag).  \newline Förderungen aus öffentlichen Mitteln sind  \newline anzurechnen. Die Maßnahmepauschale ist  \newline nach Gruppen für Leistungsberechtigte mit  \newline vergleichbarem Bedarf zu kalkulieren. &  & (3) \textcolor{diffgreen}{entfällt} \\
\hline
(4) Die Absätze 1 bis 3 finden auch  \newline Anwendung, wenn volljährige  \newline Leistungsberechtigte Leistungen zur  \newline Schulbildung nach \S{} 112 Absatz 1 Nummer  \newline 1 sowie Leistungen zur schulischen  \newline Ausbildung für einen Beruf nach \S{} 112  \newline Absatz 1 Nummer 2 erhalten, soweit diese  \newline Leistungen in besonderen  \newline Ausbildungsstätten über Tag und Nacht für  \newline Menschen mit Behinderungen erbracht  \newline werden. Entsprechendes gilt bei anderen  \newline volljährigen Leistungsberechtigten, wenn  \newline 1.  \newline das Konzept des Leistungserbringers auf  \newline Minderjährige als zu betreuenden  \newline Personenkreis ausgerichtet ist,  \newline 2.  \newline der Leistungsberechtigte von diesem  \newline Leistungserbringer bereits Leistungen über  \newline Tag und Nacht auf Grundlage von  \newline Vereinbarungen nach den Absätzen 1 bis  \newline 3, \S{} 78b des Achten Buches, \S{} 75 Absatz 3  \newline des Zwölften Buches in der am 31.  \newline Dezember 2019 geltenden Fassung oder  \newline nach Maßgabe des \S{} 75 Absatz 4 des  \newline Zwölften Buches in der am 31. Dezember  \newline 2019 geltenden Fassung erhalten hat und  \newline 3.  \newline der Leistungsberechtigte nach Erreichen  \newline der Volljährigkeit für eine kurze Zeit, in der  \newline Regel nicht länger als bis zur Vollendung  \newline des 21. Lebensjahres, Leistungen von  \newline diesem Leistungserbringer weitererhält, mit  \newline denen insbesondere vor dem Erreichen der  \newline Volljährigkeit definierte Teilhabeziele  \newline erreicht werden sollen. &  & (4) \textcolor{diffgreen}{entfällt} \\
\hline
\rowcolor{gray!25}
\textbf{\S{} 136} &  & \textbf{\S{} 136} \\
\hline
Beitrag aus Einkommen zu den  \newline Aufwendungen &  & Beitrag aus Einkommen zu den  \newline Aufwendungen \\
\hline
(1) Bei den Leistungen nach diesem Teil ist  \newline ein Beitrag zu den Aufwendungen  \newline aufzubringen, wenn das Einkommen im  \newline Sinne des \S{} 135 der antragstellenden  \newline Person sowie bei minderjährigen Personen  \newline der im Haushalt lebenden Eltern oder des  \newline im Haushalt lebenden Elternteils die  \newline Beträge nach Absatz 2 übersteigt. &  & (1) Bei den Leistungen nach diesem Teil ist  \newline ein Beitrag zu den Aufwendungen  \newline aufzubringen, wenn das Einkommen im  \newline Sinne des \S{} 135 der antragstellenden  \newline Person sowie bei minderjährigen Personen  \newline der im Haushalt lebenden Eltern oder des  \newline im Haushalt lebenden Elternteils die  \newline Beträge nach Absatz 2 übersteigt. \\
\hline
Absätze 2 bis 5 &  & unverändert \\
\hline
\rowcolor{gray!25}
\textbf{\S{} 138} &  & \textbf{\S{} 138} \\
\hline
Besondere Höhe des Beitrages zu den  \newline Aufwendungen &  & Besondere Höhe des Beitrages zu den  \newline Aufwendungen \\
\hline
(1) Ein Beitrag ist nicht aufzubringen bei  \newline 1.  \newline heilpädagogischen Leistungen nach \S{} 113  \newline Absatz 2 Nummer 3,  \newline 2.  \newline Leistungen zur medizinischen  \newline Rehabilitation nach \S{} 109,  \newline 3.  \newline Leistungen zur Teilhabe am Arbeitsleben  \newline nach \S{} 111 Absatz 1,  \newline 4.  \newline Leistungen zur Teilhabe an Bildung nach \S{}  \newline 112 Absatz 1 Nummer 1,  \newline 5.  \newline Leistungen zur schulischen oder  \newline hochschulischen Ausbildung oder  \newline Weiterbildung für einen Beruf nach \S{} 112  \newline Absatz 1 Nummer 2, soweit diese  \newline Leistungen in besonderen  \newline Ausbildungsstätten über Tag und Nacht für  \newline Menschen mit Behinderungen erbracht  \newline werden,  \newline 6.  \newline Leistungen zum Erwerb und Erhalt  \newline praktischer Kenntnisse und Fähigkeiten  \newline nach \S{} 113 Absatz 2 Nummer 5, soweit  \newline diese der Vorbereitung auf die Teilhabe am  \newline Arbeitsleben nach \S{} 111 Absatz 1 dienen,  \newline 7.  \newline Leistungen nach \S{} 113 Absatz 1, die noch  \newline nicht eingeschulten leistungsberechtigten  \newline Personen die für sie erreichbare Teilnahme  \newline am Leben in der Gemeinschaft  \newline ermöglichen sollen,  \newline 8.  \newline gleichzeitiger Gewährung von Leistungen  \newline zum Lebensunterhalt nach dem Zweiten  \newline oder Zwölften Buch oder nach \S{} 27a des  \newline Bundesversorgungsgesetzes. &  & (1) Ein Beitrag ist nicht aufzubringen bei  \newline 1.  \newline heilpädagogischen Leistungen nach \S{} 113  \newline Absatz 2 Nummer 3,  \newline 2.  \newline Leistungen zur medizinischen  \newline Rehabilitation nach \S{} 109,  \newline 3.  \newline Leistungen zur Teilhabe am Arbeitsleben  \newline nach \S{} 111 Absatz 1,  \newline 4.  \newline Leistungen zur Teilhabe an Bildung nach \S{}  \newline 112 Absatz 1 Nummer 1,  \newline 5.  \newline Leistungen zur schulischen oder  \newline hochschulischen Ausbildung oder  \newline Weiterbildung für einen Beruf nach \S{} 112  \newline Absatz 1 Nummer 2, soweit diese  \newline Leistungen in besonderen  \newline Ausbildungsstätten über Tag und Nacht für  \newline Menschen mit Behinderungen erbracht  \newline werden,  \newline 6.  \newline Leistungen zum Erwerb und Erhalt  \newline praktischer Kenntnisse und Fähigkeiten  \newline nach \S{} 113 Absatz 2 Nummer 5, soweit  \newline diese der Vorbereitung auf die Teilhabe am  \newline Arbeitsleben nach \S{} 111 Absatz 1 dienen,  \newline 7.  \newline Leistungen nach \S{} 113 Absatz 1, die noch  \newline nicht eingeschulten leistungsberechtigten  \newline Personen die für sie erreichbare Teilnahme  \newline am Leben in der Gemeinschaft  \newline ermöglichen sollen,  \newline \textbf{\textcolor{diffgreen}{7}}.  \newline gleichzeitiger Gewährung von Leistungen  \newline zum Lebensunterhalt nach dem Zweiten  \newline oder Zwölften Buch oder nach \S{} 27a des  \newline Bundesversorgungsgesetzes. \\
\hline
(2) Wenn ein Beitrag nach \S{} 137  \newline aufzubringen ist, ist für weitere Leistungen  \newline im gleichen Zeitraum oder weitere  \newline Leistungen an minderjährige Kinder im  \newline gleichen Haushalt nach diesem Teil kein  \newline weiterer Beitrag aufzubringen. &  & (2) Wenn ein Beitrag nach \S{} 137  \newline aufzubringen ist, ist für weitere Leistungen  \newline im gleichen Zeitraum oder weitere  \newline Leistungen an minderjährige Kinder im  \newline gleichen Haushalt nach diesem Teil kein  \newline weiterer Beitrag aufzubringen. \\
\hline
\rowcolor{gray!25}
\textbf{\S{} 140} &  & \textbf{\S{} 140} \\
\hline
Einsatz des Vermögens &  & Einsatz des Vermögens \\
\hline
(1) Die antragstellende Person sowie bei  \newline minderjährigen Personen die im Haushalt  \newline lebenden Eltern oder ein Elternteil haben  \newline vor der Inanspruchnahme von Leistungen  \newline nach diesem Teil die erforderlichen Mittel  \newline aus ihrem Vermögen aufzubringen. &  & \textbf{(1) Die antrag}\textbf{\textcolor{diffgreen}{s}}\textbf{stellende Person }\textbf{\textcolor{diffgreen}{hat}}\textbf{ vor  \newline der Inanspruchnahme von Leistungen  \newline nach diesem Teil die erforderlichen  \newline Mittel aus ihrem Vermögen  \newline aufzubringen. } \\
\hline
Absätze 2 und 3 &  & unverändert \\
\hline
\rowcolor{gray!25}
\textbf{\S{} 142} &  & \textbf{\S{} 142} \\
\hline
Sonderregelungen für minderjährige  \newline Leistungsberechtigte und in Sonderfällen &  & Sonderregelungen für minderjährige  \newline Leistungsberechtigte und in Sonderfällen \\
\hline
(1) Minderjährigen Leistungsberechtigten  \newline und ihren Eltern oder einem Elternteil ist  \newline bei Leistungen im Sinne des \S{} 138 Absatz  \newline 1 Nummer 1, 2, 4, 5 und 7 die Aufbringung  \newline der Mittel für die Kosten des  \newline Lebensunterhalts nur in Höhe der für den  \newline häuslichen Lebensunterhalt ersparten  \newline Aufwendungen zuzumuten, soweit  \newline Leistungen über Tag und Nacht oder über  \newline Tag erbracht werden. &  & (1) \textcolor{diffgreen}{entfällt} \\
\hline
(2) Sind Leistungen von einem oder  \newline mehreren Anbietern über Tag und Nacht  \newline oder über Tag oder für ärztliche oder  \newline ärztlich verordnete Maßnahmen  \newline erforderlich, sind die Leistungen, die der  \newline Vereinbarung nach \S{} 134 Absatz 3  \newline zugrunde liegen, durch den Träger der  \newline Eingliederungshilfe auch dann in vollem  \newline Umfang zu erbringen, wenn den  \newline minderjährigen Leistungsberechtigten und  \newline ihren Eltern oder einem Elternteil die  \newline Aufbringung der Mittel nach Absatz 1 zu  \newline einem Teil zuzumuten ist. In Höhe dieses  \newline Teils haben sie zu den Kosten der  \newline erbrachten Leistungen beizutragen;  \newline mehrere Verpflichtete haften als  \newline Gesamtschuldner. &  & (2) \textcolor{diffgreen}{entfällt} \\
\hline
(3) Die Absätze 1 und 2 gelten  \newline entsprechend für volljährige  \newline Leistungsberechtigte, wenn diese  \newline Leistungen erhalten, denen  \newline Vereinbarungen nach \S{} 134 Absatz 4  \newline zugrunde liegen. In diesem Fall ist den  \newline volljährigen Leistungsberechtigten die  \newline Aufbringung der Mittel für die Kosten des  \newline Lebensunterhalts nur in Höhe der für ihren  \newline häuslichen Lebensunterhalt ersparten  \newline Aufwendungen zuzumuten.  \textbf{Erstes Buch Sozialgesetzbuch } &  & (3) \textcolor{diffgreen}{entfällt}  \textbf{Erstes Buch Sozialgesetzbuch } \\
\hline
- Aus Synopsis 2024 - \newline \textbf{( - SGB I)  \newline vom: 11.12.1975 - zuletzt geändert durch  \newline Art. 2a v. 22.12.2023 I Nr. 408 } \newline - Aus Synopsis 2026 - \newline \textbf{( - SGB I)  \newline vom: 11.12.1975 - zuletzt geändert durch  \newline Art. 62 Absatz 3 v. 4.2.2026 I Nr. 33 } & \textbf{( - SGB I)  \newline vom: 11.12.1975 - zuletzt geändert durch  \newline Art. }\textbf{\textcolor{diffred}{2a}}\textbf{ v. }\textbf{\textcolor{diffred}{22}}\textbf{.}\textbf{\textcolor{diffred}{1}}\textbf{2.202}\textbf{\textcolor{diffred}{3}}\textbf{ I Nr. }\textbf{\textcolor{diffred}{408}} & \textbf{( - SGB I)  \newline vom: 11.12.1975 - zuletzt geändert durch  \newline Art. }\textbf{\textcolor{diffgreen}{62 Absatz 3}}\textbf{ v. }\textbf{\textcolor{diffgreen}{4}}\textbf{.2.202}\textbf{\textcolor{diffgreen}{6}}\textbf{ I Nr. }\textbf{\textcolor{diffgreen}{33}} \\
\hline
... & ... & ... \\
\hline
4. Hilfe zur Erziehung, Eingliederungshilfe  \newline für seelisch behinderte Kinder und  \newline Jugendliche sowie Hilfe für junge  \newline Volljährige. & 4. \textcolor{diffred}{H}i\textcolor{diffred}{lf}e zur \textcolor{diffred}{Erziehung, Eingliederungshilfe  \newline für seelisch behinderte Kinder} und  \newline \textcolor{diffred}{Jugendliche }\textbf{\textcolor{diffred}{mit Behinderungen}}\textcolor{diffred}{ sowie  \newline Hilfe} für junge Volljährige. & 4. \textbf{\textcolor{diffgreen}{Le}}\textbf{i}\textbf{\textcolor{diffgreen}{stung}}\textbf{e}\textbf{\textcolor{diffgreen}{n}}\textbf{ zur }\textbf{\textcolor{diffgreen}{Entwicklung, zur  \newline Erziehung}}\textbf{ und }\textbf{\textcolor{diffgreen}{zur Teilhabe, Hilfen}}\textbf{ für  \newline junge Volljährige.} \\
\hline
(2) Zuständig sind die Kreise und die  \newline kreisfreien Städte, nach Maßgabe des  \newline Landesrechts auch kreisangehörige  \newline Gemeinden; sie arbeiten mit der freien  \newline Jugendhilfe zusammen.  \textbf{Zweites Buch Sozialgesetzbuch  } & (2) unverändert  \textbf{Zweites Buch Sozialgesetzbuch } & (2) unverändert  \textbf{Zweites Buch Sozialgesetzbuch } \\
\hline
- Aus Synopsis 2024 - \newline \textbf{( - SGB II)  \newline vom: 13.05.2011 - zuletzt geändert durch  \newline Art. 5 v. 27.03.2024 I Nr. 107} \newline - Aus Synopsis 2026 - \newline \textbf{( - SGB II)  \newline vom: 13.05.2011 - zuletzt geändert durch  \newline Art. 8 v. 22.12.2025 I Nr. 363} & \textbf{( - SGB II)  \newline vom: 13.05.2011 - zuletzt geändert durch  \newline Art. }\textbf{\textcolor{diffred}{5}}\textbf{ v. 2}\textbf{\textcolor{diffred}{7}}\textbf{.}\textbf{\textcolor{diffred}{03}}\textbf{.202}\textbf{\textcolor{diffred}{4}}\textbf{ I Nr. }\textbf{\textcolor{diffred}{107}} & \textbf{( - SGB II)  \newline vom: 13.05.2011 - zuletzt geändert durch  \newline Art. }\textbf{\textcolor{diffgreen}{8}}\textbf{ v. 2}\textbf{\textcolor{diffgreen}{2}}\textbf{.}\textbf{\textcolor{diffgreen}{12}}\textbf{.202}\textbf{\textcolor{diffgreen}{5}}\textbf{ I Nr. }\textbf{\textcolor{diffgreen}{363}} \\
\hline
... & ... & ... \\
\hline
(7) Nicht als Einkommen zu  \newline berücksichtigen sind Einnahmen von  \newline Schülerinnen und Schülern allgemein- oder  \newline berufsbildender Schulen, die das 25.  \newline Lebensjahr noch nicht vollendet haben, aus  \newline Erwerbstätigkeiten, die in den Schulferien  \newline ausgeübt werden. Satz 1 gilt nicht für eine  \newline Ausbildungsvergütung, auf die eine  \newline Schülerin oder ein Schüler einen Anspruch  \newline hat.  \textbf{Fünftes Buch Sozialgesetzbuch } & (7) unverändert  \textbf{Fünftes Buch Sozialgesetzbuch } & (7) unverändert  \textbf{Fünftes Buch Sozialgesetzbuch } \\
\hline
- Aus Synopsis 2024 - \newline \textbf{( - SGB V)  \newline vom: 20.12.1988 - zuletzt geändert durch  \newline Art. 33 u. 35 Absatz 10 v. 27.03.2024 I Nr.  \newline 108 } \newline - Aus Synopsis 2026 - \newline \textbf{( - SGB V)  \newline vom: 20.12.1988 - zuletzt geändert durch  \newline Art. 8 v. 3.2.2026 I Nr. 28 } & \textbf{( - SGB V)  \newline vom: 20.12.1988 - zuletzt geändert durch  \newline Art. }\textbf{\textcolor{diffred}{33 u. 35 Absatz 10}}\textbf{ v. 2}\textbf{\textcolor{diffred}{7.03}}\textbf{.202}\textbf{\textcolor{diffred}{4}}\textbf{ I Nr.} \newline \textbf{\textcolor{diffred}{10}}\textbf{8 } & \textbf{( - SGB V)  \newline vom: 20.12.1988 - zuletzt geändert durch  \newline Art. }\textbf{\textcolor{diffgreen}{8}}\textbf{ v. }\textbf{\textcolor{diffgreen}{3.}}\textbf{2.202}\textbf{\textcolor{diffgreen}{6}}\textbf{ I Nr. }\textbf{\textcolor{diffgreen}{2}}\textbf{8 } \\
\hline
... & ... & ... \\
\hline
(4) \S{} 45 Absatz 3 gilt entsprechend. Den  \newline Anspruch auf unbezahlte Freistellung von  \newline der Arbeitsleistung haben auch  \newline Arbeitnehmer, die nicht Versicherte mit  \newline Anspruch auf Krankengeld nach Absatz 1  \newline sind.  \textbf{Vierzehntes Buch Sozialgesetzbuch  } & (4) unverändert  \textbf{Vierzehntes Buch Sozialgesetzbuch } & (4) unverändert  \textbf{Vierzehntes Buch Sozialgesetzbuch } \\
\hline
- Aus Synopsis 2024 - \newline \textbf{( - SGB XIV)  \newline vom: 12.12.2019 - zuletzt geändert durch  \newline Art. 11 v. 22.12.2023 I Nr. 408 } \newline - Aus Synopsis 2026 - \newline \textbf{( - SGB XIV)  \newline vom: 12.12.2019 - zuletzt geändert durch  \newline Art. 1 der Verordnung v. 13.6.2025 I Nr.  \newline 144 } & \textbf{( - SGB XIV)  \newline vom: 12.12.2019 - zuletzt geändert durch  \newline Art. }\textbf{\textcolor{diffred}{11}}\textbf{ v. }\textbf{\textcolor{diffred}{22}}\textbf{.}\textbf{\textcolor{diffred}{12}}\textbf{.202}\textbf{\textcolor{diffred}{3}}\textbf{ I Nr. 4}\textbf{\textcolor{diffred}{08}} & \textbf{( - SGB XIV)  \newline vom: 12.12.2019 - zuletzt geändert durch  \newline Art. }\textbf{\textcolor{diffgreen}{1 der Verordnung}}\textbf{ v. }\textbf{\textcolor{diffgreen}{13}}\textbf{.}\textbf{\textcolor{diffgreen}{6}}\textbf{.202}\textbf{\textcolor{diffgreen}{5}}\textbf{ I Nr.} \newline \textbf{\textcolor{diffgreen}{1}}\textbf{4}\textbf{\textcolor{diffgreen}{4}} \\
\hline
\rowcolor{gray!25}
\textbf{\S{} 65} &  & \textbf{\S{} 65} \\
\hline
Leistungen zur Teilhabe an Bildung &  & Leistungen zur Teilhabe an Bildung \\
\hline
Geschädigte, die auf Grund der  \newline Schädigungsfolgen zum  \newline leistungsberechtigten Personenkreis im  \newline Sinne von \S{} 99 des Neunten Buches  \newline gehören, erhalten Leistungen zur Teilhabe  \newline an Bildung entsprechend Teil 2 Kapitel 5  \newline des Neunten Buches. &  & Geschädigte, die auf Grund der  \newline Schädigungsfolgen zum  \newline leistungsberechtigten Personenkreis im  \newline Sinne von \S{} 99 des Neunten Buches  \newline gehören, erhalten Leistungen zur Teilhabe  \newline an Bildung entsprechend Teil 2 Kapitel 5  \newline des Neunten Buches. \textbf{\textcolor{diffgreen}{Geschädigte Kinder  \newline oder Jugendliche, die aufgrund der  \newline Schädigungsfolgen zum  \newline leistungsberechtigten Personenkreis im  \newline Sinne von \S{} 27 Absatz 3 des Achten  \newline Buches gehören, erhalten Leistungen  \newline zur Teilhabe an Bildung entsprechend \S{}  \newline 35d des Achten Buches.}} \\
\hline
\rowcolor{gray!25}
\textbf{\S{} 66} &  & \textbf{\S{} 66} \\
\hline
Leistungen zur Sozialen Teilhabe &  & Leistungen zur Sozialen Teilhabe \\
\hline
(1) Geschädigte, die auf Grund der  \newline Schädigungsfolgen zum  \newline leistungsberechtigten Personenkreis im  \newline Sinne von \S{} 99 des Neunten Buches  \newline gehören, erhalten Leistungen zur Sozialen  \newline Teilhabe entsprechend Teil 2 Kapitel 6 des  \newline Neunten Buches. &  & (1) Geschädigte, die auf Grund der  \newline Schädigungsfolgen zum  \newline leistungsberechtigten Personenkreis im  \newline Sinne von \S{} 99 des Neunten Buches  \newline gehören, erhalten Leistungen zur Sozialen  \newline Teilhabe entsprechend Teil 2 Kapitel 6 des  \newline Neunten Buches. \textbf{\textcolor{diffgreen}{Geschädigte Kinder  \newline oder Jugendliche, die aufgrund der  \newline Schädigungsfolgen zum  \newline leistungsberechtigten Personenkreis im  \newline Sinne von \S{} 27 Absatz 3 des Achten  \newline Buches gehören, erhalten Leistungen  \newline zur Sozialen Teilhabe entsprechend \S{}\S{}  \newline 35f, 35h und 35i des Achten Buches.}} \\
\hline
(2) Abweichend von Absatz 1 werden  \newline Leistungen zur Mobilität nach \S{} 83 des  \newline Neunten Buches erbracht. Sie umfassen  \newline zudem Leistungen zum Betrieb, Unterhalt,  \newline Unterstellen und Abstellen eines  \newline Kraftfahrzeuges. &  & (2) unverändert \\
\hline
... & ... & ... \\
\hline
\textbf{Jugendschutzgesetz } &  & \textbf{Jugendschutzgesetz } \\
\hline
\textbf{Vom 23. Juli 202 (BGBL. I S. 2730),  \newline zuletzt geändert durch Art. 12 v. 6. Mai  \newline 2024 I Nr. 149 } &  & \textbf{Vom 23. Juli 202 (BGBL. I S. 2730),  \newline zuletzt geändert durch Art. 12 v. 6. Mai  \newline 2024 I Nr. 149 } \\
\hline
\rowcolor{gray!25}
\textbf{\S{} 9} &  & \textbf{\S{} 9} \\
\hline
Absatz 1 &  & Unverändert \\
\hline
(2) Absatz 1 Nummer 1 gilt nicht, wenn  \newline Jugendliche von einer  \newline personensorgeberechtigten Person  \newline begleitet werden. &  & (2) \textcolor{diffgreen}{entfällt} \\
\hline
Absätze 3 bis 4 &  & Absätze \textcolor{diffgreen}{2 und} 3 \\
\hline
\end{longtable}

\end{document}